
\Data\ID{1003171101}
\Updated{2022-04-23 05:04:11}
\Level{A}
\SubArea{Linear functions}
\LANGen{
\Question{
Given the linear function \( f(x)=kx-4 \), find the value of \( k \) so that \( f(2)= 2 \).}
{\( k=3 \)}{\( k=4 \)}{\( k=2 \)}{\( k=-2 \)}{}{}}
\LANGpl{
\Question{
Dana jest funkcja liniowa \( f(x)=kx-4 \). Wyznacz wartość  \( k \) tak, aby  \( f(2)= 2 \).}
{\( k=3 \)}{\( k=4 \)}{\( k=2 \)}{\( k=-2 \)}{}{}}
\LANGcs{
\Question{
Je dána lineární funkce \( f(x)=kx-4 \). Pro které \( k \) platí, že \( f(2)= 2 \)?}
{\( k=3 \)}{\( k=4 \)}{\( k=2 \)}{\( k=-2 \)}{}{}}
\LANGsk{
\Question{
Daná je lineárna funkcia \( f(x)=kx-4 \). Určte hodnotu \( k \), ak platí \( f(2)= 2 \).}
{\( k=3 \)}{\( k=4 \)}{\( k=2 \)}{\( k=-2 \)}{}{}}
\LANGes{
\Question{
Sea una función lineal \( f(x)=kx-4 \). Halla el valor de \( k \) para que \( f(2)= 2 \).}
{\( k=3 \)}{\( k=4 \)}{\( k=2 \)}{\( k=-2 \)}{}{}}
\End

\Data\ID{1003171102}
\Updated{2022-04-23 05:04:11}
\Level{A}
\SubArea{Linear functions}
\LANGen{
\Question{
Consider the linear function \( f(x) = -4x +3 \). Give the input value for \( f \) such that the output value of \( f \) is \( -9 \).}
{\( 3 \)}{\( -3 \)}{\( 39 \)}{\( -33 \)}{}{}}
\LANGpl{
\Question{
Rozważ funkcję liniową \( f(x) = -4x +3 \). Podaj wartość dodaną  \( f \) tak, by otrzymana wartość \( f \) wyniosła \( -9 \).}
{\( 3 \)}{\( -3 \)}{\( 39 \)}{\( -33 \)}{}{}}
\LANGcs{
\Question{
Je dána lineární funkce \( f(x) = -4x +3 \). Ve kterém bodě bude hodnota funkce \( f \) rovna \( -9 \)?}
{\( 3 \)}{\( -3 \)}{\( 39 \)}{\( -33 \)}{}{}}
\LANGsk{
\Question{
Daná je lineárna funkcia \( f(x) = -4x +3 \). V ktorom bode bude hodnota funkcie \( f \) rovná \( -9 \)?}
{\( 3 \)}{\( -3 \)}{\( 39 \)}{\( -33 \)}{}{}}
\LANGes{
\Question{
Sea una función lineal \( f(x) = -4x +3 \). Halla el valor de entrada para \( f \) si el valor de salida de \( f \) es \( -9 \).}
{\( 3 \)}{\( -3 \)}{\( 39 \)}{\( -33 \)}{}{}}
\End

\Data\ID{1103171103}
\Updated{2022-07-06 12:07:05}
\Level{A}
\SubArea{Linear functions}
\LANGen{
\Question{
Consider the linear function \( f(x)=-2x+4 \). Which of the given graphs is the graph of \( f \)?}
{}{}{}{}{}{}}
\LANGpl{
\Question{
Rozważ funkcję  \( f(x)=-2x+4 \). Który z podanych wykresów jest wykresem funkcji \( f \)?}
{}{}{}{}{}{}}
\LANGcs{
\Question{
Je dána lineární funkce \( f(x)=-2x+4 \). Vyberte obrázek, na kterém je graf této funkce.}
{}{}{}{}{}{}}
\LANGsk{
\Question{
Daná je lineárna funkcia \( f(x)=-2x+4 \). Ktorý z daných grafov je graf funkcie \( f \)?}
{}{}{}{}{}{}}
\LANGes{
\Question{
Sea la función lineal \( f(x)=-2x+4 \). ¿Cuál de las siguientes gráficas corresponde a la función \( f \)?}
{}{}{}{}{}{}}

\IMGanswer{1103171103a.pdf}
\IMGanswer{1103171103b.pdf}
\IMGanswer{1103171103c.pdf}
\IMGanswer{1103171103d.pdf}\End

\Data\ID{1103171403}
\Updated{2022-04-23 05:04:11}
\Level{A}
\SubArea{Linear functions}
\IMG{1103171403_0.pdf}
\LANGen{
\Question{
Determine whether the line drawn in the picture is the graph of a linear function of the variable \(x\). If so, find the formula for the function.}
{It is not the graph of a linear function in the picture.}{\( y=5 \)}{\( x=5 \)}{\( y=5x \)}{}{}}
\LANGpl{
\Question{
Określ czy na poniższych rysunkach przedstawiono wykres funkcji liniowej zmiennej \(x\). Jeśli tak, wyznacz wzór tej funkcji.}
{Rysunek nie przedstawia wykresu funkcji liniowej.}{\( y=5 \)}{\( x=5 \)}{\( y=5x \)}{}{}}
\LANGcs{
\Question{
Rozhodněte, zda přímka nakreslená na obrázku představuje graf lineární funkce proměnné \(x\). Pokud ano, určete její předpis.}
{Není to graf lineární funkce.}{\( y=5 \)}{\( x=5 \)}{\( y=5x \)}{}{}}
\LANGsk{
\Question{
Rozhodnite, či na obrázku sa nachádza graf lineárnej funkcie premennej \(x\). Ak áno, tak určte jej predpis.}
{Na obrázku sa nenachádza graf lineárnej funkcie.}{\( y=5 \)}{\( x=5 \)}{\( y=5x \)}{}{}}
\LANGes{
\Question{
Determina si la recta en el dibujo es la gráfica de una función lineal de la variable \(x\). Si lo es, halla la fórmula de la función.}
{No es gráfica de una función lineal.}{\( y=5 \)}{\( x=5 \)}{\( y=5x \)}{}{}}
\End

\Data\ID{1103171405}
\Updated{2020-12-30 11:12:29}
\Level{A}
\SubArea{Linear functions}
\IMG{1103171405_0.pdf}
\LANGen{
\Question{
Choose the true statement.}
{The given points do not lay on the graph of a linear function.}{The given points lay on the graph of the linear function \( f(x)=-x+3 \).}{The given points lay on the graph of the linear function \( f(x)=-x+3.5 \).}{The given points lay on the graph of the linear function \( f(x)=x+3 \).}{}{}}
\LANGpl{
\Question{
Wybierz prawdziwe stwierdzenie.}
{Podane punkty nie leżą na wykresie funkcji liniowej.}{Podane punkty leżą na wykresie funkcji liniowej \( f(x)=-x+3 \).}{Podane punkty leżą na wykresie funkcji liniowej \( f(x)=-x+3{,}5 \).}{Podane punkty leżą na wykresie funkcji liniowej \( f(x)=x+3 \).}{}{}}
\LANGcs{
\Question{
Na obrázku jsou dány čtyři body. Vyberte pravdivé tvrzení.}
{Body neleží na grafu žádné lineární funkce.}{Body leží na grafu lineární funkce \( f(x)=-x+3 \).}{Body leží na grafu lineární funkce \( f(x)=-x+3{,}5 \).}{Body leží na grafu lineární funkce \( f(x)=x+3 \).}{}{}}
\LANGsk{
\Question{
Vyberte pravdivé tvrdenie.}
{Dané body neležia na grafe lineárnej funkcie.}{Dané body ležia na grafe lineárnej funkcie \( f(x)=-x+3 \).}{Dané body ležia na grafe lineárnej funkcie \( f(x)=-x+3{,}5 \).}{Dané body ležia na grafe lineárnej funkcie \( f(x)=x+3 \).}{}{}}
\LANGes{
\Question{
Elige la declaración correcta.}
{Los puntos dados no pertenecen a la gráfica de ninguna función lineal.}{Los puntos dados pertenecen a la gráfica de la función lineal \( f(x)=-x+3 \).}{Los puntos dados pertenecen a la gráfica de la función lineal \( f(x)=-x+3.5 \).}{Los puntos dados pertenecen a la gráfica de la función lineal \( f(x)=x+3 \).}{}{}}
\End

\Data\ID{1103171406}
\Updated{2022-04-23 05:04:11}
\Level{A}
\SubArea{Linear functions}
\IMG{1103171406.pdf}
\LANGen{
\Question{
Choose the formula of the function whose graph is shown in the picture.}
{\( f(x)=-\frac54x-\frac74;\ x\in(-3;1] \)}{\( f(x)=-\frac54x-\frac74;\ x\in[-3;2) \)}{\( f(x)=\frac54x-\frac{17}4;\ x\in(-3;1] \)}{\( f(x)=-\frac54x+\frac74;\ x\in(-3;1] \)}{}{}}
\LANGpl{
\Question{
Wybierz wzór dla funkcji przedstawionego na rysunku poniżej.}
{\( f(x)=-\frac54x-\frac74;\ x\in(-3;1\rangle \)}{\( f(x)=-\frac54x-\frac74;\ x\in\langle-3;2) \)}{\( f(x)=\frac54x-\frac{17}4;\ x\in(-3;1\rangle \)}{\( f(x)=-\frac54x+\frac74;\ x\in(-3;1\rangle \)}{}{}}
\LANGcs{
\Question{
Vyberte předpis funkce, jejíž graf je na obrázku.}
{\( f(x)=-\frac54x-\frac74;\ x\in(-3;1\rangle \)}{\( f(x)=-\frac54x-\frac74;\ x\in\langle-3;2) \)}{\( f(x)=\frac54x-\frac{17}4;\ x\in(-3;1\rangle \)}{\( f(x)=-\frac54x+\frac74;\ x\in(-3;1\rangle \)}{}{}}
\LANGsk{
\Question{
Vyberte predpis funkcie, ktorej graf je na obrázku.}
{\( f(x)=-\frac54x-\frac74;\ x\in(-3;1\rangle \)}{\( f(x)=-\frac54x-\frac74;\ x\in\langle-3;2) \)}{\( f(x)=\frac54x-\frac{17}4;\ x\in(-3;1\rangle \)}{\( f(x)=-\frac54x+\frac74;\ x\in(-3;1\rangle \)}{}{}}
\LANGes{
\Question{
Elige la fórmula de la función lineal representada en el dibujo.}
{\( f(x)=-\frac54x-\frac74;\ x\in(-3;1] \)}{\( f(x)=-\frac54x-\frac74;\ x\in[-3;2) \)}{\( f(x)=\frac54x-\frac{17}4;\ x\in(-3;1] \)}{\( f(x)=-\frac54x+\frac74;\ x\in(-3;1] \)}{}{}}
\End

\Data\ID{2000000802}
\Updated{2022-01-17 06:01:45}
\Level{A}
\SubArea{Linear functions}
\LANGen{
\Question{
For a linear function \(f\) is given \(f(0)=4\) and \(f(2)=0\). Find the corresponding function \(f\).}
{\(f(x)=-2x+4\)}{\(f(x)=2x+4\)}{\(f(x)=4x+2\)}{\(f(x)=2x-4\)}{}{}}
\LANGpl{
\Question{
Dla funkcji \(f\) dane jest \(f(0)=4\) i \(f(2)=0\). Wskaż odpowiedni wzór funkcji \(f\).}
{\(f(x)=-2x+4\)}{\(f(x)=2x+4\)}{\(f(x)=4x+2\)}{\(f(x)=2x-4\)}{}{}}
\LANGcs{
\Question{
Pro kterou z níže uvedených funkcí \(f\) platí, že \(f(0)=4\) a \(f(2)=0\).}
{\(f(x)=-2x+4\)}{\(f(x)=2x+4\)}{\(f(x)=4x+2\)}{\(f(x)=2x-4\)}{}{}}
\LANGsk{
\Question{
Pre lineárnu funkciu \(f\) platí \(f(0)=4\) a \(f(2)=0\). Nájdite predpis funkcie \(f\).}
{\(f(x)=-2x+4\)}{\(f(x)=2x+4\)}{\(f(x)=4x+2\)}{\(f(x)=2x-4\)}{}{}}
\LANGes{
\Question{
Para una función lineal \(f\) tenemos que \(f(0)=4\) y \(f(2)=0\). Halla la función correspondiente \(f\).}
{\(f(x)=-2x+4\)}{\(f(x)=2x+4\)}{\(f(x)=4x+2\)}{\(f(x)=2x-4\)}{}{}}
\End

\Data\ID{2000000803}
\Updated{2022-01-17 06:01:12}
\Level{A}
\SubArea{Linear functions}
\LANGen{
\Question{
For a linear function \(f\) is given \(f(2)=4\) and \(f(1)=0\). Find the intersection point \(P\) of the graph of \(f\) with the \(y\)-axis.}
{\(P=[0;-4]\)}{\(P=[-4;0]\)}{\(P=[0;1]\)}{\(P=[1;0]\)}{}{}}
\LANGpl{
\Question{
Dla funkcji liniowej \(f\) dane jest \(f(2)=4\) i \(f(1)=0\). Znajdź współrzędne punktu \(P\), gdzie wykres funkcji \(f\) przecina się z osią \(y\).}
{\(P=[0;-4]\)}{\(P=[-4;0]\)}{\(P=[0;1]\)}{\(P=[1;0]\)}{}{}}
\LANGcs{
\Question{
Je dána lineární funkce \(f\), pro kterou platí, že \(f(2)=4\) a \(f(1)=0\). Najděte průsečík \(P\) grafu funkce \(f\) se souřadnou osou \(y\).}
{\(P=[0;-4]\)}{\(P=[-4;0]\)}{\(P=[0;1]\)}{\(P=[1;0]\)}{}{}}
\LANGsk{
\Question{
Pre lineárnu funkciu \(f\) platí \(f(2)=4\) a \(f(1)=0\). Nájdite súradnice priesečníka  \(P\), v ktorom sa graf funkcie \(f\) pretína s  \(y\)-osou.}
{\(P=[0;-4]\)}{\(P=[-4;0]\)}{\(P=[0;1]\)}{\(P=[1;0]\)}{}{}}
\LANGes{
\Question{
Para una función lineal \(f\) tenemos que \(f(2)=4\) y \(f(1)=0\). Halla el punto de intersección \(P\) de la gráfica de \(f\) con el eje \(y\).}
{\(P=[0;-4]\)}{\(P=[-4;0]\)}{\(P=[0;1]\)}{\(P=[1;0]\)}{}{}}
\End

\Data\ID{2000000804}
\Updated{2022-01-17 06:01:09}
\Level{A}
\SubArea{Linear functions}
\LANGen{
\Question{
Given \(f(x)=2x+1\) and \(g(x)=x+5\), find the point \(P\) at which graphs \(f\) and \(g\) intersect.}
{\(P=[4;9]\)}{\(P=[2;5]\)}{\(P=[9;-4]\)}{\(P=[9;4]\)}{}{}}
\LANGpl{
\Question{
Dane są funkcje  \(f(x)=2x+1\) i  \(g(x)=x+5\). Znajdź współrzędne punktu \(P\), w którym przecinają się funkcje \(f\) i  \(g\).}
{\(P=[4;9]\)}{\(P=[2;5]\)}{\(P=[9;-4]\)}{\(P=[9;4]\)}{}{}}
\LANGcs{
\Question{
Jsou dány funkce \(f(x)=2x+1\) a \(g(x)=x+5\). Najděte průsečík \(P\) grafů obou funkcí.}
{\(P=[4;9]\)}{\(P=[2;5]\)}{\(P=[9;-4]\)}{\(P=[9;4]\)}{}{}}
\LANGsk{
\Question{
Dané sú funkcie  \(f(x)=2x+1\) a \(g(x)=x+5\). Nájdite súradnice bodu  \(P\), v ktorom sa funkcie \(f\) a \(g\) pretínajú.}
{\(P=[4;9]\)}{\(P=[2;5]\)}{\(P=[9;-4]\)}{\(P=[9;4]\)}{}{}}
\LANGes{
\Question{
Sea \(f(x)=2x+1\) y \(g(x)=x+5\), halla el punto \(P\) dónde intersecan las gráficas de \(f\) y \(g\).}
{\(P=[4;9]\)}{\(P=[2;5]\)}{\(P=[9;-4]\)}{\(P=[9;4]\)}{}{}}
\End

\Data\ID{2000000805}
\Updated{2021-12-01 03:12:17}
\Level{A}
\SubArea{Linear functions}
\LANGen{
\Question{
Determine the coefficient \(k\) so that the graph of the function \(f(x)=kx+2\) passes through the point \(A=[-2;6]\).}
{\(k=-2\)}{\(k=\frac{2}{3}\)}{\(k=2\)}{\(k=-4\)}{}{}}
\LANGpl{
\Question{
Wyznacz współczynnik \(k\), tak aby wykres funkcji  \(f(x)=kx+2\) przechodził przez punkt \(A=[-2;6]\).}
{\(k=-2\)}{\(k=\frac{2}{3}\)}{\(k=2\)}{\(k=-4\)}{}{}}
\LANGcs{
\Question{
Urete koeficient \(k\) tak, aby graf funkce \(f(x)=kx+2\) procházel  bodem \(A=[-2;6]\).}
{\(k=-2\)}{\(k=\frac{2}{3}\)}{\(k=2\)}{\(k=-4\)}{}{}}
\LANGsk{
\Question{
Určte koeficient \(k\) tak, aby graf funkcie  \(f(x)=kx+2\) prechádzal bodom  \(A=[-2;6]\).}
{\(k=-2\)}{\(k=\frac{2}{3}\)}{\(k=2\)}{\(k=-4\)}{}{}}
\LANGes{
\Question{
Determina el coeficiente \(k\) para que la gráfica de la función \(f(x)=kx+2\) pase por el punto \(A=[-2;6]\).}
{\(k=-2\)}{\(k=\frac{2}{3}\)}{\(k=2\)}{\(k=-4\)}{}{}}
\End

\Data\ID{2000000806}
\Updated{2022-05-05 07:05:42}
\Level{A}
\SubArea{Linear functions}
\LANGen{
\Question{
Consider the linear function \(f(x)=-x+3\). Choose the incorrect statement.}
{The function \(f\) is increasing.}{The intercept of \(f\) with \(x\)-axis has coordinates \([3;0]\).}{The function \(f\) is one-to-one function.}{The intercept of \(f\) with \(y\)-axis has coordinates \([0;3]\).}{}{}}
\LANGpl{
\Question{
Dana jest funkcja liniowa \(f(x)=-x+3\). Wybierz niepoprawne stwierdzenie.}
{Funkcja \(f\) jest rosnąca.}{Punkt przecięcia funkcji \(f\) z osią \(x\) ma współrzędne \([3;0]\).}{Funkcja \(f\) jest różnowartościowa.}{Punkt przecięcia funkcji \(f\) z osią \(y\) ma współrzędne \([0;3]\).}{}{}}
\LANGcs{
\Question{
Je dána funkce \(f(x)=-x+3\). Které tvrzení je nepravdivé?}
{Funkce \(f\) je rostoucí.}{Průsečík grafu funkce \(f\) s osou \(x\) má souřadnice \([3;0]\).}{Funkce \(f\) je prostá.}{Průsečík grafu funkce \(f\) s osou \(y\) má souřadnice \([0;3]\).}{}{}}
\LANGsk{
\Question{
Daná je lineárna funkcia \(f(x)=-x+3\). Vyberte nesprávne tvrdenie.}
{Funkcia \(f\) je rastúca.}{Priesečník funkcie \(f\)  s  \(x\)-osou má súradnice \([3;0]\).}{Funkcia  \(f\) je prostá.}{Priesečník funkcie  \(f\)  s  \(y\)-osou  má súradnice  \([0;3]\).}{}{}}
\LANGes{
\Question{
Considera la función lineal \(f(x)=-x+3\). Elije la declaración falsa.}
{La función \(f\) es creciente.}{La intersección de \(f\) con el eje \(x\) tiene coordinadas \([3;0]\).}{La función \(f\) es inyectiva.}{La intersección de \(f\) con el eje \(y\) tiene coordenadas \([0;3]\).}{}{}}
\End

\Data\ID{2000000807}
\Updated{2021-12-01 03:12:05}
\Level{A}
\SubArea{Linear functions}
\LANGen{
\Question{
Consider the linear function \(f(x)=-2x+6\). Choose the correct statement.}
{The function \(f\) is one-to-one.}{The function \(f\) is even.}{The function \(f\) is odd.}{The function \(f\) is increasing.}{}{}}
\LANGpl{
\Question{
Dana jest funkcja liniowa \(f(x)=-2x+6\). Wybierz poprawne stwierdzenie.}
{Funkcja \(f\) jest różnowartościowa.}{Funkcja \(f\) jest parzysta.}{Funkcja \(f\) jest nieparzysta.}{Funkcja \(f\) jest rosnąca.}{}{}}
\LANGcs{
\Question{
Je dána lineární funkce \(f(x)=-2x+6\). Které tvrzení je pravdivé?}
{Funkce \(f\) je prostá.}{Funkce \(f\) je sudá.}{Funkce \(f\) je lichá.}{Funkce \(f\) je rostoucí.}{}{}}
\LANGsk{
\Question{
Daná je lineárna funkcia \(f(x)=-2x+6\). Vyberte správne tvrdenie.}
{Funkcia  \(f\) je prostá.}{Funkcia  \(f\) je párna.}{Funkcia   \(f\) je nepárna.}{Funkcia  \(f\) je rastúca.}{}{}}
\LANGes{
\Question{
Considera la función lineal \(f(x)=-2x+6\). Elije la declaración verdadera.}
{La función \(f\) es inyectiva.}{La función \(f\) es par.}{La función \(f\) es impar.}{La función \(f\) es creciente.}{}{}}
\End

\Data\ID{2000000808}
\Updated{2021-12-01 03:12:59}
\Level{A}
\SubArea{Linear functions}
\LANGen{
\Question{
Consider the linear function \(f(x)=-6x-2\). Choose the correct statement.}
{\(f(3)=-20\)}{\(f(-2)=-14\)}{\(f(10)=-2\)}{\(f(-2)=0\)}{}{}}
\LANGpl{
\Question{
Dana jest funkcja liniowa \(f(x)=-6x-2\). Wybierz poprawne stwierdzenie.}
{\(f(3)=-20\)}{\(f(-2)=-14\)}{\(f(10)=-2\)}{\(f(-2)=0\)}{}{}}
\LANGcs{
\Question{
Je dána lineární funkce \(f(x)=-6x-2\). Které tvrzení je pravdivé?}
{\(f(3)=-20\)}{\(f(-2)=-14\)}{\(f(10)=-2\)}{\(f(-2)=0\)}{}{}}
\LANGsk{
\Question{
Daná je lineárna funkcia \(f(x)=-6x-2\). Vyberte správne tvrdenie.}
{\(f(3)=-20\)}{\(f(-2)=-14\)}{\(f(10)=-2\)}{\(f(-2)=0\)}{}{}}
\LANGes{
\Question{
Considera la función lineal \(f(x)=-6x-2\). Elije la declaración correcta.}
{\(f(3)=-20\)}{\(f(-2)=-14\)}{\(f(10)=-2\)}{\(f(-2)=0\)}{}{}}
\End

\Data\ID{2000000809}
\Updated{2021-12-01 03:12:51}
\Level{A}
\SubArea{Linear functions}
\LANGen{
\Question{
Consider the linear function \(f(x)=3x-2\). Which of the given points belongs to the graph of \(f\)?}
{\([-2;-8]\)}{\([-2;0]\)}{\(\left[0;-\frac{2}{3}\right]\)}{\(\left[0;-\frac{4}{3}\right]\)}{}{}}
\LANGpl{
\Question{
Dana jest funkcja liniowa \(f(x)=3x-2\). Który z podanych punktów należy do wykresu funkcji \(f\)?}
{\([-2;-8]\)}{\([-2;0]\)}{\(\left[0;-\frac{2}{3}\right]\)}{\(\left[0;-\frac{4}{3}\right]\)}{}{}}
\LANGcs{
\Question{
Je dána lineární funkce \(f(x)=3x-2\). Který z níže uvedených bodů náleží grafu této funkce \(f\)?}
{\([-2;-8]\)}{\([-2;0]\)}{\(\left[0;-\frac{2}{3}\right]\)}{\(\left[0;-\frac{4}{3}\right]\)}{}{}}
\LANGsk{
\Question{
Daná je lineárna funkcia \(f(x)=3x-2\). Ktorý z uvedených bodov leží na grafe funkcie \(f\)?}
{\([-2;-8]\)}{\([-2;0]\)}{\(\left[0;-\frac{2}{3}\right]\)}{\(\left[0;-\frac{4}{3}\right]\)}{}{}}
\LANGes{
\Question{
Sea una función lineal \(f(x)=3x-2\). ¿Cuál de los puntos pertenece a la gráfica de \(f\)?}
{\([-2;-8]\)}{\([-2;0]\)}{\(\left[0;-\frac{2}{3}\right]\)}{\(\left[0;-\frac{4}{3}\right]\)}{}{}}
\End

\Data\ID{2000000810}
\Updated{2021-12-01 03:12:45}
\Level{A}
\SubArea{Linear functions}
\LANGen{
\Question{
Consider a function \(f(x)=ax+b\), where \(x \in (-1;4)\) and the range of \(f\) is \((-3;2)\). Find the coefficients \(a\) and \(b\).}
{\(a=-1,~b=1\)}{\(a=-1,~b=4\)}{\(a=2,~b=-3\)}{\(a=1,~b=2\)}{}{}}
\LANGpl{
\Question{
Dana jest funkcja \(f(x)=ax+b\), gdzie \(x \in (-1;4)\) a zbiór wartości funkcji \(f\) to \((-3;2)\). Wyznacz współczynniki \(a\) i \(b\).}
{\(a=-1,~b=1\)}{\(a=-1,~b=4\)}{\(a=2,~b=-3\)}{\(a=1,~b=2\)}{}{}}
\LANGcs{
\Question{
Je dána funkce \(f(x)=ax+b\), pro kterou platí, že \(x \in (-1;4)\) a obor hodnot funkce \(f\) je \((-3;2)\). Určete koeficienty \(a\) a \(b\).}
{\(a=-1,~b=1\)}{\(a=-1,~b=4\)}{\(a=2,~b=-3\)}{\(a=1,~b=2\)}{}{}}
\LANGsk{
\Question{
Daná je funkcia  \(f(x)=ax+b\), kde \(x \in (-1;4)\) a obor hodnôt funkcie \(f\) je \((-3;2)\). Nájdite koeficienty \(a\) a \(b\).}
{\(a=-1,~b=1\)}{\(a=-1,~b=4\)}{\(a=2,~b=-3\)}{\(a=1,~b=2\)}{}{}}
\LANGes{
\Question{
Considera la función \(f(x)=ax+b\),  dónde \(x \in (-1;4)\) y el Rango de \(f\) es \((-3;2)\). Halla los coficientes \(a\) y \(b\).}
{\(a=-1,~b=1\)}{\(a=-1,~b=4\)}{\(a=2,~b=-3\)}{\(a=1,~b=2\)}{}{}}
\End

\Data\ID{2000003101}
\Updated{2021-12-01 04:12:50}
\Level{A}
\SubArea{Linear functions}
\LANGen{
\Question{
Which of the functions 
\[f (x)=5,\] 
\[g(x)= 0.3x-3,\] 
\[h (x)=-0.4x+5,\]
\[k (x)=-3x+2,\] 
\[m (x)=12x+4\] 
are increasing?}
{\(g\), \(m\)}{\(h\), \(k\)}{\(f\), \(g\), \(m\)}{\(f\), \(k\), \(m\)}{}{}}
\LANGpl{
\Question{
Która z podanych funkcji jest funkcją rosnącą? 
\[f (x)=5\] 
\[g(x)= 0{,}3x-3\] 
\[h (x)=-0{,}4x+5\]
\[k (x)=-3x+2\] 
\[m (x)=12x+4\]}
{\(g\), \(m\)}{\(h\), \(k\)}{\(f\), \(g\), \(m\)}{\(f\), \(k\), \(m\)}{}{}}
\LANGcs{
\Question{
Které z funkcí 
\[f (x)=5,\] 
\[g(x)= 0{,}3x-3,\] 
\[h (x)=-0{,}4x+5,\]
\[k (x)=-3x+2,\] 
\[m (x)=12x+4\] 
jsou rostoucí?}
{\(g\), \(m\)}{\(h\), \(k\)}{\(f\), \(g\), \(m\)}{\(f\), \(k\), \(m\)}{}{}}
\LANGsk{
\Question{
Ktoré z funkcií 
\[f (x)=5,\] 
\[g(x)= 0{,}3x-3,\] 
\[h (x)=-0{,}4x+5,\]
\[k (x)=-3x+2,\] 
\[m (x)=12x+4\] 
sú rastúce?}
{\(g\), \(m\)}{\(h\), \(k\)}{\(f\), \(g\), \(m\)}{\(f\), \(k\), \(m\)}{}{}}
\LANGes{
\Question{
¿Cuál de las funciones
\[f (x)=5,\] 
\[g(x)= 0.3x-3,\] 
\[h (x)=-0.4x+5,\]
\[k (x)=-3x+2,\] 
\[m (x)=12x+4\] 
es creciente?}
{\(g\), \(m\)}{\(h\), \(k\)}{\(f\), \(g\), \(m\)}{\(f\), \(k\), \(m\)}{}{}}
\End

\Data\ID{2000003102}
\Updated{2022-01-17 06:01:33}
\Level{A}
\SubArea{Linear functions}
\LANGen{
\Question{
Which of the functions \(f\), \(g\), \(h\), \(k\), \(m\) are increasing as well as bounded functions?
\[f(x)=5,~x\in [ 0; \infty)\]

\[g (x)=0.3x-3,~x\in [ 0;6 ]\]

\[h (x)=-0.4+5,~x\in (-\infty; 3]\]

\[k (x)=3x+2,~x\in [ -3;5)\]

\[m (x)=12x+4,~x\in [ 0; \infty)\]}
{\(g\), \(k\)}{\(f\), \(g\), \(k\), \(m\)}{\(g\), \(k\), \(m\)}{\(f\), \(h\)}{}{}}
\LANGpl{
\Question{
Które z podanych funkcji \(f\), \(g\), \(h\), \(k\), \(m\) są funkcjami zarówno rosnącymi, jak i ograniczonymi?
\[f(x)=5,~x\in \langle 0; \infty)\]

\[g (x)=0{,}3x-3,~x\in \langle 0;6 \rangle\]

\[h (x)=-0{,}4+5,~x\in (-\infty; 3\rangle\]

\[k (x)=3x+2,~x\in \langle -3;5)\]

\[m (x)=12x+4,~x\in \langle 0; \infty)\]}
{\(g\), \(k\)}{\(f\), \(g\), \(k\), \(m\)}{\(g\), \(k\), \(m\)}{\(f\), \(h\)}{}{}}
\LANGcs{
\Question{
Které z funkcí \(f\), \(g\), \(h\), \(k\), \(m\) jsou rostoucí a zároveň ohraničené?
\[f(x)=5,~x\in \langle 0; \infty)\]

\[g (x)=0{,}3x-3,~x\in \langle 0;6 \rangle\]

\[h (x)=-0{,}4+5,~x\in (-\infty; 3\rangle\]

\[k (x)=3x+2,~x\in \langle -3;5)\]

\[m (x)=12x+4,~x\in \langle 0; \infty)\]}
{\(g\), \(k\)}{\(f\), \(g\), \(k\), \(m\)}{\(g\), \(k\), \(m\)}{\(f\), \(h\)}{}{}}
\LANGsk{
\Question{
Ktoré z funkcií  \(f\), \(g\), \(h\), \(k\), \(m\) sú rastúce a sčasne ohraničené?
\[f(x)=5,~x\in \langle 0; \infty)\]

\[g (x)=0{,}3x-3,~x\in \langle 0;6 \rangle\]

\[h (x)=-0{,}4+5,~x\in (-\infty; 3\rangle\]

\[k (x)=3x+2,~x\in \langle -3;5)\]

\[m (x)=12x+4,~x\in \langle 0; \infty)\]}
{\(g\), \(k\)}{\(f\), \(g\), \(k\), \(m\)}{\(g\), \(k\), \(m\)}{\(f\), \(h\)}{}{}}
\LANGes{
\Question{
¿Cuáles de las funciones \(f\), \(g\), \(h\), \(k\), \(m\) son crecientes y acotadas a la vez?
\[f(x)=5,~x\in [ 0; \infty)\]

\[g (x)=0.3x-3,~x\in [ 0;6 ]\]

\[h (x)=-0.4+5,~x\in (-\infty; 3]\]

\[k (x)=3x+2,~x\in [ -3;5)\]

\[m (x)=12x+4,~x\in [ 0; \infty)\]}
{\(g\), \(k\)}{\(f\), \(g\), \(k\), \(m\)}{\(g\), \(k\), \(m\)}{\(f\), \(h\)}{}{}}
\End

\Data\ID{2000003103}
\Updated{2022-01-17 06:01:55}
\Level{A}
\SubArea{Linear functions}
\LANGen{
\Question{
Which of the functions \(f\), \(g\), \(h\), \(k\), \(m\), \(n\) are decreasing, have minimum and are bounded functions?
\[f (x)=-3,~x \in \mathbb{R}\]

\[g (x)=-0.3x-3,~x \in [ 0;6 ]\]

\[h (x)=-0.4x+5,~x \in (-\infty ;3 ]\]

\[k (x)=3x+2,~x \in [ -3;5)\]

\[m (x)=-12x+4,~x \in [ 0;\infty)\]

\[n (x)=-2x+4,~x \in (0;7 ]\]}
{\(g\), \(n\)}{\(f\), \(g\), \(h\), \(m\), \(n\)}{\(g\)}{\(k\), \(n\)}{}{}}
\LANGpl{
\Question{
Które z podanych funkcji \(f\), \(g\), \(h\), \(k\), \(m\), \(n\) są funkcjami malejącymi, ograniczonymi i mają minium?
\[f (x)=-3,~x \in \mathbb{R}\]

\[g (x)=-0{,}3x-3,~x \in \langle 0;6 \rangle\]

\[h (x)=-0{,}4x+5,~x \in (-\infty ;3 \rangle\]

\[k (x)=3x+2,~x \in \langle -3;5)\]

\[m (x)=-12x+4,~x \in \langle 0;\infty)\]

\[n (x)=-2x+4,~x \in (0;7 \rangle\]}
{\(g\), \(n\)}{\(f\), \(g\), \(h\), \(m\), \(n\)}{\(g\)}{\(k\), \(n\)}{}{}}
\LANGcs{
\Question{
Které z funkcí \(f\), \(g\), \(h\), \(k\), \(m\), \(n\) jsou klesající, mají minimum a jsou ohraničené?
\[f (x)=-3,~x \in \mathbb{R}\]

\[g (x)=-0{,}3x-3,~x \in \langle 0;6 \rangle\]

\[h (x)=-0{,}4x+5,~x \in (-\infty ;3 \rangle\]

\[k (x)=3x+2,~x \in \langle -3;5)\]

\[m (x)=-12x+4,~x \in \langle 0;\infty)\]

\[n (x)=-2x+4,~x \in (0;7 \rangle\]}
{\(g\), \(n\)}{\(f\), \(g\), \(h\), \(m\), \(n\)}{\(g\)}{\(k\), \(n\)}{}{}}
\LANGsk{
\Question{
Ktoré z funkcií  \(f\), \(g\), \(h\), \(k\), \(m\), \(n\) sú klesajúce, majú minimum a sú ohraničené?
\[f (x)=-3,~x \in \mathbb{R}\]

\[g (x)=-0{,}3x-3,~x \in \langle 0;6 \rangle\]

\[h (x)=-0{,}4x+5,~x \in (-\infty ;3 \rangle\]

\[k (x)=3x+2,~x \in \langle -3;5)\]

\[m (x)=-12x+4,~x \in \langle 0;\infty)\]

\[n (x)=-2x+4,~x \in (0;7 \rangle\]}
{\(g\), \(n\)}{\(f\), \(g\), \(h\), \(m\), \(n\)}{\(g\)}{\(k\), \(n\)}{}{}}
\LANGes{
\Question{
¿Cuáles de las funciones  \(f\), \(g\), \(h\), \(k\), \(m\), \(n\) son decrecientes, acotadas y tienen mínimo?
\[f (x)=-3,~x \in \mathbb{R}\]

\[g (x)=-0.3x-3,~x \in [ 0;6 ]\]

\[h (x)=-0.4x+5,~x \in (-\infty ;3 ]\]

\[k (x)=3x+2,~x \in [ -3;5)\]

\[m (x)=-12x+4,~x \in [ 0;\infty)\]

\[n (x)=-2x+4,~x \in (0;7 ]\]}
{\(g\), \(n\)}{\(f\), \(g\), \(h\), \(m\), \(n\)}{\(g\)}{\(k\), \(n\)}{}{}}
\End

\Data\ID{2000003104}
\Updated{2021-12-01 04:12:38}
\Level{A}
\SubArea{Linear functions}
\LANGen{
\Question{
Find range of the function \(m (x)= 7x+1,~x \in (0;7]\).}
{\(H(m) = (1;50 ] \)}{\(H(m) = [ 1;50 ) \)}{\(H(m) =  [ 1;50 ] \)}{\(H(m) = (0;7 ] \)}{}{}}
\LANGpl{
\Question{
Znajdź zbiór wartości funkcji \(m (x)= 7x+1,~x \in (0;7\rangle\).}
{\(H(m) = (1;50 \rangle \)}{\(H(m) = \langle 1;50 ) \)}{\(H(m) =  \langle 1;50 \rangle \)}{\(H(m) = (0;7 \rangle \)}{}{}}
\LANGcs{
\Question{
Jaký je obor hodnot funkce \(m (x)= 7x+1,~x \in (0;7\rangle\)?}
{\(H(m) = (1;50 \rangle \)}{\(H(m) = \langle 1;50 ) \)}{\(H(m) =  \langle 1;50 \rangle \)}{\(H(m) = (0;7 \rangle \)}{}{}}
\LANGsk{
\Question{
Nájdite obor hodnôt funkcie  \(m (x)= 7x+1,~x \in (0;7\rangle\).}
{\(H(m) = (1;50 \rangle \)}{\(H(m) = \langle 1;50 ) \)}{\(H(m) =  \langle 1;50 \rangle \)}{\(H(m) = (0;7 \rangle \)}{}{}}
\LANGes{
\Question{
¿Cuál es el Rango de la función \(m (x)= 7x+1,~x \in (0;7]\)?}
{\(H(m) = (1;50 ] \)}{\(H(m) = [ 1;50 ) \)}{\(H(m) =  [ 1;50 ] \)}{\(H(m) = (0;7 ] \)}{}{}}
\End

\Data\ID{2000003105}
\Updated{2022-01-17 06:01:17}
\Level{A}
\SubArea{Linear functions}
\IMG{2000003101_10-figure0.pdf}
\LANGen{
\Question{
There is a graph of a linear function in the picture. Find the formula (equation) for this function.}
{\(y=7x\)}{\( y=x+7\)}{\( y=7x+1\)}{\(7y=x\)}{}{}}
\LANGpl{
\Question{
Dany jest wykres funkcji liniowej. Znajdź wzór dla tej funkcji.}
{\(y=7x\)}{\( y=x+7\)}{\( y=7x+1\)}{\(7y=x\)}{}{}}
\LANGcs{
\Question{
Na obrázku je dán graf lineární funkce. Určete její rovnici.}
{\(y=7x\)}{\( y=x+7\)}{\( y=7x+1\)}{\(7y=x\)}{}{}}
\LANGsk{
\Question{
Na obrázku je graf lineárnej funkcie. Nájdite predpis tejto funkcie.}
{\(y=7x\)}{\( y=x+7\)}{\( y=7x+1\)}{\(7y=x\)}{}{}}
\LANGes{
\Question{
En el dibujo aparece la gráfica de una función lineal. Averigua su ecuación.}
{\(y=7x\)}{\( y=x+7\)}{\( y=7x+1\)}{\(7y=x\)}{}{}}
\End

\Data\ID{2000003106}
\Updated{2022-01-17 06:01:34}
\Level{A}
\SubArea{Linear functions}
\IMG{2000003101_10-figure1.pdf}
\LANGen{
\Question{
There is a graph of a linear function in the picture. Find the formula (equation) for this function.}
{\(y=-x+2\)}{\( y=2x+4\)}{\( y=-2x+4\)}{\(y=-2x+2\)}{}{}}
\LANGpl{
\Question{
Dany jest wykres funkcji liniowej. Znajdź wzór dla tej funkcji.}
{\(y=-x+2\)}{\( y=2x+4\)}{\( y=-2x+4\)}{\(y=-2x+2\)}{}{}}
\LANGcs{
\Question{
Na obrázku je dán graf lineární funkce. Určete její rovnici.}
{\(y=-x+2\)}{\( y=2x+4\)}{\( y=-2x+4\)}{\(y=-2x+2\)}{}{}}
\LANGsk{
\Question{
Na obrázku je graf lineárnej funkcie. Nájdite predpis tejto funkcie.}
{\(y=-x+2\)}{\( y=2x+4\)}{\( y=-2x+4\)}{\(y=-2x+2\)}{}{}}
\LANGes{
\Question{
En el dibujo está la gráfica de una función. Elige su ecuación.}
{\(y=-x+2\)}{\( y=2x+4\)}{\( y=-2x+4\)}{\(y=-2x+2\)}{}{}}
\End

\Data\ID{2000003107}
\Updated{2022-01-17 06:01:17}
\Level{A}
\SubArea{Linear functions}
\IMG{2000003101fig234_0.pdf}
\LANGen{
\Question{
There are three graphs of the linear functions in the picture. Choose the formula (equation) of the function which is not in the picture.}
{\(y=-2x+2\)}{\(y=-x+2\)}{\(y=2x\)}{\(y=2x+2\)}{}{}}
\LANGpl{
\Question{
Dane są trzy wykresy funkcji liniowych. Wybierz wzór funkcji, która nie jest przedstawiona na żadnym wykresie.}
{\(y=-2x+2\)}{\(y=-x+2\)}{\(y=2x\)}{\(y=2x+2\)}{}{}}
\LANGcs{
\Question{
Na obrázku jsou tři grafy lineárních funkcí. Určete rovnici náležící funkci, která na obrázku není.}
{\(y=-2x+2\)}{\(y=-x+2\)}{\(y=2x\)}{\(y=2x+2\)}{}{}}
\LANGsk{
\Question{
Na obrázku sú grafy troch lineárnych funkcií. Vyberte predpis funkcie, ktorej graf na obrázku nie je.}
{\(y=-2x+2\)}{\(y=-x+2\)}{\(y=2x\)}{\(y=2x+2\)}{}{}}
\LANGes{
\Question{
En el dibujo aparecen las gráficas de tres funciones lineales. Elige la ecuación que describe una función cuya gráfica no está en el dibujo.}
{\(y=-2x+2\)}{\(y=-x+2\)}{\(y=2x\)}{\(y=2x+2\)}{}{}}
\End

\Data\ID{2000003108}
\Updated{2022-01-17 06:01:20}
\Level{A}
\SubArea{Linear functions}
\IMG{2000003101_10-figure5_0.pdf}
\LANGen{
\Question{
There is a graph of a function \(f\) in the picture. The function \(f\) is restriction of a linear function such that the domain of \(f\) is \([ -2;\infty)\). What are properties of \(f\)?}
{The function \(f\) is bounded above, injective (one-to-one) and decreasing.}{The function \(f\) has maximum and minimum, is decreasing and bounded.}{The function \(f\) is odd, decreasing and has maximum.}{The function \(f\) does not have minimum, is even and bounded above.}{}{}}
\LANGpl{
\Question{
Dany jest wykres funkcji  \(f\).  Funkcja \(f\) jest takim ograniczeniem funkcji liniowej, że dziedziną funkcji  \(f\) jest przedział \(\langle -2;\infty)\). Jakie własności ma funkcja  \(f\)?}
{Funkcja \(f\) jest ograniczona od góry, jest iniekcją i malejąca.}{Funkcja \(f\) ma maximum i minimum, jest malejąca i ograniczona.}{Funkcja \(f\) jest nieparzysta, malejąca i ma maximum.}{Funkcja \(f\) nie ma minimum, jest stała i ograniczona od góry.}{}{}}
\LANGcs{
\Question{
Na obrázku je graf lineární funkce \(f\). Definiční obor funkce \(f\) je \(\langle -2;\infty)\). Jaké jsou vlastnosti této funkce \(f\)?}
{Funkce \(f\) je omezená shora, prostá a klesající.}{Funkce \(f\) má maximum a minimum, je klesající a ohraničená.}{Funkce \(f\) je lichá, klesající a má maximum.}{Funkce \(f\) nemá minimum, je sudá a shora ohraničená.}{}{}}
\LANGsk{
\Question{
Na obrázku je graf funkcie  \(f\) . Definičným oborom funkcie  \(f\) je interval \(\langle -2;\infty)\). Aké vlastnosti má funkcia \(f\)?}
{Funkcia \(f\) je zhora ohraničená, prostá, klesajúca.}{Funkcia \(f\) má maximum aj minimum, je klesajúca a ohraničená.}{Funkcia \(f\) je nepárna, klesajúca a má maximum.}{Funkcia \(f\) nemá minimum, je párna a zhora ohraničená.}{}{}}
\LANGes{
\Question{
En el dibujo se muestra la gráfica de la funcion lineal \(f\). El Dominio de la función \(f\) es \([ -2;\infty)\). ¿Qué propiedades tiene la función \(f\)?}
{La función \(f\) es acotada superiormente, decreciente e inyectiva.}{La función \(f\) tiene máximo y mínimo, es decreciente y acotada.}{La función \(f\) es impar, decreciente y tiene máximo.}{La función \(f\) no tiene mánimo, es par y acotada superiormente.}{}{}}
\End

\Data\ID{2010009301}
\Updated{2022-11-02 09:11:27}
\Level{A}
\SubArea{Linear functions}
\IMG{2010009301-figure0.pdf}
\LANGen{
\Question{
The linear function \(g\) 
is graphed in the picture. Find the analytic expression for the function
\(g\).}
{\(y = \frac{1}
{4}x\)}{\(y = -\frac{1}
{4}x\)}{\(y = 4x\)}{\(y =-4x\)}{}{}}
\LANGpl{
\Question{
Na rysunku przedstawiono funkcję liniową \(g\). Podaj wzór funkcji \(g\).}
{\(y = \frac{1}
{4}x\)}{\(y = -\frac{1}
{4}x\)}{\(y = 4x\)}{\(y =-4x\)}{}{}}
\LANGcs{
\Question{
Na obrázku je znázorněný graf lineární funkce \(g\). Jaký je předpis této funkce
\(g\)?}
{\(y = \frac{1}
{4}x\)}{\(y = -\frac{1}
{4}x\)}{\(y = 4x\)}{\(y =-4x\)}{}{}}
\LANGsk{
\Question{
Na obrázku je znázornený graf funkcie  \(g\).
Nájdite predpis tejto funkcie.}
{\(y = \frac{1}
{4}x\)}{\(y = -\frac{1}
{4}x\)}{\(y = 4x\)}{\(y =-4x\)}{}{}}
\LANGes{
\Question{
La imagen muestra la gráfica de la función lineal \(g\). Halla su expresión analítica.}
{\(y = \frac{1}
{4}x\)}{\(y = -\frac{1}
{4}x\)}{\(y = 4x\)}{\(y =-4x\)}{}{}}
\End

\Data\ID{2010009302}
\Updated{2022-11-02 09:11:27}
\Level{A}
\SubArea{Linear functions}
\LANGen{
\Question{
Given the linear function \(f(x) = -5x + 3\),
evaluate \(f(-2) + f(2)\).}
{\( 6\)}{\( -14\)}{\( 0\)}{\( -20\)}{}{}}
\LANGpl{
\Question{
Dana jest funkcja \(f(x) = -5x + 3\),
oblicz \(f(-2) + f(2)\).}
{\( 6\)}{\( -14\)}{\( 0\)}{\( -20\)}{}{}}
\LANGcs{
\Question{
Je dána lineární funkce \(f(x) = -5x + 3\).
Vypočtěte \(f(-2) + f(2)\).}
{\( 6\)}{\( -14\)}{\( 0\)}{\( -20\)}{}{}}
\LANGsk{
\Question{
Daná je lineárna funkcia  \(f(x) = -5x + 3\),
nájdite  \(f(-2) + f(2)\).}
{\( 6\)}{\( -14\)}{\( 0\)}{\( -20\)}{}{}}
\LANGes{
\Question{
Dada la función lineal \(f(x) = -5x + 3\),
evalúa \(f(-2) + f(2)\).}
{\( 6\)}{\( -14\)}{\( 0\)}{\( -20\)}{}{}}
\End

\Data\ID{2010009303}
\Updated{2022-11-02 09:11:27}
\Level{A}
\SubArea{Linear functions}
\LANGen{
\Question{
Let the function \(g\) 
be defined as a linear function with graph passing through the points
\(A = [-3;2]\) and
\(B = [-2;4]\). Find an analytic
expression for the function \(g\).}
{\(g(x)= 2x + 8\)}{\(g(x)= \frac12 x -4\)}{\(g(x)= -\frac74 x + \frac12\)}{\(g(x)= 2x -4\)}{}{}}
\LANGpl{
\Question{
Funkcja \(g\) to funkcja liniowa z wykresem przechodzącym przez punkty \(A = [-3;2]\) i
\(B = [-2;4]\). Podaj wzór funkcji \(g\).}
{\(g(x)= 2x + 8\)}{\(g(x)= \frac12 x -4\)}{\(g(x)= -\frac74 x + \frac12\)}{\(g(x)= 2x -4\)}{}{}}
\LANGcs{
\Question{
Graf lineární funkce \(g\) 
prochází body 
\(A = [-3;2]\) a
\(B = [-2;4]\). Jaký je předpis této funkce \(g\)?}
{\(g(x)= 2x + 8\)}{\(g(x)= \frac12 x -4\)}{\(g(x)= -\frac74 x + \frac12\)}{\(g(x)= 2x -4\)}{}{}}
\LANGsk{
\Question{
Graf lineárnej funkcie \(g\) 
prechádza bodmi 
\(A = [-3;2]\) a
\(B = [-2;4]\). Nájdite predpis danej funkcie \(g\).}
{\(g(x)= 2x + 8\)}{\(g(x)= \frac12 x -4\)}{\(g(x)= -\frac74 x + \frac12\)}{\(g(x)= 2x -4\)}{}{}}
\LANGes{
\Question{
Sea la función \(g\) 
definida como una función lineal cuya gráfica pasa por los puntos
\(A = [-3;2]\) y
\(B = [-2;4]\). Halla la expresión analítica de la función \(g\).}
{\(g(x)= 2x + 8\)}{\(g(x)= \frac12 x -4\)}{\(g(x)= -\frac74 x + \frac12\)}{\(g(x)= 2x -4\)}{}{}}
\End

\Data\ID{2010009304}
\Updated{2022-11-02 09:11:27}
\Level{A}
\SubArea{Linear functions}
\LANGen{
\Question{
Consider the linear function \(f(x)= -\frac{2}
{5}x + 3\).
Find the intersection point of the graph of
\(f\) with
\(x\)-axis.}
{\(\left[\frac{15}2;0\right]\)}{\(\left[-\frac{15}2;0\right]\)}{\([0;3]\)}{\([13;0]\)}{}{}}
\LANGpl{
\Question{
Dana jest funckja \(f(x)= -\frac{2}
{5}x + 3\).
Znajdź punkt przecięcia wykresu
\(f\) z osią
\(x\).}
{\(\left[\frac{15}2;0\right]\)}{\(\left[-\frac{15}2;0\right]\)}{\([0;3]\)}{\([13;0]\)}{}{}}
\LANGcs{
\Question{
Je dána lineární funkce \(f(x)= -\frac{2}
{5}x + 3\).
Najděte průsečík grafu funkce
\(f\) s osou
\(x\).}
{\(\left[\frac{15}2;0\right]\)}{\(\left[-\frac{15}2;0\right]\)}{\([0;3]\)}{\([13;0]\)}{}{}}
\LANGsk{
\Question{
Daná je lineárna funkcia  \(f(x)= -\frac{2}
{5}x + 3\).
Nájdite priesečník  funkcie 
\(f\) s osou
\(x\).}
{\(\left[\frac{15}2;0\right]\)}{\(\left[-\frac{15}2;0\right]\)}{\([0;3]\)}{\([13;0]\)}{}{}}
\LANGes{
\Question{
Sea la función lineal \(f(x)= -\frac{2}
{5}x + 3\).
Halla el punto de intersección de la gráfica de 
\(f\) con el eje
\(x\).}
{\(\left[\frac{15}2;0\right]\)}{\(\left[-\frac{15}2;0\right]\)}{\([0;3]\)}{\([13;0]\)}{}{}}
\End

\Data\ID{2010009305}
\Updated{2022-11-02 09:11:27}
\Level{A}
\SubArea{Linear functions}
\LANGen{
\Question{
Consider the linear function \(f(x) = -3x + 9\).
Find the intersection point of the graph of
\(f\) with
\(y\)-axis.}
{\([0;9]\)}{\([9;0]\)}{\([0;3]\)}{\([3;0]\)}{}{}}
\LANGpl{
\Question{
Dana jest funkcja \(f(x) = -3x + 9\).
Znajdź punkt przecięcia wykresu
\(f\) z osią
\(y\).}
{\([0;9]\)}{\([9;0]\)}{\([0;3]\)}{\([3;0]\)}{}{}}
\LANGcs{
\Question{
Je dána lineární funkce \(f(x) = -3x + 9\).
Najděte průsečík grafu funkce
\(f\) s osou
\(y\).}
{\([0;9]\)}{\([9;0]\)}{\([0;3]\)}{\([3;0]\)}{}{}}
\LANGsk{
\Question{
Daná je lineárna funkcia  \(f(x) = -3x + 9\).
Nájdite priesečník funkcie 
\(f\) s osou
\(y\).}
{\([0;9]\)}{\([9;0]\)}{\([0;3]\)}{\([3;0]\)}{}{}}
\LANGes{
\Question{
Sea la función lineal \(f(x) = -3x + 9\).
Halla el punto de intersección de la gráfica de
\(f\) con el eje
\(y\).}
{\([0;9]\)}{\([9;0]\)}{\([0;3]\)}{\([3;0]\)}{}{}}
\End

\Data\ID{2010009306}
\Updated{2022-11-02 09:11:27}
\Level{A}
\SubArea{Linear functions}
\LANGen{
\Question{
Given the linear function \(f(x) = -2x + 1\),
evaluate 
\[ 
 f(a) + f(a-1).
\]}
{\(- 4a +4\)}{\(- 4a +3\)}{\(4\)}{\(- 4a +2\)}{}{}}
\LANGpl{
\Question{
Dana jest funkcja \(f(x) = -2x + 1\). Oblicz 
\[ 
 f(a) + f(a-1).
\]}
{\(- 4a +4\)}{\(- 4a +3\)}{\(4\)}{\(- 4a +2\)}{}{}}
\LANGcs{
\Question{
Je dána lineární funkce \(f(x) = -2x + 1\).
Vypočtěte 
\[ 
 f(a) + f(a-1).
\]}
{\(- 4a +4\)}{\(- 4a +3\)}{\(4\)}{\(- 4a +2\)}{}{}}
\LANGsk{
\Question{
Daná je lineárna funkcia  \(f(x) = -2x + 1\).
Nájdite hodnotu
\[ 
 f(a) + f(a-1).
\]}
{\(- 4a +4\)}{\(- 4a +3\)}{\(4\)}{\(- 4a +2\)}{}{}}
\LANGes{
\Question{
Dada la función lineal \(f(x) = -2x + 1\),
evalúa 
\[ 
 f(a) + f(a-1).
\]}
{\(- 4a +4\)}{\(- 4a +3\)}{\(4\)}{\(- 4a +2\)}{}{}}
\End

\Data\ID{2010016601}
\Updated{2022-11-02 09:11:38}
\Level{A}
\SubArea{Linear functions}
\LANGen{
\Question{
Given the linear function \( f(x)=kx+3\), find the value of \( k \) so that \( f(6)= 15 \).}
{\( k=2 \)}{\( k=\frac15 \)}{\( k=3 \)}{\( k=5 \)}{}{}}
\LANGpl{
\Question{
Dana jest funkcja liniowa \( f(x)=kx+3\), znajdź wartość \( k \) tak, aby \( f(6)= 15 \).}
{\( k=2 \)}{\( k=\frac15 \)}{\( k=3 \)}{\( k=5 \)}{}{}}
\LANGcs{
\Question{
Je dána lineární funkce \( f(x)=kx+3\). Jaká je hodnota \( k \), pokud platí, že \( f(6)= 15 \)?}
{\( k=2 \)}{\( k=\frac15 \)}{\( k=3 \)}{\( k=5 \)}{}{}}
\LANGsk{
\Question{
Daná je lineárna funkcia \( f(x)=kx+3\). Nájdite hodnotu \( k \), ak platí  \( f(6)= 15 \).}
{\( k=2 \)}{\( k=\frac15 \)}{\( k=3 \)}{\( k=5 \)}{}{}}
\LANGes{
\Question{
Dada la función lineal \( f(x)=kx+3\), halla el valor de \( k \) tal que \( f(6)= 15 \).}
{\( k=2 \)}{\( k=\frac15 \)}{\( k=3 \)}{\( k=5 \)}{}{}}
\End

\Data\ID{2010016604}
\Updated{2022-11-02 09:11:38}
\Level{A}
\SubArea{Linear functions}
\IMG{2010016601-figure4.pdf}
\LANGen{
\Question{
Determine whether the line drawn in the picture is the graph of a linear function of the variable \(x\). If so, find the formula for the function.}
{There is no graph of the linear function in the picture.}{\( y=-2\)}{\( x=-2\)}{\( y=2x\)}{}{}}
\LANGpl{
\Question{
Określ, czy linia narysowana na rysunku jest wykresem funkcji liniowej ze zmienną \(x\). Jeśli tak, podaj wzór funkcji.}
{Na rysunku nie ma wykresu funkcji liniowej.}{\( y=-2\)}{\( x=-2\)}{\( y=2x\)}{}{}}
\LANGcs{
\Question{
Zjistěte, zda přímka na obrázku je grafem lineární funkce proměnné \(x\). Pokud ano, najděte předpis funkce.}
{Na obrázku není graf lineární funkce.}{\( y=-2\)}{\( x=-2\)}{\( y=2x\)}{}{}}
\LANGsk{
\Question{
Zistite, či priamka nakreslená na obrázku je grafom lineárnej funkcie s premennou \(x\). Ak áno, nájdite predpis  funkcie.}
{Na obrázku nie je graf lineárnej funkcie.}{\( y=-2\)}{\( x=-2\)}{\( y=2x\)}{}{}}
\LANGes{
\Question{
Determina si la recta dibujada en la imagen es la gráfica de una función lineal de la variable \(x\). Si es así, encuentra la fórmula de la función.}
{La imagen no muestra la gráfica de una función lineal.}{\( y=-2\)}{\( x=-2\)}{\( y=2x\)}{}{}}
\End

\Data\ID{2010016605}
\Updated{2022-11-02 09:11:38}
\Level{A}
\SubArea{Linear functions}
\IMG{2010016601-figure5.pdf}
\LANGen{
\Question{
Choose the formula of the function whose graph is shown in the picture.}
{\( f(x)=x+1;\ x\in [ -2;3)\)}{\( f(x)=x+1;\ x\in ( -2;3]\)}{\( f(x)=-x+1;\ x\in [ -2;3)\)}{\( f(x)=x-1;\ x\in [ -2;3)\)}{}{}}
\LANGpl{
\Question{
Wybierz wzór funkcji, której wykres jest pokazany na rysunku.}
{\( f(x)=x+1;\ x\in \langle -2;3)\)}{\( f(x)=x+1;\ x\in ( -2;3\rangle\)}{\( f(x)=-x+1;\ x\in \langle -2;3)\)}{\( f(x)=x-1;\ x\in \langle -2;3)\)}{}{}}
\LANGcs{
\Question{
Vyberte předpis funkce, jejíž graf je zakreslen na obrázku.}
{\( f(x)=x+1;\ x\in \langle -2;3)\)}{\( f(x)=x+1;\ x\in ( -2;3\rangle\)}{\( f(x)=-x+1;\ x\in \langle -2;3)\)}{\( f(x)=x-1;\ x\in \langle -2;3)\)}{}{}}
\LANGsk{
\Question{
Vyberte predpis funkcie, ktorej graf je zakreslený na obrázku.}
{\( f(x)=x+1;\ x\in \langle -2;3)\)}{\( f(x)=x+1;\ x\in ( -2;3\rangle\)}{\( f(x)=-x+1;\ x\in \langle -2;3)\)}{\( f(x)=x-1;\ x\in \langle -2;3)\)}{}{}}
\LANGes{
\Question{
Elige la fórmula de la función cuya gráfica hay en la imagen.}
{\( f(x)=x+1;\ x\in [ -2;3)\)}{\( f(x)=x+1;\ x\in ( -2;3]\)}{\( f(x)=-x+1;\ x\in [ -2;3)\)}{\( f(x)=x-1;\ x\in [ -2;3)\)}{}{}}
\End

\Data\ID{2010016607}
\Updated{2022-11-02 09:11:38}
\Level{A}
\SubArea{Linear functions}
\LANGen{
\Question{
Given the function \(f(x) = -2x +3\), where \(x\in (-3 ;4] \), determine the monotonicity of \(f\).}
{decreasing}{increasing}{constant}{nondecreasing}{}{}}
\LANGpl{
\Question{
Dana jest funkcja \(f(x) = -2x +3\), gdzie \(x\in (-3 ;4\rangle \). Wyznacz monotoniczność funkcji \(f\).}
{Funkcja jest malejąca.}{Funkcja jest rosnąca.}{Funkcja jest stała.}{Funkcja jest niemalejąca.}{}{}}
\LANGcs{
\Question{
Je dána funkce \(f(x) = -2x +3\), pro \(x\in (-3 ;4\rangle \). Jaká je monotónnost funkce \(f\)?}
{klesající}{rostoucí}{konstantní}{neklesající}{}{}}
\LANGsk{
\Question{
Daná je funkcia \(f(x) = -2x +3\), kde \(x\in (-3 ;4\rangle \). Určte monotónnosť funkcie  \(f\).}
{klesajúca}{rastúca}{konštantná}{neklesajúca}{}{}}
\LANGes{
\Question{
Dada la función \(f(x) = -2x +3\), donde \(x\in (-3 ;4] \), halla la monotonía de \(f\).}
{decreciente}{creciente}{constante}{no decreciente}{}{}}
\End

\Data\ID{2010016608}
\Updated{2022-11-02 09:11:38}
\Level{A}
\SubArea{Linear functions}
\LANGen{
\Question{
Given the function \(f(x) = 3x -4\), where \(x\in ( -4;3]  \), determine the monotonicity of \(f\).}
{increasing}{decreasing}{constant}{nonincreasing}{}{}}
\LANGpl{
\Question{
Dana jest funkcja \(f(x) = 3x -4\), gdzie \(x\in ( -4;3\rangle  \). Wyznacz monotoniczność funkcji \(f\).}
{Funkcja jest rosnąca.}{Funkcja jest malejąca.}{Funkcja jest stała.}{Funkcja jest nierosnąca.}{}{}}
\LANGcs{
\Question{
Je dána funkce \(f(x) = 3x -4\), kde \(x\in ( -4;3\rangle  \). Jaká je monotónnost funkce \(f\)?}
{rostoucí}{klesající}{konstantní}{nerostoucí}{}{}}
\LANGsk{
\Question{
Daná je funkcia \(f(x) = 3x -4\), kde \(x\in ( -4;3\rangle  \). Určte monotónnosť funkcie \(f\).}
{rastúca}{klesajúca}{konštantná}{nerastúca}{}{}}
\LANGes{
\Question{
Dada la función \(f(x) = 3x -4\), donde \(x\in ( -4;3]  \), halla la monotonía de \(f\).}
{creciente}{decreciente}{constante}{no creciente}{}{}}
\End

\Data\ID{2010016609}
\Updated{2022-11-02 09:11:38}
\Level{A}
\SubArea{Linear functions}
\LANGen{
\Question{
Given the function \(f(x) = 3x - 15\),
\(x\in \mathbb{R}\), find the
\(x\)-intercept.}
{\(x = 5\)}{\(x = -5\)}{\(x = \frac15\)}{\(x = -\frac15\)}{}{}}
\LANGpl{
\Question{
Dana jest funkcja \(f(x) = 3x - 15\),
gdzie \(x\in \mathbb{R}\). Znajdź punkt przecięcia 
\(x\).}
{\(x = 5\)}{\(x = -5\)}{\(x = \frac15\)}{\(x = -\frac15\)}{}{}}
\LANGcs{
\Question{
Je dána funkce \(f(x) = 3x - 15\),
\(x\in \mathbb{R}\). Najděte průsečík jejího grafu s osou
\(x\).}
{\(x = 5\)}{\(x = -5\)}{\(x = \frac15\)}{\(x = -\frac15\)}{}{}}
\LANGsk{
\Question{
Daná je funkcia  \(f(x) = 3x - 15\),
\(x\in \mathbb{R}\). Nájdite priesečník funkcie s osou 
\(x\).}
{\(x = 5\)}{\(x = -5\)}{\(x = \frac15\)}{\(x = -\frac15\)}{}{}}
\LANGes{
\Question{
Dada la función \(f(x) = 3x - 15\),
\(x\in \mathbb{R}\), halla el punto de intersección con el eje
\(x\).}
{\(x = 5\)}{\(x = -5\)}{\(x = \frac15\)}{\(x = -\frac15\)}{}{}}
\End

\Data\ID{2010016610}
\Updated{2022-11-02 09:11:38}
\Level{A}
\SubArea{Linear functions}
\LANGen{
\Question{
Given the function \(f(x) = -2x - 10\),
\(x\in (-\infty ;3] \),
solve
\[ 
 f(x) = -8.
\]}
{\(x = -1\)}{\(x = 6\)}{\(x = -26\)}{\(x = 9\)}{}{}}
\LANGpl{
\Question{
Dana jest funkcja \(f(x) = -2x - 10\)
gdzie \(x\in (-\infty ;3\rangle \).
Rozwiąż 
\[ 
 f(x) = -8.
\]}
{\(x = -1\)}{\(x = 6\)}{\(x = -26\)}{\(x = 9\)}{}{}}
\LANGcs{
\Question{
Je dána funkce \(f(x) = -2x - 10\),
\(x\in (-\infty ;3\rangle \).
Řešte
\[ 
 f(x) = -8.
\]}
{\(x = -1\)}{\(x = 6\)}{\(x = -26\)}{\(x = 9\)}{}{}}
\LANGsk{
\Question{
Daná je funkcia \(f(x) = -2x - 10\),
\(x\in (-\infty ;3\rangle \),
vyriešte
\[ 
 f(x) = -8.
\]}
{\(x = -1\)}{\(x = 6\)}{\(x = -26\)}{\(x = 9\)}{}{}}
\LANGes{
\Question{
Dada la función \(f(x) = -2x - 10\),
\(x\in (-\infty ;3] \),
resuelve
\[ 
 f(x) = -8.
\]}
{\(x = -1\)}{\(x = 6\)}{\(x = -26\)}{\(x = 9\)}{}{}}
\End

\Data\ID{2010016611}
\Updated{2022-11-02 09:11:38}
\Level{A}
\SubArea{Linear functions}
\IMG{2010016601-figure6.pdf}
\LANGen{
\Question{
The range of the function \(g\) graphed in the picture is \([ -1;\infty)\). Find the domain of \(g\).}
{\([ -5;\infty)\)}{\(\mathbb{R}\)}{\([ -1;\infty)\)}{\((-5;\infty)\)}{}{}}
\LANGpl{
\Question{
Zakres funkcji \(g\) przedstawionej na rysunku to \(\langle -1;\infty)\). Znajdź dziedzinę funkcji \(g\).}
{\(\langle -5;\infty)\)}{\(\mathbb{R}\)}{\(\langle -1;\infty)\)}{\((-5;\infty)\)}{}{}}
\LANGcs{
\Question{
Na obrázku je funkce \(g\) (část lineární funkce), jejíž obor hodnot je \(\langle -1;\infty)\). Jaký je definiční obor této funkce \(g\)?}
{\(\langle -5;\infty)\)}{\(\mathbb{R}\)}{\(\langle -1;\infty)\)}{\((-5;\infty)\)}{}{}}
\LANGsk{
\Question{
Na obrázku je funkcia \(g\), ktorej oborom hodnôt je \(\langle -1;\infty)\). Nájdite definičný obor funkcie \(g\).}
{\(\langle -5;\infty)\)}{\(\mathbb{R}\)}{\(\langle -1;\infty)\)}{\((-5;\infty)\)}{}{}}
\LANGes{
\Question{
El Rango de la función \(g\), cuya gráfica está en el dibujo, es \([ -1;\infty)\). Halla el dominio de \(g\).}
{\([ -5;\infty)\)}{\(\mathbb{R}\)}{\([ -1;\infty)\)}{\((-5;\infty)\)}{}{}}
\End

\Data\ID{2100000801}
\Updated{2021-12-01 03:12:52}
\Level{A}
\SubArea{Linear functions}
\LANGen{
\Question{
Consider \(f(x)= -x+3 \) a linear function. Which of the given graphs is the graph of \(f\)?}
{}{}{}{}{}{}}
\LANGpl{
\Question{
Dana jest funkcja liniowa \(f(x)= -x+3 \).  Który z podanych wykresów jest wykresem funkcji \(f\)?}
{}{}{}{}{}{}}
\LANGcs{
\Question{
Je dána lineární funkce \(f(x)= -x+3 \). Který z uvedených grafů je grafem této funkce \(f\)?}
{}{}{}{}{}{}}
\LANGsk{
\Question{
Daná je lineárna funkcia \(f(x)= -x+3 \). Ktorý z uvedených grafov je grafom funkcie \(f\)?}
{}{}{}{}{}{}}
\LANGes{
\Question{
Sea \(f(x)= -x+3 \) una función lineal. ¿Cuál de las gráficas es la de \(f\)?}
{}{}{}{}{}{}}

\IMGanswer{linfce-figure0_0.pdf}
\IMGanswer{linfce-figure1.pdf}
\IMGanswer{linfce-figure2.pdf}
\IMGanswer{linfce-figure3_0.pdf}\End

\Data\ID{2110016603}
\Updated{2022-11-02 09:11:00}
\Level{A}
\SubArea{Linear functions}
\LANGen{
\Question{
Consider the linear function \( f(x)=3x-6 \). Which of the given graphs is the graph of \( f \)?}
{}{}{}{}{}{}}
\LANGpl{
\Question{
Dana jest funkcja liniowa \( f(x)=3x-6 \). Wskaż wykres funkcji \( f \).}
{}{}{}{}{}{}}
\LANGcs{
\Question{
Je dána lineární funkce \( f(x)=3x-6 \). Který z grafů je grafem funkce \( f \)?}
{}{}{}{}{}{}}
\LANGsk{
\Question{
Daná je lineárna funkcia  \( f(x)=3x-6 \). Ktorý z uvedených grafov je grafom funkcie \( f \)?}
{}{}{}{}{}{}}
\LANGes{
\Question{
Consideramos la función lineal \( f(x)=3x-6 \). ¿Cuál de las siguientes gráficas corresponde a la función \( f \)?}
{}{}{}{}{}{}}

\IMGanswer{2010016601-figure0.pdf}
\IMGanswer{2010016601-figure1.pdf}
\IMGanswer{2010016601-figure2.pdf}
\IMGanswer{2010016601-figure3.pdf}\End

\Data\ID{9000004201}
\Updated{2022-04-23 05:04:11}
\Level{A}
\SubArea{Linear functions}
\LANGen{
\Question{
Consider the function \(f(x) = 3x - 6\),
\(x\in (-\infty ;3] \). Find the
range of \(f\).}
{\((-\infty ;3] \)}{\([ 3;\infty )\)}{\(\mathbb{R}\)}{\((-\infty ;3)\)}{}{}}
\LANGpl{
\Question{
Rozważ funkcję \(f\colon y = 3x - 6\),
\(x\in (-\infty ;3] \). Znajdź zakres \(f\).}
{\((-\infty ;3] \)}{\([ 3;\infty )\)}{\(\mathbb{R}\)}{\((-\infty ;3)\)}{}{}}
\LANGcs{
\Question{
Je dána funkce \(f\colon y = 3x - 6, \ x\in (-\infty ;3\rangle \).
Obor hodnot funkce \(f\) 
je:}
{\((-\infty ;3\rangle \)}{\(\langle 3;\infty )\)}{\(\mathbb{R}\)}{\((-\infty ;3)\)}{}{}}
\LANGsk{
\Question{
Daná je funkcia \(f\colon y = 3x - 6, \ x\in (-\infty ;3\rangle \).
Určte obor hodnôt funkcie \(f\).}
{\((-\infty ;3\rangle \)}{\(\langle 3;\infty )\)}{\(\mathbb{R}\)}{\((-\infty ;3)\)}{}{}}
\LANGes{
\Question{
Dada la función \(f(x) = 3x - 6\),
\(x\in (-\infty ;3] \). Halla el Rango de \(f\).}
{\((-\infty ;3] \)}{\([ 3;\infty )\)}{\(\mathbb{R}\)}{\((-\infty ;3)\)}{}{}}
\End

\Data\ID{9000004203}
\Updated{2022-04-23 05:04:11}
\Level{A}
\SubArea{Linear functions}
\LANGen{
\Question{
Given the function \(f(x) = 3x - 6\),
\(x\in (-\infty ;3] \), find the
\(x\)-intercept.}
{\(x = 2\)}{\(x = -2\)}{\(x = \frac{1} 
{2}\)}{\(x = -\frac{1}
{2}\)}{}{}}
\LANGpl{
\Question{
Dana jest funkcja \(f\colon y = 3x - 6\),
\(x\in (-\infty ;3] \). Znajdź punkt przecięcia prostej z osią współrzędnych
\(x\).}
{\(x = 2\)}{\(x = -2\)}{\(x = \frac{1} 
{2}\)}{\(x = -\frac{1}
{2}\)}{}{}}
\LANGcs{
\Question{
Je dána funkce \(f\colon y = 3x - 6,\ x\in (-\infty ;3\rangle \).
Průsečík grafu funkce \(f\) 
s osou \(x\) 
má souřadnice:}
{\([2;0]\)}{\([-2;0]\)}{\(\left [\frac{1}
{2};0\right ]\)}{\(\left [-\frac{1}
{2};0\right ]\)}{}{}}
\LANGsk{
\Question{
Daná je funkcia \(f\colon y = 3x - 6,\ x\in (-\infty ;3\rangle \).
Priesečník grafu funkcie \(f\) 
s osou \(x\) 
má súradnice:}
{\([2;0]\)}{\([-2;0]\)}{\(\left [\frac{1}
{2};0\right ]\)}{\(\left [-\frac{1}
{2};0\right ]\)}{}{}}
\LANGes{
\Question{
Dada la función \(f(x) = 3x - 6\),
\(x\in (-\infty ;3] \), halla su intersección con el eje
\(x\).}
{\(x = 2\)}{\(x = -2\)}{\(x = \frac{1} 
{2}\)}{\(x = -\frac{1}
{2}\)}{}{}}
\End

\Data\ID{9000004204}
\Updated{2020-09-11 12:09:26}
\Level{A}
\SubArea{Linear functions}
\LANGen{
\Question{
Given the function \(f(x)= 3x - 6\),
\(x\in (-\infty ;3] \), find the
\(y\)-intercept.}
{\(y = -6\)}{\(y = 6\)}{\(y = 2\)}{\(y = -2\)}{}{}}
\LANGpl{
\Question{
Dana jest funkcja \(f\colon y = 3x - 6\),
\(x\in (-\infty ;3] \). Znajdź punkt przecięcia prostej z osią współrzędnych 
\(y\).}
{\(y = -6\)}{\(y = 6\)}{\(y = 2\)}{\(y = -2\)}{}{}}
\LANGcs{
\Question{
Je dána funkce \(f\colon y = 3x - 6,\ x\in (-\infty ;3\rangle \).
Průsečík grafu funkce \(f\) 
s osou \(y\) 
má souřadnice:}
{\([0;-6]\)}{\([0;6]\)}{\([0;2]\)}{\([0;-2]\)}{}{}}
\LANGsk{
\Question{
Daná je funkcia \(f\colon y = 3x - 6,\ x\in (-\infty ;3\rangle \).
Priesečník grafu funkcie \(f\) 
s osou \(y\) 
má súradnice:}
{\([0;-6]\)}{\([0;6]\)}{\([0;2]\)}{\([0;-2]\)}{}{}}
\LANGes{
\Question{
Dada la función \(f(x)= 3x - 6\),
\(x\in (-\infty ;3] \), halla su intersección con el eje 
\(y\).}
{\(y = -6\)}{\(y = 6\)}{\(y = 2\)}{\(y = -2\)}{}{}}
\End

\Data\ID{9000004205}
\Updated{2020-09-11 12:09:01}
\Level{A}
\SubArea{Linear functions}
\LANGen{
\Question{
Given the function \(f(x) = 3x - 6\),
\(x\in (-\infty ;3] \), find
\(f(-4)\).}
{\(- 18\)}{\(\frac{2}
{3}\)}{\(6\)}{\(- 6\)}{}{}}
\LANGpl{
\Question{
Dana jest funkcja \(f\colon y = 3x - 6\),
\(x\in (-\infty ;3] \). Znajdź
\(f(-4)\).}
{\(- 18\)}{\(\frac{2}
{3}\)}{\(6\)}{\(- 6\)}{}{}}
\LANGcs{
\Question{
Je dána funkce \(f\colon y = 3x - 6,\ x\in (-\infty ;3\rangle \).
Funkční hodnota funkce \(f\) 
v bodě \(- 4\) 
je:}
{\(- 18\)}{\(\frac{2}
{3}\)}{\(6\)}{\(- 6\)}{}{}}
\LANGsk{
\Question{
Daná je funkcia \(f\colon y = 3x - 6,\ x\in (-\infty ;3\rangle \).
Funkčná hodnota funkcie \(f\) 
v bode \(- 4\) 
je:}
{\(- 18\)}{\(\frac{2}
{3}\)}{\(6\)}{\(- 6\)}{}{}}
\LANGes{
\Question{
Dada la función \(f(x) = 3x - 6\),
\(x\in (-\infty ;3] \), halla
\(f(-4)\).}
{\(- 18\)}{\(\frac{2}
{3}\)}{\(6\)}{\(- 6\)}{}{}}
\End

\Data\ID{9000004206}
\Updated{2022-04-23 05:04:11}
\Level{A}
\SubArea{Linear functions}
\LANGen{
\Question{
Given the function \(f(x) = 3x - 6\),
\(x\in (-\infty ;3] \),
solve
\[ 
 f(x) = -8.
\]}
{\(x = -\frac{2}
{3}\)}{\(x = -\frac{3}
{2}\)}{\(x = -30\)}{\(x = -18\)}{}{}}
\LANGpl{
\Question{
Dana jest funkcja \(f\colon y = 3x - 6\),
\(x\in (-\infty ;3] \). Rozwiąż \[ 
 f(x) = -8.
\]}
{\(x = -\frac{2}
{3}\)}{\(x = -\frac{3}
{2}\)}{\(x = -30\)}{\(x = -18\)}{}{}}
\LANGcs{
\Question{
Je dána funkce \(f\colon y = 3x - 6,\ x\in (-\infty ;3\rangle \).
Hodnota funkce je \(- 8\) 
pro:}
{\(x = -\frac{2}
{3}\)}{\(x = -\frac{3}
{2}\)}{\(x = -30\)}{\(x = -18\)}{}{}}
\LANGsk{
\Question{
Daná je funkcia \(f\colon y = 3x - 6,\ x\in (-\infty ;3\rangle \).
Funkčná hodnota \( 
 f(x) = -8
\), nájdite \(x\).}
{\(x = -\frac{2}
{3}\)}{\(x = -\frac{3}
{2}\)}{\(x = -30\)}{\(x = -18\)}{}{}}
\LANGes{
\Question{
Dada la función \(f(x) = 3x - 6\),
\(x\in (-\infty ;3] \),
resuelve
\[ 
 f(x) = -8.
\]}
{\(x = -\frac{2}
{3}\)}{\(x = -\frac{3}
{2}\)}{\(x = -30\)}{\(x = -18\)}{}{}}
\End

\Data\ID{9000004207}
\Updated{2022-09-09 11:09:58}
\Level{A}
\SubArea{Linear functions}
\IMG{000042_07-figure0.pdf}
\LANGen{
\Question{
The range of the function \(g\) 
graphed in the picture is \((-\infty ;3] \). Find the domain of \(g\).}
{\([ - 2;\infty )\)}{\(\mathbb{R}\)}{\((-\infty ;3] \)}{\((-2;\infty )\)}{}{}}
\LANGpl{
\Question{
Zakres funkcji \(g\) 
przedstawiony na rysunku wynosi \((-\infty ;3] \). Znajdź dziedzinę  \(g\).}
{\([ - 2;\infty )\)}{\(\mathbb{R}\)}{\((-\infty ;3] \)}{\((-2;\infty )\)}{}{}}
\LANGcs{
\Question{
Na obrázku je dána funkce \(g\) (část lineární funkce), která má obor hodnot
\((-\infty ;3\rangle \). Určete definiční obor funkce \(g\).}
{\(\langle - 2;\infty )\)}{\(\mathbb{R}\)}{\((-\infty ;3\rangle \)}{\((-2;\infty )\)}{}{}}
\LANGsk{
\Question{
Obor hodnôt funkcie \(g\),
ktorej graf vidíme na obrázku, je \((-\infty ;3\rangle \).
Určte definičný obor funkcie \(g\).}
{\(\langle - 2;\infty )\)}{\(\mathbb{R}\)}{\((-\infty ;3\rangle \)}{\((-2;\infty )\)}{}{}}
\LANGes{
\Question{
El Rango de la función \(g\), cuya gráfica está en el dibujo, es \((-\infty ;3] \). Halla el Dominio de \(g\).}
{\([ - 2;\infty )\)}{\(\mathbb{R}\)}{\((-\infty ;3] \)}{\((-2;\infty )\)}{}{}}
\End

\Data\ID{9000004208}
\Updated{2020-09-11 12:09:57}
\Level{A}
\SubArea{Linear functions}
\IMG{000042_08-figure0.pdf}
\LANGen{
\Question{
The domain of the part of the linear function \(g\) 
graphed in the picture is \([ - 2;\infty )\).
Find the range of \(g\).}
{\([ - 1;\infty )\)}{\(\mathbb{R}\)}{\((-2;\infty )\)}{\((-1;\infty )\)}{}{}}
\LANGpl{
\Question{
Dziedziną funkcji liniowej  \(g\) 
przedstawionej na rysunku jest \([ - 2;\infty )\).
Znajdź zakres \(g\).}
{\([ - 1;\infty )\)}{\(\mathbb{R}\)}{\((-2;\infty )\)}{\((-1;\infty )\)}{}{}}
\LANGcs{
\Question{
Na obrázku je dána funkce \(g\) (část lineární funkce), která má definiční obor \(\langle - 2;\infty )\).
Určete obor hodnot funkce \(g\).}
{\(\langle - 1;\infty )\)}{\(\mathbb{R}\)}{\((-2;\infty )\)}{\((-1;\infty )\)}{}{}}
\LANGsk{
\Question{
Definičný obor funkcie \(g\),
ktorej graf vidíme na obrázku je
\(\langle - 2;\infty )\). Určte obor hodnôt funkcie \(g\).}
{\(\langle - 1;\infty )\)}{\(\mathbb{R}\)}{\((-2;\infty )\)}{\((-1;\infty )\)}{}{}}
\LANGes{
\Question{
El Dominio de la parte de la función \(g\), cuya gráfica está en el dibujo, es \([ - 2;\infty )\).
Halla el Rango de \(g\).}
{\([ - 1;\infty )\)}{\(\mathbb{R}\)}{\((-2;\infty )\)}{\((-1;\infty )\)}{}{}}
\End

\Data\ID{9000004209}
\Updated{2022-04-23 05:04:37}
\Level{A}
\SubArea{Linear functions}
\IMG{000042_09-figure0.pdf}
\LANGen{
\Question{
The linear function \(g\) 
is graphed in the picture. Find the analytic expression for the function
\(g\).}
{\(y = -\frac{3}
{2}x\)}{\(y = \frac{3} 
{2}x\)}{\(y = \frac{2} 
{3}x\)}{\(y = -\frac{2}
{3}x\)}{}{}}
\LANGpl{
\Question{
Wykres funkcji liniowej \(g\) 
przedstawiono na rysunku poniżej. Wyznacz wzór funkcji \(g\).}
{\(y = -\frac{3}
{2}x\)}{\(y = \frac{3} 
{2}x\)}{\(y = \frac{2} 
{3}x\)}{\(y = -\frac{2}
{3}x\)}{}{}}
\LANGcs{
\Question{
Lineární funkce \(g\),
jejíž graf vidíme na obrázku, je dána předpisem:}
{\(y = -\frac{3}
{2}x\)}{\(y = \frac{3} 
{2}x\)}{\(y = \frac{2} 
{3}x\)}{\(y = -\frac{2}
{3}x\)}{}{}}
\LANGsk{
\Question{
Určte predpis lineárnej funkcie \(g\),
ktorej graf vidíme na obrázku.}
{\(y = -\frac{3}
{2}x\)}{\(y = \frac{3} 
{2}x\)}{\(y = \frac{2} 
{3}x\)}{\(y = -\frac{2}
{3}x\)}{}{}}
\LANGes{
\Question{
En el dibujo se representa la gráfica de la función lineal \(g\) 
Halla la expresión analítica de
\(g\).}
{\(y = -\frac{3}
{2}x\)}{\(y = \frac{3} 
{2}x\)}{\(y = \frac{2} 
{3}x\)}{\(y = -\frac{2}
{3}x\)}{}{}}
\End

\Data\ID{9000004210}
\Updated{2021-01-27 05:01:42}
\Level{A}
\SubArea{Linear functions}
\IMG{000042_10-figure0.pdf}
\LANGen{
\Question{
The function \(g\) is
graphed in the picture. Find \(g(0)\).}
{\(0\)}{\(3\)}{\(- 2\)}{\(1\)}{}{}}
\LANGpl{
\Question{
Funkcja \(g\) została przedstawiona na rysunku. Znajdź \(g(0)\).}
{\(0\)}{\(3\)}{\(- 2\)}{\(1\)}{}{}}
\LANGcs{
\Question{
Funkce \(g\),
jejíž graf vidíme na obrázku, má funkční hodnotu v bodě
\(0\) rovnu
číslu:}
{\(0\)}{\(3\)}{\(- 2\)}{\(1\)}{}{}}
\LANGsk{
\Question{
Určte funkčnú hodnotu v bode \(0\)
funkcie \(g\), 
ktorej graf vidíme na obrázku.}
{\(0\)}{\(3\)}{\(- 2\)}{\(1\)}{}{}}
\LANGes{
\Question{
El dibujo representa la gráfica de la función \(g\).
Halla \(g(0)\).}
{\(0\)}{\(3\)}{\(- 2\)}{\(1\)}{}{}}
\End

\Data\ID{9000005701}
\Updated{2020-09-11 12:09:18}
\Level{A}
\SubArea{Linear functions}
\LANGen{
\Question{
Given the linear function \(f(x) = 3x - 2\),
evaluate \(f\left (\frac{1}
{6}\right )\).}
{\(-\frac{3} 
{2}\)}{\(- 1\)}{\(\frac{1}
{6}\)}{\(\frac{5}
{2}\)}{}{}}
\LANGpl{
\Question{
Dana jest funkcja liniowa \(f\colon y = 3x - 2\).
Oblicz \(f\left (\frac{1}
{6}\right )\).}
{\(-\frac{3} 
{2}\)}{\(- 1\)}{\(\frac{1}
{6}\)}{\(\frac{5}
{2}\)}{}{}}
\LANGcs{
\Question{
Je dána lineární funkce \(f\colon y = 3x - 2\).
Hodnota funkce \(f\) 
v bodě \(\frac{1}
{6}\) 
je rovna:}
{\(-\frac{3} 
{2}\)}{\(- 1\)}{\(\frac{1}
{6}\)}{\(\frac{5}
{2}\)}{}{}}
\LANGsk{
\Question{
Daná je lineárna funkcia \(f\colon y = 3x - 2\).
Funkčná hodnota funkcie \(f\) 
v bode \(\frac{1}
{6}\) 
je:}
{\(-\frac{3} 
{2}\)}{\(- 1\)}{\(\frac{1}
{6}\)}{\(\frac{5}
{2}\)}{}{}}
\LANGes{
\Question{
Dada la función lineal \(f(x) = 3x - 2\),
evalúa \(f\left (\frac{1}
{6}\right )\).}
{\(-\frac{3} 
{2}\)}{\(- 1\)}{\(\frac{1}
{6}\)}{\(\frac{5}
{2}\)}{}{}}
\End

\Data\ID{9000005702}
\Updated{2022-04-23 05:04:37}
\Level{A}
\SubArea{Linear functions}
\LANGen{
\Question{
Given the linear function \(f(x) = -2x + 3\),
evaluate \(f(2) + f(-2)\).}
{\(6\)}{\(0\)}{\(3\)}{\(- 8\)}{}{}}
\LANGpl{
\Question{
Dana jest funkcja liniowa \(f\colon y = -2x + 3\). 
Oblicz \(f(2) + f(-2)\).}
{\(6\)}{\(0\)}{\(3\)}{\(- 8\)}{}{}}
\LANGcs{
\Question{
Je dána lineární funkce \(f\colon y = -2x + 3\).
Hodnota \(f(2) + f(-2)\) 
je rovna:}
{\(6\)}{\(0\)}{\(3\)}{\(- 8\)}{}{}}
\LANGsk{
\Question{
Daná je lineárna funkcia \(f\colon y = -2x + 3\).
Určte hodnotu: \(f(2) + f(-2)\).}
{\(6\)}{\(0\)}{\(3\)}{\(- 8\)}{}{}}
\LANGes{
\Question{
Dada la función lineal \(f(x) = -2x + 3\),
evalúa \(f(2) + f(-2)\).}
{\(6\)}{\(0\)}{\(3\)}{\(- 8\)}{}{}}
\End

\Data\ID{9000005703}
\Updated{2020-09-11 12:09:14}
\Level{A}
\SubArea{Linear functions}
\LANGen{
\Question{
Given the linear function \(f(x)= \frac{1} 
{2}x - 2\),
evaluate \(f(-4) - f(4)\).}
{\(- 4\)}{\(- 6\)}{\(0\)}{\(4\)}{}{}}
\LANGpl{
\Question{
Dana jest funkcja liniowa \(f\colon y = \frac{1} 
{2}x - 2\).
Oblicz \(f(-4) - f(4)\).}
{\(- 4\)}{\(- 6\)}{\(0\)}{\(4\)}{}{}}
\LANGcs{
\Question{
Je dána lineární funkce \(f\colon y = \frac{1} 
{2}x - 2\).
Hodnota \(f(-4) - f(4)\) 
je rovna:}
{\(- 4\)}{\(- 6\)}{\(0\)}{\(4\)}{}{}}
\LANGsk{
\Question{
Daná je lineárna funkcia \(f\colon y = \frac{1} 
{2}x - 2\). Určte hodnotu: \(f(-4) - f(4)\).}
{\(- 4\)}{\(- 6\)}{\(0\)}{\(4\)}{}{}}
\LANGes{
\Question{
Dada la función lineal \(f(x)= \frac{1} 
{2}x - 2\),
evalúa \(f(-4) - f(4)\).}
{\(- 4\)}{\(- 6\)}{\(0\)}{\(4\)}{}{}}
\End

\Data\ID{9000005704}
\Updated{2020-09-11 12:09:51}
\Level{A}
\SubArea{Linear functions}
\LANGen{
\Question{
Given the linear function \(f(x) = 5x - 3\),
solve \(f(x) = -8\).}
{\(- 1\)}{\(- 43\)}{\(- 16\)}{\(11\)}{}{}}
\LANGpl{
\Question{
Dana jest funkcja liniowa  \(f\colon y = 5x - 3\). 
Oblicz \(f(x) = -8\).}
{\(- 1\)}{\(- 43\)}{\(- 16\)}{\(11\)}{}{}}
\LANGcs{
\Question{
Je dána lineární funkce \(f\colon y = 5x - 3\).
Určete, pro které \(x\in \mathbb{R}\) 
platí, že \(f(x) = -8\).}
{\(- 1\)}{\(- 43\)}{\(- 16\)}{\(11\)}{}{}}
\LANGsk{
\Question{
Daná je lineárna funkcia \(f\colon y = 5x - 3\).
Určte, pre ktoré \(x\in \mathbb{R}\) 
platí, že \(f(x) = -8\).}
{\(- 1\)}{\(- 43\)}{\(- 16\)}{\(11\)}{}{}}
\LANGes{
\Question{
Dada la función lineal \(f(x) = 5x - 3\),
resuelve \(f(x) = -8\).}
{\(- 1\)}{\(- 43\)}{\(- 16\)}{\(11\)}{}{}}
\End

\Data\ID{9000005705}
\Updated{2021-01-27 05:01:35}
\Level{A}
\SubArea{Linear functions}
\LANGen{
\Question{
Let \(f(x) = -\frac{1}
{2}x + a\). Which value of
the real parameter \(a\) in the
definition of the function \(f\) 
guarantees \(f(2) = 2\)?}
{\(3\)}{\(1\)}{\(- 4\)}{\(5\)}{}{}}
\LANGpl{
\Question{
Dana jest funkcja liniowa \(f\colon y = -\frac{1}
{2}x + a\). Jakie wartości parametru  \(a\) funkcji \(f\) 
spełniają \(f(2) = 2\)?}
{\(3\)}{\(1\)}{\(- 4\)}{\(5\)}{}{}}
\LANGcs{
\Question{
Je dána lineární funkce \(f\colon y = -\frac{1}
{2}x + a\).
Určete, pro které \(a\in \mathbb{R}\) 
platí, že \(f(2) = 2\).}
{\(3\)}{\(1\)}{\(- 4\)}{\(5\)}{}{}}
\LANGsk{
\Question{
Daná je lineárna funkcia \(f\colon y = -\frac{1}
{2}x + a\).
Určte, pre ktoré \(a\in \mathbb{R}\) 
platí, že \(f(2) = 2\).}
{\(3\)}{\(1\)}{\(- 4\)}{\(5\)}{}{}}
\LANGes{
\Question{
Sea \(f(x) = -\frac{1}
{2}x + a\). ¿Para qué valor del parámetro real \(a\) se cumple que \(f(2) = 2\)?}
{\(3\)}{\(1\)}{\(- 4\)}{\(5\)}{}{}}
\End

\Data\ID{9000005706}
\Updated{2022-04-23 05:04:37}
\Level{A}
\SubArea{Linear functions}
\LANGen{
\Question{
Let the function \(f\) 
be defined as a linear function with graph passing through the points
\(A = [2;3]\) and
\(B = [-1;6]\). Find an analytic
expression for the function \(f\).}
{\(f(x)= -x + 5\)}{\(f(x) = x + 1\)}{\(f(x)= 2x - 1\)}{\(f(x) = -5x + 1\)}{}{}}
\LANGpl{
\Question{
Dana jest funkcja liniowa \(f\), której wykres przechodzi przez punkty
\(A = [2;3]\) i
\(B = [-1;6]\). Wyznacz wzór funkcji \(f\).}
{\(f\colon y = -x + 5\)}{\(f\colon y = x + 1\)}{\(f\colon y = 2x - 1\)}{\(f\colon y = -5x + 1\)}{}{}}
\LANGcs{
\Question{
Funkční předpis lineární funkce
\(f\), jejíž graf
prochází body \(A = [2;3]\),
\(B = [-1;6]\), je:}
{\(f\colon y = -x + 5\)}{\(f\colon y = x + 1\)}{\(f\colon y = 2x - 1\)}{\(f\colon y = -5x + 1\)}{}{}}
\LANGsk{
\Question{
Graf lineárnej funkcie
\(f\), prechádza bodmi \(A = [2;3]\),
\(B = [-1;6]\). Určte predpis danej funkcie.}
{\(f\colon y = -x + 5\)}{\(f\colon y = x + 1\)}{\(f\colon y = 2x - 1\)}{\(f\colon y = -5x + 1\)}{}{}}
\LANGes{
\Question{
Sea una función lineal \(f\) 
que pasa por los puntos
\(A = (2;3)\) y
\(B = (-1;6)\). Halla la expresión analítica de la función \(f\).}
{\(f(x)= -x + 5\)}{\(f(x) = x + 1\)}{\(f(x)= 2x - 1\)}{\(f(x) = -5x + 1\)}{}{}}
\End

\Data\ID{9000005707}
\Updated{2020-09-18 04:09:59}
\Level{A}
\SubArea{Linear functions}
\LANGen{
\Question{
Let \(f\) be the linear
function \(f(x) = -x + 4\) restricted
to the interval \(x\in [ - 3;2] \).
Find the range of \(f\).}
{\([ 2;7] \)}{\([ 1;6] \)}{\([ - 3;3] \)}{\([ - 1;2] \)}{}{}}
\LANGpl{
\Question{
Dana jest funkcja liniowa
 \(f\colon y = -x + 4\) ograniczona w przedziale \(x\in [ - 3;2] \).
Wyznacz zakres funkcji \(f\).}
{\([ 2;7] \)}{\([ 1;6] \)}{\([ - 3;3] \)}{\([ - 1;2] \)}{}{}}
\LANGcs{
\Question{
Je dána lineární funkce \(f\colon y = -x + 4\). Určete obor hodnot funkce, která je restrikcí funkce \(f\) na interval \(\langle - 3;2\rangle \).}
{\(\langle 2;7\rangle \)}{\(\langle 1;6\rangle \)}{\(\langle - 3;3\rangle \)}{\(\langle - 1;2\rangle \)}{}{}}
\LANGsk{
\Question{
Daná je lineárna funkcia \(f\colon y = -x + 4\),
na intervale \(x\in \langle - 3;2\rangle \). 
Určte obor hodnôt funkcie \(f\).}
{\(\langle 2;7\rangle \)}{\(\langle 1;6\rangle \)}{\(\langle - 3;3\rangle \)}{\(\langle - 1;2\rangle \)}{}{}}
\LANGes{
\Question{
Sea \(f\) la función lineal \(f(x) = -x + 4\) restringida al intervalo
 \(x\in [ - 3;2] \).
Halla el Rango de \(f\).}
{\([ 2;7] \)}{\([ 1;6] \)}{\([ - 3;3] \)}{\([ - 1;2] \)}{}{}}
\End

\Data\ID{9000005708}
\Updated{2021-01-27 05:01:14}
\Level{A}
\SubArea{Linear functions}
\LANGen{
\Question{
Consider the linear function \(f(x)= -5x + 4\) 
and points \(A = [1;-1]\),
\(B = [-2;-14]\),
\(C = [3;-11]\),
\(D = [-4;24]\).
How many of these points lies on the graph of the function
\(f\)?}
{\(3\)}{\(1\)}{\(2\)}{\(4\)}{}{}}
\LANGpl{
\Question{
Rozważ funkcję liniową  \(f\colon y = -5x + 4\) 
i punkty \(A = [1;-1]\),
\(B = [-2;-14]\),
\(C = [3;-11]\),
\(D = [-4;24]\).
Ile punktów należy do wykresu funkcji 
\(f\)?}
{\(3\)}{\(1\)}{\(2\)}{\(4\)}{}{}}
\LANGcs{
\Question{
Je dána funkce \(f\colon y = -5x + 4\) 
a body \(A = [1;-1]\),
\(B = [-2;-14]\),
\(C = [3;-11]\),
\(D = [-4;24]\).
Kolik z uvedených bodů leží na grafu funkce
\(f\)?}
{\(3\)}{\(1\)}{\(2\)}{\(4\)}{}{}}
\LANGsk{
\Question{
Daná je lineárna funkcia \(f\colon y = -5x + 4\) 
a body \(A = [1;-1]\),
\(B = [-2;-14]\),
\(C = [3;-11]\),
\(D = [-4;24]\).
Koľko z daných bodov leží na grafe funkcie
\(f\)?}
{\(3\)}{\(1\)}{\(2\)}{\(4\)}{}{}}
\LANGes{
\Question{
Considera la función lineal \(f(x)= -5x + 4\) 
y los puntos \(A = (1;-1)\),
\(B = (-2;-14)\),
\(C = (3;-11)\),
\(D = (-4;24)\).
¿Cuántos de estos puntos pertenecen a la gráfica de 
\(f\)?}
{\(3\)}{\(1\)}{\(2\)}{\(4\)}{}{}}
\End

\Data\ID{9000005709}
\Updated{2022-04-23 05:04:37}
\Level{A}
\SubArea{Linear functions}
\LANGen{
\Question{
Consider the linear function \(f(x)= -\frac{4}
{3}x + 4\).
Find the intersection point of the graph of
\(f\) with
\(x\)-axis.}
{\([3;0]\)}{\([0;-6]\)}{\([0;-4]\)}{\([6;0]\)}{}{}}
\LANGpl{
\Question{
Rozważ funkcję liniową \(f\colon y = -\frac{4}
{3}x + 4\).
Znajdź punkt przecięcia wykresu  
\(f\) z osią \(x\).}
{\([3;0]\)}{\([0;-6]\)}{\([0;-4]\)}{\([6;0]\)}{}{}}
\LANGcs{
\Question{
Je dána funkce \(f\colon y = -\frac{4}
{3}x + 4\). Průsečík
grafu této funkce s osou \(x\) 
má souřadnice:}
{\([3;0]\)}{\([0;-6]\)}{\([0;-4]\)}{\([6;0]\)}{}{}}
\LANGsk{
\Question{
Daná je lineárna funkcia \(f\colon y = -\frac{4}
{3}x + 4\). Priesečník grafu funkcie \(f\) s osou \(x\) 
má súradnice:}
{\([3;0]\)}{\([0;-6]\)}{\([0;-4]\)}{\([6;0]\)}{}{}}
\LANGes{
\Question{
Sea \(f(x)= -\frac{4}
{3}x + 4\).
Halla el punto de intersección de 
\(f\) con el eje 
\(x\).}
{\((3;0)\)}{\((0;-6)\)}{\((0;-4)\)}{\((6;0)\)}{}{}}
\End

\Data\ID{9000005710}
\Updated{2022-04-23 05:04:37}
\Level{A}
\SubArea{Linear functions}
\LANGen{
\Question{
Consider the linear function \(f(x) = 4x + 4\).
Find the intersection point of the graph of
\(f\) with
\(y\)-axis.}
{\([0;4]\)}{\([-1;0]\)}{\([-4;0]\)}{\([0;0]\)}{}{}}
\LANGpl{
\Question{
Rozważ funkcję liniową  \(f\colon y = 4x + 4\).
Znajdź punkt przecięcia wykresu funkcji
\(f\) z osią \(y\).}
{\([0;4]\)}{\([-1;0]\)}{\([-4;0]\)}{\([0;0]\)}{}{}}
\LANGcs{
\Question{
Je dána funkce \(f\colon y = 4x + 4\). Průsečík
grafu této funkce s osou \(y\) 
má souřadnice:}
{\([0;4]\)}{\([-1;0]\)}{\([-4;0]\)}{\([0;0]\)}{}{}}
\LANGsk{
\Question{
Daná je lineárna funkcia \(f\colon y = 4x + 4\).
Priesečník grafu funkcie
\(f\) s osou
\(y\) má súradnice:}
{\([0;4]\)}{\([-1;0]\)}{\([-4;0]\)}{\([0;0]\)}{}{}}
\LANGes{
\Question{
Considera la función lineal \(f(x) = 4x + 4\).
Calcula el punto de intersección de 
\(f\) con el eje
\(y\).}
{\((0;4)\)}{\((-1;0)\)}{\((-4;0)\)}{\((0;0)\)}{}{}}
\End

\Data\ID{9000005801}
\Updated{2022-04-23 05:04:37}
\Level{A}
\SubArea{Linear functions}
\LANGen{
\Question{
Given the linear function \(f(x) = -3x + 1\),
evaluate 
\[ 
 f(a) + f(1 - a).
\]}
{\(- 1\)}{\(- 3a\)}{\(- 6a - 3\)}{\(- 2\)}{}{}}
\LANGpl{
\Question{
Dana jest funkcja liniowa \(f\colon y = -3x + 1\).
Oblicz \[ 
 f(a) + f(1 - a).
\]}
{\(- 1\)}{\(- 3a\)}{\(- 6a - 3\)}{\(- 2\)}{}{}}
\LANGcs{
\Question{
Je dána lineární funkce \(f\colon y = -3x + 1\).
Hodnota \(f(a) + f(1 - a)\) 
je rovna:}
{\(- 1\)}{\(- 3a\)}{\(- 6a - 3\)}{\(- 2\)}{}{}}
\LANGsk{
\Question{
Daná je lineárna funkcia \(f\colon y = -3x + 1\).
Určte hodnotu \(f(a) + f(1 - a)\).}
{\(- 1\)}{\(- 3a\)}{\(- 6a - 3\)}{\(- 2\)}{}{}}
\LANGes{
\Question{
Dada la función lineal \(f(x) = -3x + 1\),
evalúa 
\[ 
 f(a) + f(1 - a).
\]}
{\(- 1\)}{\(- 3a\)}{\(- 6a - 3\)}{\(- 2\)}{}{}}
\End

\Data\ID{9000005802}
\Updated{2020-09-18 05:09:31}
\Level{A}
\SubArea{Linear functions}
\LANGen{
\Question{
Given the linear function \(f(x) = -\frac{1}
{4}x + 4\),
evaluate 
\[ 
 f(2a)\cdot f(-2a).
\]}
{\(16 -\frac{a^{2}} 
 {4} \)}{\(0\)}{\(4 - a^{2}\)}{\(- 4 + a^{2}\)}{}{}}
\LANGpl{
\Question{
Dana jest funkcja liniowa \(f\colon y = -\frac{1}
{4}x + 4\).
Oblicz \[ 
 f(2a)\cdot f(-2a).
\]}
{\(16 -\frac{a^{2}} 
 {4} \)}{\(0\)}{\(4 - a^{2}\)}{\(- 4 + a^{2}\)}{}{}}
\LANGcs{
\Question{
Je dána lineární funkce \(f\colon y = -\frac{1}
{4}x + 4\).
Hodnota \(f(2a)\cdot f(-2a)\) 
je rovna:}
{\(16 -\frac{a^{2}} 
 {4} \)}{\(0\)}{\(4 - a^{2}\)}{\(- 4 + a^{2}\)}{}{}}
\LANGsk{
\Question{
Daná je lineárna funkcia \(f\colon y = -\frac{1}
{4}x + 4\).
Určte hodnotu \(f(2a)\cdot f(-2a)\).}
{\(16 -\frac{a^{2}} 
 {4} \)}{\(0\)}{\(4 - a^{2}\)}{\(- 4 + a^{2}\)}{}{}}
\LANGes{
\Question{
Dada la función lineal \(f(x) = -\frac{1}
{4}x + 4\),
evalúa
\[ 
 f(2a)\cdot f(-2a).
\]}
{\(16 -\frac{a^{2}} 
 {4} \)}{\(0\)}{\(4 - a^{2}\)}{\(- 4 + a^{2}\)}{}{}}
\End

\Data\ID{9000005803}
\Updated{2021-01-27 05:01:09}
\Level{A}
\SubArea{Linear functions}
\LANGen{
\Question{
Identify an analytic expression of a linear function
\(f\) which
satisfies \(f(-2) = 5\) 
and \(f(4) = 2\).}
{\(f\colon y = -\frac{1}
{2}x + 4\)}{\(f\colon y = x - 2\)}{\(f\colon y = -x + 6\)}{\(f\colon y = -2x + 1\)}{}{}}
\LANGpl{
\Question{
Wyznacz wzór funkcji liniowej 
\(f\), która spełnia warunki \(f(-2) = 5\) 
i \(f(4) = 2\).}
{\(f\colon y = -\frac{1}
{2}x + 4\)}{\(f\colon y = x - 2\)}{\(f\colon y = -x + 6\)}{\(f\colon y = -2x + 1\)}{}{}}
\LANGcs{
\Question{
Funkční předpis lineární funkce
\(f\), pro kterou
platí \(f(-2) = 5 \wedge f(4) = 2\),
je:}
{\(f\colon y = -\frac{1}
{2}x + 4\)}{\(f\colon y = x - 2\)}{\(f\colon y = -x + 6\)}{\(f\colon y = -2x + 1\)}{}{}}
\LANGsk{
\Question{
Určte predpis lineárnej funkcie
\(f\) pre ktorú 
platí \(f(-2) = 5 \wedge f(4) = 2\).}
{\(f\colon y = -\frac{1}
{2}x + 4\)}{\(f\colon y = x - 2\)}{\(f\colon y = -x + 6\)}{\(f\colon y = -2x + 1\)}{}{}}
\LANGes{
\Question{
Halla la expresión analítica de la función lineal
\(f\) para la cuál se cumple que \(f(-2) = 5\) 
y \(f(4) = 2\).}
{\(f\colon y = -\frac{1}
{2}x + 4\)}{\(f\colon y = x - 2\)}{\(f\colon y = -x + 6\)}{\(f\colon y = -2x + 1\)}{}{}}
\End

\Data\ID{9000007204}
\Updated{2021-01-27 06:01:20}
\Level{A}
\SubArea{Linear functions}
\LANGen{
\Question{
Among the following functions identify an odd function
\(f\) which
satisfies \(f(-2) = 4\).}
{\(f(x) = -2x\)}{\(f(x) = -2x + 1\)}{\(f(x) = x + 2\)}{\(f(x) = 2x + 8\)}{}{}}
\LANGpl{
\Question{
Spośród podanych funkcji wybierz funkcję nieparzystą
\(f\), która spełnia  warunek \(f(-2) = 4\).}
{\(f(x) = -2x\)}{\(f(x) = -2x + 1\)}{\(f(x) = x + 2\)}{\(f(x) = 2x + 8\)}{}{}}
\LANGcs{
\Question{
Která lineární funkce je lichá a platí, že
\(f(-2) = 4\)?}
{\(f(x) = -2x\)}{\(f(x) = -2x + 1\)}{\(f(x)= x + 2\)}{\(f(x) = 2x\)}{}{}}
\LANGsk{
\Question{
Ktorá z nasledujúcich funkcií \(f\) je nepárna
a zároveň platí: \(f(-2) = 4\)?}
{\(f(x) = -2x\)}{\(f(x) = -2x + 1\)}{\(f(x) = x + 2\)}{\(f(x) = 2x + 8\)}{}{}}
\LANGes{
\Question{
Entre las siguientes funciones, encuentra la función
\(f\) que es impar y satisface \(f(-2) = 4\).}
{\(f(x) = -2x\)}{\(f(x) = -2x + 1\)}{\(f(x) = x + 2\)}{\(f(x) = 2x + 8\)}{}{}}
\End

\Data\ID{9000007205}
\Updated{2020-09-21 09:09:33}
\Level{A}
\SubArea{Linear functions}
\LANGen{
\Question{
Given the linear function 
\[ 
 f(x) = -3x + 2
\]
solve the inequality \(f(x)\geq 1\).}
{\(x\leq \frac{1}
{3}\)}{\(x\geq \frac{1}
{3}\)}{\(x\leq \frac{2}
{3}\)}{\(x\geq 2\)}{}{}}
\LANGpl{
\Question{
Dana jest funkcja liniowa
\[ 
 f\colon y = -3x + 2.
\] Rozwiąż podaną nierówność \(f(x)\geq 1\).}
{\(x\leq \frac{1}
{3}\)}{\(x\geq \frac{1}
{3}\)}{\(x\leq \frac{2}
{3}\)}{\(x\geq 2\)}{}{}}
\LANGcs{
\Question{
Je dána funkce \(f\colon y = -3x + 2\).
Pro která \(x\in \mathbb{R}\) 
platí, že \(f(x)\geq 1\)?}
{\(x\leq \frac{1}
{3}\)}{\(x\geq \frac{1}
{3}\)}{\(x\leq \frac{2}
{3}\)}{\(x\geq 2\)}{}{}}
\LANGsk{
\Question{
Daná je lineárna funkcia \(f\colon y = -3x + 2\).
Pre ktoré \(x\in \mathbb{R}\) 
platí, že \(f(x)\geq 1\)?}
{\(x\leq \frac{1}
{3}\)}{\(x\geq \frac{1}
{3}\)}{\(x\leq \frac{2}
{3}\)}{\(x\geq 2\)}{}{}}
\LANGes{
\Question{
Dada la función lineal
\[ 
 f(x) = -3x + 2
\]
resuelve la inecuación \(f(x)\geq 1\).}
{\(x\leq \frac{1}
{3}\)}{\(x\geq \frac{1}
{3}\)}{\(x\leq \frac{2}
{3}\)}{\(x\geq 2\)}{}{}}
\End

\Data\ID{9000007802}
\Updated{2021-04-05 06:04:42}
\Level{A}
\SubArea{Linear functions}
\LANGen{
\Question{
Consider linear functions \(f(x) = ax - 2\) 
and \(g\colon y = -4x + 3\). Find the value of
the real parameter \(a\) which
ensure that the graphs of \(f\) 
and \(g\) 
are two parallel lines.}
{\(- 4\)}{\(4\)}{\(- 2\)}{\(2\)}{}{}}
\LANGpl{
\Question{
Rozważ funkcje liniowe \(f\colon y = ax - 2\) 
i \(g\colon y = -4x + 3\). Znajdź wartość rzeczywistego parametru \(a\), który zapewnia, że wykresy funkcji \(f\) 
i \(g\)  są dwiema liniami równoległymi.}
{\(- 4\)}{\(4\)}{\(- 2\)}{\(2\)}{}{}}
\LANGcs{
\Question{
Jsou dány lineární funkce \(f\colon y = ax - 2,\ g\colon y = -4x + 3\).
Určete koeficient \(a\) 
tak, aby grafy obou funkcí byly rovnoběžné přímky.}
{\(- 4\)}{\(4\)}{\(- 2\)}{\(2\)}{}{}}
\LANGsk{
\Question{
Dané sú lineárne funkcie \(f\colon y = ax - 2,\ g\colon y = -4x + 3\).
Určte koeficient \(a\in \mathbb{R}\)
tak, aby grafy obidvoch funkcií boli rovnobežné priamky.}
{\(- 4\)}{\(4\)}{\(- 2\)}{\(2\)}{}{}}
\LANGes{
\Question{
Considera las funciones lineales \(f(x) = ax - 2\) 
y \(g\colon y = -4x + 3\). Halla el valor del parámetro real
 \(a\) que garantiza que las gráficas de 
 \(f\) 
y \(g\) 
son dos rectas paralelas.}
{\(- 4\)}{\(4\)}{\(- 2\)}{\(2\)}{}{}}
\End

\Data\ID{9000007804}
\Updated{2021-04-05 06:04:24}
\Level{A}
\SubArea{Linear functions}
\LANGen{
\Question{
Three of the points \(A = [2;-4]\),
\(B = [0;-3]\),
\(C = [-2;-1]\),
\(D = [-4;1]\) lie on
the graph of the same linear function. Identify these points. (Colinear points.)}
{\(B\),
\(C\),
\(D\)}{\(A\),
\(B\),
\(C\)}{\(A\),
\(B\),
\(D\)}{\(A\),
\(C\),
\(D\)}{}{}}
\LANGpl{
\Question{
Trzy punkty  \(A = [2;-4]\),
\(B = [0;-3]\),
\(C = [-2;-1]\),
\(D = [-4;1]\) należą do wykresu tej samej funkcji liniowej. Znajdź te punkty.}
{\(B\),
\(C\),
\(D\)}{\(A\),
\(B\),
\(C\)}{\(A\),
\(B\),
\(D\)}{\(A\),
\(C\),
\(D\)}{}{}}
\LANGcs{
\Question{
Jsou dány body \(A = [2;-4]\),
\(B = [0;-3]\),
\(C = [-2;-1]\),
\(D = [-4;1]\).
Tři z nich leží na grafu téže lineární funkce. Které to jsou?}
{\(B\),
\(C\),
\(D\)}{\(A\),
\(B\),
\(C\)}{\(A\),
\(B\),
\(D\)}{\(A\),
\(C\),
\(D\)}{}{}}
\LANGsk{
\Question{
Tri z daných bodov \(A = [2;-4]\),
\(B = [0;-3]\),
\(C = [-2;-1]\),
\(D = [-4;1]\) ležia
na grafe tej istej funkcie. Určte ich.}
{\(B\),
\(C\),
\(D\)}{\(A\),
\(B\),
\(C\)}{\(A\),
\(B\),
\(D\)}{\(A\),
\(C\),
\(D\)}{}{}}
\LANGes{
\Question{
Tres de los puntos \(A = [2;-4]\),
\(B = [0;-3]\),
\(C = [-2;-1]\),
\(D = [-4;1]\) pertencen a la gráfica de la misma función. ¿Cuáles son? (La función es una recta)}
{\(B\),
\(C\),
\(D\)}{\(A\),
\(B\),
\(C\)}{\(A\),
\(B\),
\(D\)}{\(A\),
\(C\),
\(D\)}{}{}}
\End

\Data\ID{9000028105}
\Updated{2021-04-05 07:04:38}
\Level{A}
\SubArea{Linear functions}
\IMG{000281_05-figure0.pdf}
\LANGen{
\Question{
Given graph of the linear function \(g\),
find the solution set of the inequality \(g(x)\leq 0\).}
{\([ 6;\infty )\)}{\(\mathbb{R}\)}{\((-\infty ;2.4] \)}{\((-\infty ;-2.3] \)}{}{}}
\LANGpl{
\Question{
Mając podany wykres funkcji liniowej  \(g\),
znajdź zbiór rozwiązań nierówności \(g(x)\leq 0\).}
{\([ 6;\infty )\)}{\(\mathbb{R}\)}{\((-\infty ;2.4] \)}{\((-\infty ;-2.3] \)}{}{}}
\LANGcs{
\Question{
Vyberte množinu, na které pro lineární funkci
\(g\) platí,
že \(g(x)\leq 0\).}
{\(\langle 6;\infty )\)}{\(\mathbb{R}\)}{\((-\infty ;2{,}4\rangle \)}{\((-\infty ;-2{,}3\rangle \)}{}{}}
\LANGsk{
\Question{
Daný je graf lineárnej funkcie \(g\).
Určte množinu všetkých \(x\in \mathbb{R}\), aby platilo \(g(x)\leq 0\).}
{\(\langle 6;\infty )\)}{\(\mathbb{R}\)}{\((-\infty ;2{,}4\rangle \)}{\((-\infty ;-2{,}3\rangle \)}{}{}}
\LANGes{
\Question{
Dada la gráfica de la función lineal \(g\),
halla el conjunto dónde \(g(x)\leq 0\).}
{\([ 6;\infty )\)}{\(\mathbb{R}\)}{\((-\infty ;2.4] \)}{\((-\infty ;-2.3] \)}{}{}}
\End

\Data\ID{1103171104}
\Updated{2020-12-30 11:12:14}
\Level{B}
\SubArea{Linear functions}
\LANGen{
\Question{
From the given pictures choose the one that shows graphs of three functions corresponding to the formula \( f(x)=kx+2 \), where \( k\in\mathbb{R}^+ \).}
{}{}{}{}{}{}}
\LANGpl{
\Question{
Z podanych rysunków wybierz ten, który przedstawia wykresy trzech funkcji odnoszących się do wzoru  \( f(x)=kx+2 \), gdzie \( k\in\mathbb{R}^+ \).}
{}{}{}{}{}{}}
\LANGcs{
\Question{
Vyberte obrázek, na kterém jsou grafy tří funkcí, které vyhovují předpisu \( f(x)=kx+2 \), kde \( k\in\mathbb{R}^+ \).}
{}{}{}{}{}{}}
\LANGsk{
\Question{
Vyberte obrázok, na ktorom sú grafy troch funkcií, ktoré vyhovujú predpisu \( f(x)=kx+2 \), kde \( k\in\mathbb{R}^+ \).}
{}{}{}{}{}{}}
\LANGes{
\Question{
Elige el dibujo con gráficas de las funciones correspondientes a la fórmula \( f(x)=kx+2 \), dónde \( k\in\mathbb{R}^+ \).}
{}{}{}{}{}{}}

\IMGanswer{1103171104a.pdf}
\IMGanswer{1103171104b.pdf}
\IMGanswer{1103171104c.pdf}
\IMGanswer{1103171104d.pdf}\End

\Data\ID{1103171401}
\Updated{2020-12-30 11:12:39}
\Level{B}
\SubArea{Linear functions}
\IMG{1103171401.pdf}
\LANGen{
\Question{
There are graphs of three linear functions in the picture. Choose the formula which corresponds to all three functions.}
{\( y=m;\ m\in\mathbb{R}^- \)}{\( y=m;\ m\in\mathbb{R}^+ \)}{\( y=mx;\ m\in\mathbb{R}\setminus\{0\} \)}{\( y=x+m;\ m\in\mathbb{R}^- \)}{}{}}
\LANGpl{
\Question{
Przedstawiono wykresy trzech funkcji liniowych na rysunku poniżej. Wybierz wzór, który odnosi się do wszystkich trzech funkcji.}
{\( y=m;\ m\in\mathbb{R}^- \)}{\( y=m;\ m\in\mathbb{R}^+ \)}{\( y=mx;\ m\in\mathbb{R}\setminus\{0\} \)}{\( y=x+m;\ m\in\mathbb{R}^- \)}{}{}}
\LANGcs{
\Question{
Na obrázku jsou grafy tří lineárních funkcí. Z následujících možností vyberte předpis, který odpovídá všem funkcím na obrázku.}
{\( y=m;\ m\in\mathbb{R}^- \)}{\( y=m;\ m\in\mathbb{R}^+ \)}{\( y=mx;\ m\in\mathbb{R}\setminus\{0\} \)}{\( y=x+m;\ m\in\mathbb{R}^- \)}{}{}}
\LANGsk{
\Question{
Dané sú grafy troch lineárnych funkcií na obrázku. Vyberte predpis, ktorý zodpovedá všetkým trom funkciám.}
{\( y=m;\ m\in\mathbb{R}^- \)}{\( y=m;\ m\in\mathbb{R}^+ \)}{\( y=mx;\ m\in\mathbb{R}\setminus\{0\} \)}{\( y=x+m;\ m\in\mathbb{R}^- \)}{}{}}
\LANGes{
\Question{
En el dibujo hay gráficas de tres funciones lineales. Elige la fórmula que corresponde a todas las funciones.}
{\( y=m;\ m\in\mathbb{R}^- \)}{\( y=m;\ m\in\mathbb{R}^+ \)}{\( y=mx;\ m\in\mathbb{R}\setminus\{0\} \)}{\( y=x+m;\ m\in\mathbb{R}^- \)}{}{}}
\End

\Data\ID{1103171402}
\Updated{2020-12-30 11:12:56}
\Level{B}
\SubArea{Linear functions}
\IMG{1103171402.pdf}
\LANGen{
\Question{
There are graphs of three linear functions in the picture. Choose the formula which corresponds to all three functions.}
{\( y=dx-2;\ d\in\mathbb{R}^- \)}{\( y=dx-2;\ d\in\mathbb{R}^+ \)}{\( y=2x+d;\ d\in\mathbb{R}^- \)}{\( y=-2x+d;\ d\in\mathbb{R}^+ \)}{}{}}
\LANGpl{
\Question{
Przedstawiono wykresy trzech funkcji liniowych na rysunku poniżej. Wybierz wzór, który odnosi się do wszystkich trzech funkcji.}
{\( y=dx-2;\ d\in\mathbb{R}^- \)}{\( y=dx-2;\ d\in\mathbb{R}^+ \)}{\( y=2x+d;\ d\in\mathbb{R}^- \)}{\( y=-2x+d;\ d\in\mathbb{R}^+ \)}{}{}}
\LANGcs{
\Question{
Na obrázku jsou grafy 3 lineárních funkcí. Z následujících možností vyberte předpis, který odpovídá všem funkcím na obrázku.}
{\( y=dx-2;\ d\in\mathbb{R}^- \)}{\( y=dx-2;\ d\in\mathbb{R}^+ \)}{\( y=2x+d;\ d\in\mathbb{R}^- \)}{\( y=-2x+d;\ d\in\mathbb{R}^+ \)}{}{}}
\LANGsk{
\Question{
Dané sú grafy troch lineárnych funkcií na obrázku. Vyberte predpis, ktorý zodpovedá všetkým trom funkciám.}
{\( y=dx-2;\ d\in\mathbb{R}^- \)}{\( y=dx-2;\ d\in\mathbb{R}^+ \)}{\( y=2x+d;\ d\in\mathbb{R}^- \)}{\( y=-2x+d;\ d\in\mathbb{R}^+ \)}{}{}}
\LANGes{
\Question{
En el dibujo hay gráficas de tres funciones lineales. Elige la fórmula que describe todas las funciones.}
{\( y=dx-2;\ d\in\mathbb{R}^- \)}{\( y=dx-2;\ d\in\mathbb{R}^+ \)}{\( y=2x+d;\ d\in\mathbb{R}^- \)}{\( y=-2x+d;\ d\in\mathbb{R}^+ \)}{}{}}
\End

\Data\ID{1103171502}
\Updated{2022-07-21 05:07:50}
\Level{B}
\SubArea{Linear functions}
\IMG{1103171502.pdf}
\LANGen{
\Question{
Current-Voltage characteristic of a conductor is a relationship, usually represented as a graph, between the current \( I \) through the conductor and the voltage \( U \) across the conductor. The picture shows two characteristics: of a metal conductor (red colour) and of an electrolytic conductor (black colour). At what voltages need to be conductors connected so that the current is the same in both of them?}
{\( 10\,\mathrm{V} \)}{\( 16\,\mathrm{V} \)}{\( 8\,\mathrm{V} \)}{\( 9.5\,\mathrm{V} \)}{}{}}
\LANGpl{
\Question{
Charakterystyka prądowo-napięciowa przewodnika jest relacją, zwykle reprezentowaną jako wykres, pomiędzy prądem \(I\) przez przewodnik i napięciem \(U\) na przewodzie. Rysunek przedstawia dwie charakterystyki przewodnika: metalowego (czerwony) i elektrolitycznego (czarny). Przy jakim napięciu muszą być podłączone przewody, aby prąd był taki sam w obu?}
{\( 10\,\mathrm{V} \)}{\( 16\,\mathrm{V} \)}{\( 8\,\mathrm{V} \)}{\( 9{,}5\,\mathrm{V} \)}{}{}}
\LANGcs{
\Question{
Voltampérová charakteristika vodiče je graficky prezentována závislost mezi elektrickým proudem \( I \), který jím protéká a napětím \( U \) mezi jeho konci. Na obrázku jsou tyto charakteristiky pro kovový vodič (červená barva) a elektrolytický vodič (černá barva). Na jaké napětí musíme připojit oba vodiče, aby jimi protékal stejný proud?}
{\( 10\,\mathrm{V} \)}{\( 16\,\mathrm{V} \)}{\( 8\,\mathrm{V} \)}{\( 9{,}5\,\mathrm{V} \)}{}{}}
\LANGsk{
\Question{
Voltampérová charakteristika vodiča je grafické znázornenie závislosti elektrického prúdu \( I \), prechádzajúceho súčiastkou od napätia medzi jej vývodmi \( U \). Na obrázku sú znázornené tieto charakteristiky pre kovový vodič (červená farba) a elektrolytický vodič (čierna farba). Na aké napätie musíme pripojiť obidva vodiče, aby cez nich pretekal rovnaký prúd?}
{\( 10\,\mathrm{V} \)}{\( 16\,\mathrm{V} \)}{\( 8\,\mathrm{V} \)}{\( 9{,}5\,\mathrm{V} \)}{}{}}
\LANGes{
\Question{
La proporción de Corriente-Voltaje de un conductor es la relación entre la corriente \( I \)  y el voltaje \( U \) en el conductor y normalmente se representa por una gráfica. En el dibujo hay dos características: de un conductor de metal (de color rojo) y de un conductor electrolítico (de color negro). ¿A qué voltajes necesitan estar conectados los conductores para que ambos tengan la misma corriente?}
{\( 10\,\mathrm{V} \)}{\( 16\,\mathrm{V} \)}{\( 8\,\mathrm{V} \)}{\( 9.5\,\mathrm{V} \)}{}{}}
\End

\Data\ID{9000005805}
\Updated{2020-09-14 12:09:09}
\Level{B}
\SubArea{Linear functions}
\LANGen{
\Question{
Consider the linear function \(f(x) = x\).
Identify a linear function \(g\) 
such that the graphs of \(f\) 
and \(g\) are symmetric
about the \(x\)-axis.}
{\(g(x) = -x\)}{\(g(x)= x\)}{\(g(x) = x + 1\)}{\(g(x) = x - 1\)}{}{}}
\LANGpl{
\Question{
Rozważ funkcję liniową \(f\colon y = x\).
Wybierz funkcję liniową \(g\) 
taką, by wykresy funkcji \(f\) 
i \(g\) były symetryczne względem osi
 \(x\) układu współrzędnych.}
{\(g\colon y = -x\)}{\(g\colon y = x\)}{\(g\colon y = x + 1\)}{\(g\colon y = x - 1\)}{}{}}
\LANGcs{
\Question{
Je dána lineární funkce \(f\colon y = x\).
Předpis funkce \(g\), jejíž
graf je s grafem funkce \(f\) 
souměrný podle osy \(x\),
je:}
{\(g\colon y = -x\)}{\(g\colon y = x\)}{\(g\colon y = x + 1\)}{\(g\colon y = x - 1\)}{}{}}
\LANGsk{
\Question{
Daná je lineárna funkcia \(f\colon y = x\).
Určte lineárnu funkciu \(g\), ak
grafy funkcií \(f\) a \(g\)
sú súmerné podľa osi \(x\).}
{\(g\colon y = -x\)}{\(g\colon y = x\)}{\(g\colon y = x + 1\)}{\(g\colon y = x - 1\)}{}{}}
\LANGes{
\Question{
Considera la función lineal \(f(x) = x\).
Encuentra una función lineal \(g\) 
tal que las gráficas de \(f\) 
y  \(g\) sean simétricas respecto al eje \(x\).}
{\(g(x) = -x\)}{\(g(x)= x\)}{\(g(x) = x + 1\)}{\(g(x) = x - 1\)}{}{}}
\End

\Data\ID{9000005806}
\Updated{2021-01-27 05:01:00}
\Level{B}
\SubArea{Linear functions}
\LANGen{
\Question{
Consider the linear function \(f(x) = -x + 2\).
Identify a linear function \(g\) 
such that the graphs of \(f\) 
and \(g\) are symmetric
about the line \(y = x\).}
{\(g(x) = -x + 2\)}{\(g(x)= -x - 2\)}{\(g(x)= x - 2\)}{\(g(x) = x + 2\)}{}{}}
\LANGpl{
\Question{
Rozważ funkcję liniową \(f\colon y = -x + 2\).
Wybierz funkcję liniową \(g\) 
taką, by wykresy funkcji \(f\) 
i \(g\) były symetryczne względem
prostej \(y = x\).}
{\(g\colon y = -x + 2\)}{\(g\colon y = -x - 2\)}{\(g\colon y = x - 2\)}{\(g\colon y = x + 2\)}{}{}}
\LANGcs{
\Question{
Je dána lineární funkce \(f\colon y = -x + 2\).
Předpis funkce \(g\), jejíž
graf je s grafem funkce \(f\) 
souměrný podle osy I. a III. kvadrantu, je:}
{\(g\colon y = -x + 2\)}{\(g\colon y = -x - 2\)}{\(g\colon y = x - 2\)}{\(g\colon y = x + 2\)}{}{}}
\LANGsk{
\Question{
Daná je lineárna funkcia \(f\colon y = -x + 2\).
Určte lineárnu funkciu \(g\), 
ak grafy funkcií \(f\) a \(g\) sú súmerné
podľa priamky \(y = x\).}
{\(g\colon y = -x + 2\)}{\(g\colon y = -x - 2\)}{\(g\colon y = x - 2\)}{\(g\colon y = x + 2\)}{}{}}
\LANGes{
\Question{
Considera la función lineal \(f(x) = -x + 2\).
Encuentra una función lineal \(g\) 
tal que las gráficas de \(f\) 
y \(g\) sean simétricas respecto a la recta \(y = x\).}
{\(g(x) = -x + 2\)}{\(g(x)= -x - 2\)}{\(g(x)= x - 2\)}{\(g(x) = x + 2\)}{}{}}
\End

\Data\ID{9000007808}
\Updated{2021-04-05 07:04:10}
\Level{B}
\SubArea{Linear functions}
\LANGen{
\Question{
Given a function \(f(x) = \frac{x} 
{3} + 1\), find
the function \(g\) such that
the graph of \(g\) is symmetric
with the graph of \(f\) 
about the \(y\)-axis.}
{\(g\colon y = -\frac{x}
{3} + 1\)}{\(g\colon y = 3x + 1\)}{\(g\colon y = -3x + 1\)}{\(g\colon y = -\frac{x}
{3} - 1\)}{Such a function does not exist.}{}}
\LANGpl{
\Question{
Dana jest funkcja \(f\colon y = \frac{x} 
{3} + 1\). Wyznacz taką funkcję  \(g\), by jej wykres
był symetryczny do wykresu
funkcji \(f\) 
względem osi \(y\) układu współrzędnych.}
{\(g\colon y = -\frac{x}
{3} + 1\)}{\(g\colon y = 3x + 1\)}{\(g\colon y = -3x + 1\)}{\(g\colon y = -\frac{x}
{3} - 1\)}{Taka funkcja nie istnieje.}{}}
\LANGcs{
\Question{
Je dána funkce \(f\colon y = \frac{x} 
{3} + 1\).
Určete předpis funkce \(g\),
jejíž graf je souměrný s grafem funkce
\(f\) podle
osy \(y\).}
{\(g\colon y = -\frac{x}
{3} + 1\)}{\(g\colon y = 3x + 1\)}{\(g\colon y = -3x + 1\)}{\(g\colon y = -\frac{x}
{3} - 1\)}{Taková funkce neexistuje.}{}}
\LANGsk{
\Question{
Daná je lineárna funkcia \(f\colon y = \frac{x} 
{3} + 1\). Určte lineárnu funkciu \(g\), ak grafy funkcií \(f\) a \(g\) 
sú súmerné podľa osi \(y\).}
{\(g\colon y = -\frac{x}
{3} + 1\)}{\(g\colon y = 3x + 1\)}{\(g\colon y = -3x + 1\)}{\(g\colon y = -\frac{x}
{3} - 1\)}{Taká funkcia neexistuje.}{}}
\LANGes{
\Question{
Dada la función \(f(x) = \frac{x} 
{3} + 1\), halla
la función \(g\) de manera que la gráfica de \(g\) y la gráfica de \(f\) 
sean simétricas respecto al eje \(y\).}
{\(g\colon y = -\frac{x}
{3} + 1\)}{\(g\colon y = 3x + 1\)}{\(g\colon y = -3x + 1\)}{\(g\colon y = -\frac{x}
{3} - 1\)}{No existe dicha función.}{}}
\End

\Data\ID{9000028101}
\Updated{2020-09-11 12:09:57}
\Level{B}
\SubArea{Linear functions}
\IMG{000281_01-figure0.pdf}
\LANGen{
\Question{
Given graphs of the linear functions \(f\) 
and \(g\), find the solution
set of the inequality \(f(x) > g(x)\).}
{\((1;\infty )\)}{\(\mathbb{R}\)}{\((0;1)\)}{\((-\infty ;1)\)}{}{}}
\LANGpl{
\Question{
Dane są wykresy funkcji liniowych \(f\) 
i \(g\). Znajdź zbiór rozwiązań nierówności \(f(x) > g(x)\).}
{\((1;\infty )\)}{\(\mathbb{R}\)}{\((0;1)\)}{\((-\infty ;1)\)}{}{}}
\LANGcs{
\Question{
Na obrázku jsou vyznačeny grafy lineárních funkcí \(f\) a \(g\). Určete, pro která \(x\) platí, že \(f(x) > g(x)\).}
{\((1;\infty )\)}{\(\mathbb{R}\)}{\((0;1)\)}{\((-\infty ;1)\)}{}{}}
\LANGsk{
\Question{
Dané sú grafy lineárnych funkcií
\(f\) a
\(g\). Určte množinu všetkých \(x\in \mathbb{R}\),
aby platilo \(f(x) > g(x)\).}
{\((1;\infty )\)}{\(\mathbb{R}\)}{\((0;1)\)}{\((-\infty ;1)\)}{}{}}
\LANGes{
\Question{
Dadas las gráficas de las funciones lineales \(f\) y \(g\), halla el conjunto dónde se cumple \(f(x) > g(x)\).}
{\((1;\infty )\)}{\(\mathbb{R}\)}{\((0;1)\)}{\((-\infty ;1)\)}{}{}}
\End

\Data\ID{9000028102}
\Updated{2020-09-11 12:09:17}
\Level{B}
\SubArea{Linear functions}
\IMG{000281_02-figure0.pdf}
\LANGen{
\Question{
Given graph of the linear function \(f\),
find the solution set of the inequality \(f(x) > 0\).}
{\((-4;\infty )\)}{\(\emptyset \)}{\((-4;2)\)}{\((-2;\infty )\)}{}{}}
\LANGpl{
\Question{
Dany jest wykres funkcji liniowej \(f\).
Znajdź zbiór rozwiązań nierówności \(f(x) > 0\).}
{\((-4;\infty )\)}{\(\emptyset \)}{\((-4;2)\)}{\((-2;\infty )\)}{}{}}
\LANGcs{
\Question{
Najděte množinu všech \( x \in \mathbb{R} \), pro která platí \(f(x) > 0\).}
{\((-4;\infty )\)}{\(\emptyset \)}{\((-4;2)\)}{\((-2;\infty )\)}{}{}}
\LANGsk{
\Question{
Daný je graf lineárnej funkcie \(f\).
Určte množinu všetkých \(x\in \mathbb{R}\), aby platilo \(f(x) > 0\).}
{\((-4;\infty )\)}{\(\emptyset \)}{\((-4;2)\)}{\((-2;\infty )\)}{}{}}
\LANGes{
\Question{
Dada la gráfica de la función lineal \(f\),
halla el conjunto dónde se cumple \(f(x) > 0\).}
{\((-4;\infty )\)}{\(\emptyset \)}{\((-4;2)\)}{\((-2;\infty )\)}{}{}}
\End

\Data\ID{9000028103}
\Updated{2021-04-05 07:04:58}
\Level{B}
\SubArea{Linear functions}
\IMG{000281_03-figure0.pdf}
\LANGen{
\Question{
Given graphs of the linear functions \(f\) 
and \(g\), find the solution
set of the inequality \(f(x) > g(x)\).}
{\((-\infty ;-2)\)}{\(\emptyset \)}{\((-4;2)\)}{\((-2;\infty )\)}{}{}}
\LANGpl{
\Question{
Dane są wykresy funkcji liniowych \(f\) 
i \(g\). Znajdź zbiór rozwiązań nierówności  \(f(x) > g(x)\).}
{\((-\infty ;-2)\)}{\(\emptyset \)}{\((-4;2)\)}{\((-2;\infty )\)}{}{}}
\LANGcs{
\Question{
Vyberte množinu, na které pro lineární funkce
\(f\) a
\(g\) platí,
že \(f(x) > g(x)\).}
{\((-\infty ;-2)\)}{\(\emptyset \)}{\((-4;2)\)}{\((-2;\infty )\)}{}{}}
\LANGsk{
\Question{
Dané sú grafy lineárnych funkcií \(f\) a \(g\).
Určte množinu všetkých \(x\in \mathbb{R}\), aby platilo \(f(x) > g(x)\).}
{\((-\infty ;-2)\)}{\(\emptyset \)}{\((-4;2)\)}{\((-2;\infty )\)}{}{}}
\LANGes{
\Question{
Dadas las gráficas de las funciones lineales \(f\) y \(g\), halla el conjunto dónde  \(f(x) > g(x)\).}
{\((-\infty ;-2)\)}{\(\emptyset \)}{\((-4;2)\)}{\((-2;\infty )\)}{}{}}
\End

\Data\ID{9000028104}
\Updated{2021-04-05 07:04:17}
\Level{B}
\SubArea{Linear functions}
\IMG{000281_04-figure0.pdf}
\LANGen{
\Question{
Given graphs of the linear functions \(f\) 
and \(g\), find the solution
set of the inequality \(f(x)\leq g(x)\).}
{\((-\infty ;2.4] \)}{\(\emptyset \)}{\((-\infty ;-2.3] \)}{\([ 6;\infty )\)}{}{}}
\LANGpl{
\Question{
Dane są wykresy funkcji liniowych \(f\) 
i \(g\). Znajdź zbiór rozwiązań nierówności \(f(x)\leq g(x)\).}
{\((-\infty ;2.4] \)}{\(\emptyset \)}{\((-\infty ;-2.3] \)}{\([ 6;\infty )\)}{}{}}
\LANGcs{
\Question{
Vyberte množinu, na které pro lineární funkce
\(f\) a
\(g\) platí,
že \(f(x)\leq g(x)\).}
{\((-\infty ;2{,}4\rangle \)}{\(\emptyset \)}{\((-\infty ;-2{,}3\rangle \)}{\(\langle 6;\infty )\)}{}{}}
\LANGsk{
\Question{
Dané sú grafy lineárnych funkcií
\(f\) a
\(g\). Určte množinu všetkých \(x\in \mathbb{R}\),
aby platilo \(f(x)\leq g(x)\).}
{\((-\infty ;2{,}4\rangle \)}{\(\emptyset \)}{\((-\infty ;-2{,}3\rangle \)}{\(\langle 6;\infty )\)}{}{}}
\LANGes{
\Question{
Dadas las gráficas de las funciones lineales \(f\) y \(g\), halla el conjunto dnde \(f(x)\leq g(x)\).}
{\((-\infty ;2.4] \)}{\(\emptyset \)}{\((-\infty ;-2.3] \)}{\([ 6;\infty )\)}{}{}}
\End

\Data\ID{9000028106}
\Updated{2021-04-05 07:04:33}
\Level{B}
\SubArea{Linear functions}
\IMG{000281_06-figure0.pdf}
\LANGen{
\Question{
Given graphs of the linear functions \(f\) 
and \(g\), find the solution
set of the inequality \(f(x)\leq g(x)\).}
{\(\emptyset \)}{\((-\infty ;0] \)}{\(\mathbb{R}\)}{\([ 0;\infty )\)}{}{}}
\LANGpl{
\Question{
Mając podane wykresy funkcji \(f\) 
i \(g\), znajdź zbiór rozwiązań nierówności \(f(x)\leq g(x)\).}
{\(\emptyset \)}{\((-\infty ;0] \)}{\(\mathbb{R}\)}{\([ 0;\infty )\)}{}{}}
\LANGcs{
\Question{
Vyberte množinu, na které pro lineární funkce
\(f\) a
\(g\) platí,
že \(f(x)\leq g(x)\).}
{\(\emptyset \)}{\((-\infty ;0\rangle \)}{\(\mathbb{R}\)}{\(\langle 0;\infty )\)}{}{}}
\LANGsk{
\Question{
Dané sú grafy lineárnych funkcií
\(f\) a
\(g\). Určte množinu všetkých \(x\in \mathbb{R}\),
aby platilo \(f(x)\leq g(x)\).}
{\(\emptyset \)}{\((-\infty ;0\rangle \)}{\(\mathbb{R}\)}{\(\langle 0;\infty )\)}{}{}}
\LANGes{
\Question{
Dadas las gráficas de las funciones lineales \(f\) y 
\(g\), halla el conjunto dónde  \(f(x)\leq g(x)\).}
{\(\emptyset \)}{\((-\infty ;0] \)}{\(\mathbb{R}\)}{\([ 0;\infty )\)}{}{}}
\End

\Data\ID{9000028107}
\Updated{2021-04-05 07:04:52}
\Level{B}
\SubArea{Linear functions}
\IMG{000281_07-figure0.pdf}
\LANGen{
\Question{
Given graphs of the linear functions \(f\) 
and \(g\), find the solution
set of the inequality \(f(x)\leq g(x)\).}
{\(\mathbb{R}\)}{\(\emptyset \)}{\((-\infty ;0] \)}{\([ 0;\infty )\)}{}{}}
\LANGpl{
\Question{
Dane są wykresy funkcji liniowych \(f\) 
i \(g\). Znajdź zbiór rozwiązań nierówności 
 \(f(x)\leq g(x)\).}
{\(\mathbb{R}\)}{\(\emptyset \)}{\((-\infty ;0] \)}{\([ 0;\infty )\)}{}{}}
\LANGcs{
\Question{
Vyberte množinu, na které pro lineární funkce
\(f\) a
\(g\) platí,
že \(f(x)\leq g(x)\).}
{\(\mathbb{R}\)}{\(\emptyset \)}{\((-\infty ;0\rangle \)}{\(\langle 0;\infty )\)}{}{}}
\LANGsk{
\Question{
Dané sú grafy lineárnych funkcií
\(f\) a
\(g\). Určte množinu všetkých \(x\in \mathbb{R}\),
aby platilo \(f(x)\leq g(x)\).}
{\(\mathbb{R}\)}{\(\emptyset \)}{\((-\infty ;0\rangle \)}{\(\langle 0;\infty )\)}{}{}}
\LANGes{
\Question{
Dadas las gráficas de las funciones lineales \(f\) y 
\(g\), halla el conjunto dónde \(f(x)\leq g(x)\).}
{\(\mathbb{R}\)}{\(\emptyset \)}{\((-\infty ;0] \)}{\([ 0;\infty )\)}{}{}}
\End

\Data\ID{9000028108}
\Updated{2021-04-05 07:04:13}
\Level{B}
\SubArea{Linear functions}
\IMG{000281_08-figure0.pdf}
\LANGen{
\Question{
Given graph of the linear function \(f\),
find the solution set of the inequality \(f(x) < 0\).}
{\(\emptyset \)}{\((-\infty ;0] \)}{\(\mathbb{R}\)}{\([ 0;\infty )\)}{}{}}
\LANGpl{
\Question{
Dany jest wykres funkcji liniowych \(f\). 
Znajdź zbiór rozwiązań nierówności \(f(x) < 0\).}
{\(\emptyset \)}{\((-\infty ;0] \)}{\(\mathbb{R}\)}{\([ 0;\infty )\)}{}{}}
\LANGcs{
\Question{
Vyberte množinu, na které pro lineární funkci
\(f\) platí,
že \(f(x) < 0\).}
{\(\emptyset \)}{\((-\infty ;0\rangle \)}{\(\mathbb{R}\)}{\(\langle 0;\infty )\)}{}{}}
\LANGsk{
\Question{
Daný je graf lineárnej funkcie \(f\).
Určte množinu všetkých \(x\in \mathbb{R}\), aby platilo \(f(x) < 0\).}
{\(\emptyset \)}{\((-\infty ;0\rangle \)}{\(\mathbb{R}\)}{\(\langle 0;\infty )\)}{}{}}
\LANGes{
\Question{
Dada la gráfica de la función lineal \(f\),
halla el conjunto dónde \(f(x) < 0\).}
{\(\emptyset \)}{\((-\infty ;0] \)}{\(\mathbb{R}\)}{\([ 0;\infty )\)}{}{}}
\End

\Data\ID{9000028109}
\Updated{2021-04-05 07:04:32}
\Level{B}
\SubArea{Linear functions}
\IMG{000281_09-figure0.pdf}
\LANGen{
\Question{
Given graphs of the linear functions \(f\),
\(g\) and
\(h\), find the solution set
of the inequality \(f(x)\leq g(x) < h(x)\).}
{\((3.73;\infty )\)}{\([ - 1.04;1)\)}{\((1;\infty )\)}{\([ 1;3.73)\)}{}{}}
\LANGpl{
\Question{
Dane są wykresy funkcji liniowych \(f\),
\(g\) i \(h\). Znajdź zbiór rozwiązań
nierówności  \(f(x)\leq g(x) < h(x)\).}
{\((3.73;\infty )\)}{\([ - 1.04;1)\)}{\((1;\infty )\)}{\([ 1;3.73)\)}{}{}}
\LANGcs{
\Question{
Vyberte množinu, na které pro lineární funkce
\(f\),
\(g\) a
\(h\) platí,
že \(f(x)\leq g(x) < h(x)\).}
{\((3{,}73;\infty )\)}{\(\langle - 1{,}04;1)\)}{\((1;\infty )\)}{\(\langle 1;3{,}73)\)}{}{}}
\LANGsk{
\Question{
Dané sú grafy lineárnych funkcií
\(f\), \(g\) a \(h\). Určte množinu všetkých \(x\in \mathbb{R}\),
aby platilo \(f(x)\leq g(x) < h(x)\).}
{\((3{,}73;\infty )\)}{\(\langle - 1{,}04;1)\)}{\((1;\infty )\)}{\(\langle 1;3{,}73)\)}{}{}}
\LANGes{
\Question{
Dadas las gráficas de las funciones lineales \(f\),
\(g\) y
\(h\), halla el conjunto dónde \(f(x)\leq g(x) < h(x)\).}
{\((3.73;\infty )\)}{\([ - 1.04;1)\)}{\((1;\infty )\)}{\([ 1;3.73)\)}{}{}}
\End

\Data\ID{9000028110}
\Updated{2021-04-05 07:04:00}
\Level{B}
\SubArea{Linear functions}
\IMG{000281_10-figure0.pdf}
\LANGen{
\Question{
Given graphs of linear functions \(f\),
\(g\) and
\(h\), find the solution set
of the inequality \(f(x)\leq g(x) < h(x)\).}
{\([ 4;7)\)}{\((-\infty ;4] \)}{\([ 1;7)\)}{\([ 7;\infty )\)}{}{}}
\LANGpl{
\Question{
Dane są wykresy funkcji liniowych \(f\),
\(g\) i \(h\). Znajdź zbiór rozwiązań nierówności \(f(x)\leq g(x) < h(x)\).}
{\([ 4;7)\)}{\((-\infty ;4] \)}{\([ 1;7)\)}{\([ 7;\infty )\)}{}{}}
\LANGcs{
\Question{
Vyberte množinu, na které pro lineární funkce
\(f\),
\(g\) a
\(h\) platí,
že \(f(x)\leq g(x) < h(x)\).}
{\(\langle 4;7)\)}{\((-\infty ;4\rangle \)}{\(\langle 1;7)\)}{\(\langle 7;\infty )\)}{}{}}
\LANGsk{
\Question{
Dané sú grafy lineárnych funkcií
\(f\), \(g\) a \(h\). Určte množinu všetkých \(x\in \mathbb{R}\),
aby platilo \(f(x)\leq g(x) < h(x)\).}
{\(\langle 4;7)\)}{\((-\infty ;4\rangle \)}{\(\langle 1;7)\)}{\(\langle 7;\infty )\)}{}{}}
\LANGes{
\Question{
Dadas las gráficas de las funciones lineales \(f\),
\(g\) y
\(h\), halla el conjunto dónde \(f(x)\leq g(x) < h(x)\).}
{\([ 4;7)\)}{\((-\infty ;4] \)}{\([ 1;7)\)}{\([ 7;\infty )\)}{}{}}
\End

\Data\ID{1003160901}
\Updated{2020-12-30 11:12:43}
\Level{C}
\SubArea{Linear functions}
\LANGen{
\Question{
Suppose \( f \) is a linear function. If the value of the independent variable \( x \) increases by \( 6 \), the function value increases by \( 18 \). Choose the correct form of the function \( f \), such that the given property is preserved.}
{\( f(x)=3x+1 \)}{\( f(x)=-3x \)}{\( f(x)=\frac13x+18 \)}{\( f(x)=\frac13x \)}{}{}}
\LANGpl{
\Question{
Załóżmy, że \( f \) jest funkcją liniową. Jeżeli wartość niezależnej zmiennej \( x \) wzrośnie o  \( 6 \), wartość funkcji wzrośnie o \( 18 \). Wybierz poprawny wzór funkcji \( f \), taki, w którym podane właściwości zostaną zachowane.}
{\( f(x)=3x+1 \)}{\( f(x)=-3x \)}{\( f(x)=\frac13x+18 \)}{\( f(x)=\frac13x \)}{}{}}
\LANGcs{
\Question{
Je dána lineární funkce \( f \). Jestliže se nezávisle proměnná \( x \) zvětší o \( 6 \), zvětší se funkční hodnota o \( 18 \). Z následujících možností vyberte předpis, který uvedenou vlastnost zohledňuje.}
{\( f(x)=3x+1 \)}{\( f(x)=-3x \)}{\( f(x)=\frac13x+18 \)}{\( f(x)=\frac13x \)}{}{}}
\LANGsk{
\Question{
Daná je lineárna funkcia \( f \). Ak sa hodnota nezávislej premennej \( x \) zväčší o \( 6 \), zväčší sa funkčná hodnota o \( 18 \). Vyberte správny predpis funkcie \( f \), ktorý zachováva danú vlastnosť.}
{\( f(x)=3x+1 \)}{\( f(x)=-3x \)}{\( f(x)=\frac13x+18 \)}{\( f(x)=\frac13x \)}{}{}}
\LANGes{
\Question{
Sea \( f \) una función lineal. Si el valor de la variable independiente \( x \) aumenta en \( 6 \), el valor de la función aumenta en \( 18 \). Elije la función \( f \) correcta para que se cumpla la condición,}
{\( f(x)=3x+1 \)}{\( f(x)=-3x \)}{\( f(x)=\frac13x+18 \)}{\( f(x)=\frac13x \)}{}{}}
\End

\Data\ID{1003160902}
\Updated{2020-12-30 11:12:27}
\Level{C}
\SubArea{Linear functions}
\LANGen{
\Question{
Suppose \( f \) is a linear function. If the value of the independent variable \( x \) increases by \( 4 \), the function value decreases by \( 12 \).  Choose the correct form of the function \( f \), such that the given property is preserved.}
{\( f(x)=-3x \)}{\( f(x)=3x \)}{\( f(x)=3x-12 \)}{\( f(x)=-\frac13x \)}{}{}}
\LANGpl{
\Question{
Załóżmy, że\( f \) jest funkcją liniową. Jeśli wartość niezależnej zmiennej \( x \) wzrośnie o  \( 4 \), wartość funkcji wzrośnie o  \( 12 \).  Wybierz odpowiedni wzór funkcji \( f \), taki, w którym dane właściwości zostaną zachowane.}
{\( f(x)=-3x \)}{\( f(x)=3x \)}{\( f(x)=3x-12 \)}{\( f(x)=-\frac13x \)}{}{}}
\LANGcs{
\Question{
Je dána lineární funkce \( f \). Jestliže se nezávisle proměnná \( x \) zvětší o \( 4 \), zmenší se funkční hodnota o \( 12 \). Z následujících možností vyberte předpis, který uvedenou vlastnost zohledňuje.}
{\( f(x)=-3x \)}{\( f(x)=3x \)}{\( f(x)=3x-12 \)}{\( f(x)=-\frac13x \)}{}{}}
\LANGsk{
\Question{
Daná je lineárna funkcia \( f \). Ak sa hodnota nezávislej premennej \( x \) zväčší o \( 4 \), zmenší sa funkčná hodnota o \( 12 \). Vyberte správny predpis funkcie \( f \), ktorý zachováva danú vlastnosť.}
{\( f(x)=-3x \)}{\( f(x)=3x \)}{\( f(x)=3x-12 \)}{\( f(x)=-\frac13x \)}{}{}}
\LANGes{
\Question{
Sea \( f \) una función lineal. Si el valor de la variable independiente \( x \) aumenta en \( 4 \), el valor de la función decrece en \( 12 \). Elije la función \( f \) para que se cumpla la propiedad.}
{\( f(x)=-3x \)}{\( f(x)=3x \)}{\( f(x)=3x-12 \)}{\( f(x)=-\frac13x \)}{}{}}
\End

\Data\ID{1003160903}
\Updated{2020-12-30 11:12:53}
\Level{C}
\SubArea{Linear functions}
\LANGen{
\Question{
The graph of a linear function \( f \) passes through the point \( \left[4\sqrt3;2\right] \) and it makes an angle of \( 30^{\circ} \) with the positive direction of \( x \)-axis measured anticlockwise. Choose the correct form of the function \( f \), such that the given property is preserved.}
{\( f(x)=\frac{\sqrt3}3x-2 \)}{\( f(x)=\sqrt3x-10 \)}{\( f(x)=\frac{\sqrt3}3x+2 \)}{\( f(x)=\sqrt3x+10 \)}{}{}}
\LANGpl{
\Question{
Wykres funkcji liniowej  \( f \) przechodzi przez punkt \( \left[4\sqrt3;2\right] \) i tworzy kąt  \( 30^{\circ} \) z dodatnim kierunkiem osi współrzędnych  \( x \)  mierząc odwrotnie do wskazówek zegara. Wybierz taki wzór funkcji \( f \), by dane właściwości zostały zachowane.}
{\( f(x)=\frac{\sqrt3}3x-2 \)}{\( f(x)=\sqrt3x-10 \)}{\( f(x)=\frac{\sqrt3}3x+2 \)}{\( f(x)=\sqrt3x+10 \)}{}{}}
\LANGcs{
\Question{
Graf lineární funkce \( f \) prochází bodem \( \left[4\sqrt3;2\right] \) a se směrem kladné poloosy \( x \) svírá úhel \( 30^{\circ} \). Z následujících možností vyberte funkci, jejíž graf má dané vlastnosti.}
{\( f(x)=\frac{\sqrt3}3x-2 \)}{\( f(x)=\sqrt3x-10 \)}{\( f(x)=\frac{\sqrt3}3x+2 \)}{\( f(x)=\sqrt3x+10 \)}{}{}}
\LANGsk{
\Question{
Graf lineárnej funkcie \( f \) prechádza bodom \( \left[4\sqrt3;2\right] \) a s kladným smerom osi \( x \) zviera uhol \( 30^{\circ} \). Z daných predpisov vyberte funkciu,  ktorá spĺňa dané vlastnosti.}
{\( f(x)=\frac{\sqrt3}3x-2 \)}{\( f(x)=\sqrt3x-10 \)}{\( f(x)=\frac{\sqrt3}3x+2 \)}{\( f(x)=\sqrt3x+10 \)}{}{}}
\LANGes{
\Question{
La gráfica de la función lineal \( f \) pasa por el punto \( \left[4\sqrt3;2\right] \) y entre la gráfica y la parte positiva de la eje \( x \) hay un ángulo de\( 30^{\circ} \). Elije la función \( f \)correcta para que se cumpla la condición.}
{\( f(x)=\frac{\sqrt3}3x-2 \)}{\( f(x)=\sqrt3x-10 \)}{\( f(x)=\frac{\sqrt3}3x+2 \)}{\( f(x)=\sqrt3x+10 \)}{}{}}
\End

\Data\ID{1003171301}
\Updated{2020-12-30 11:12:08}
\Level{C}
\SubArea{Linear functions}
\LANGen{
\Question{
The freezing point and the boiling point of water (both under the normal atmospheric pressure) are the base of the most commonly used temperature scale around Europe. It is Celsius temperature scale in Celsius degrees (\( ^{\circ}\mathrm{C} \)).  Fahrenheit temperature scale in Fahrenheit degrees (\( ^{\circ}\mathrm{F} \)) is of general common use in English speaking countries especially in the USA. The basic temperature points in mentioned scales have the values:
\[ \begin{array}{l}
\text{Water freezing point } \dots\ 0\,^{\circ}\mathrm{C} / 32\,^{\circ}\mathrm{F} \\
\text{Water boiling point } \dots\ 100\,^{\circ}\mathrm{C} / 212\,^{\circ}\mathrm{F} 
\end{array} \]
From the following equations choose the one that can be used to convert a temperature from its Celsius representation to the Fahrenheit value, provided you know that the relationship between the scales is linear. (In equations, \( F \) is a numerical value of a temperature in Fahrenheits and \( C \) is a numerical value of a temperature in Celsius.)}
{\( F=\frac95 C+32 \)}{\( F=\frac59C+32 \)}{\( F=\frac59 C-\frac{160}9 \)}{\( F=32C+100 \)}{}{}}
\LANGpl{
\Question{
Temperatura zamarzania i temperatura wrzenia wody (obydwie pod normalnym ciśnieniem atmosferycznym) są podstawą najczęściej stosowanej skali temperatury w Europie. Jest to skala temperatury Celsjusza wyrażona  w stopniach Celsjusza (\( ^{\circ}\mathrm{C} \)). Skala temperatury Fahrenheita wyrażona w stopniach Fahrenheita jest powszechnie stosowana w krajach anglojęzycznych, szczególnie w USA. Podstawowe punkty temperatury w wymienionych skalach mają następujące wartości:
\[ \begin{array}{l}
\text{punkt zamarzania wody } \dots\ 0\,^{\circ}\mathrm{C} / 32\,^{\circ}\mathrm{F} \\
\text{punkt wrzenia wody } \dots\ 100\,^{\circ}\mathrm{C} / 212\,^{\circ}\mathrm{F} 
\end{array} \]
Z podanych równań wybierz równanie, którego można użyć, żeby zamienić temperaturę wyrażoną w stopniach Celsjusza na tę wyrażona za pomocą skali Fahrenheita, zakładając, że relacja pomiędzy tymi skalami jest liniowa. (Iw równaniach, \( F \) jest wartością numeryczną temperatury Fahrenheita, a \( C \) jest wartością numeryczną w stopniach Celsjusza.)}
{\( F=\frac95 C+32 \)}{\( F=\frac59C+32 \)}{\( F=\frac59 C-\frac{160}9 \)}{\( F=32C+100 \)}{}{}}
\LANGcs{
\Question{
Teploty tuhnutí a varu vody (obojí za normálního tlaku) jsou základními body v Evropě nejpoužívanější  teploměrné stupnice — Celsiovy (jednotka \( ^{\circ}\mathrm{C} \)). V USA je ale nejrozšířenější stupnicí Fahrenheitova  (jednotka \( ^{\circ}\mathrm{F} \)). Uváděné teploty jsou v obou stupnicích vyjádřeny těmito hodnotami:
\[ \begin{array}{l}
\text{Teplota tuhnutí vody } \dots\ 0\,^{\circ}\mathrm{C} / 32\,^{\circ}\mathrm{F} \\
\text{Teplota varu vody } \dots\ 100\,^{\circ}\mathrm{C} / 212\,^{\circ}\mathrm{F} 
\end{array} \]
Z následujících možností vyberte rovnici, která umožní výpočet teploty Fahrenheitovy ze známé teploty Celsiovy, když víme, že mezi těmito stupnicemi je lineární funkční vztah. (\( F \) je teplota ve Fahrenheitově stupnici a \( C \) je teplota v Celsiově stupnici.)}
{\( F=\frac95 C+32 \)}{\( F=\frac59C+32 \)}{\( F=\frac59 C-\frac{160}9 \)}{\( F=32C+100 \)}{}{}}
\LANGsk{
\Question{
Teploty tuhnutia a varu vody (obidve za normálneho tlaku) sú základnými bodmi v Europe najpoužívanejšej teplotnej stupnice — Celziovej (jednotka \( ^{\circ}\mathrm{C} \)). V USA je však rozšírenejšia  Fahrenheitova stupnica  (jednotka \( ^{\circ}\mathrm{F} \)). Uvádzané teploty sú pri oboch stupniciach vyjadrené hodnotami:
\[ \begin{array}{l}
\text{Teplota tuhnutia vody } \dots\ 0\,^{\circ}\mathrm{C} / 32\,^{\circ}\mathrm{F} \\
\text{Teplota varu vody } \dots\ 100\,^{\circ}\mathrm{C} / 212\,^{\circ}\mathrm{F} 
\end{array} \]
Z následujúcich možností vyberte rovnicu, ktorá umožní výpočet Fahrenheitovej teploty zo známej Celziovej teploty, ak vieme, že medzi týmito stupnicami je lineárny vzťah. (\( F \) je číselná hodnota teploty vo Fahrenheitovej stupnici a \( C \) je číselná hodnota teploty v Celziovej stupnici.)}
{\( F=\frac95 C+32 \)}{\( F=\frac59C+32 \)}{\( F=\frac59 C-\frac{160}9 \)}{\( F=32C+100 \)}{}{}}
\LANGes{
\Question{
Los puntos de congelación y ebullición del agua (en presión atmosférica normal) son la base de la escala de temperatura más usada en Europa. Se llama escala de temperatura de Celsius y se mide en grados Celsius (\( ^{\circ}\mathrm{C} \)).  La escala de temperatura Fahrenheit se mide en grados Fahrenheit (\( ^{\circ}\mathrm{F} \)) y es la más usada en países angloparlantes, especialmente en E.E.U.U. Los puntos básicos tienen estos valores:
\[ \begin{array}{l}
\text{Punto de congelación de agua } \dots\ 0\,^{\circ}\mathrm{C} / 32\,^{\circ}\mathrm{F} \\
\text{Punto de ebullición de agua } \dots\ 100\,^{\circ}\mathrm{C} / 212\,^{\circ}\mathrm{F} 
\end{array} \]
De las siguientes ecuaciones elige la que describe la conversión de grados Celsius a grados Fahrenheit si sabemos que la relación es lineal. (En las ecuaciones \( F \) es el valor numérico de la temperatura en la escala Fahrenheit y \( C \) es el valor numérico de la temperatura en la escala Celsius.)}
{\( F=\frac95 C+32 \)}{\( F=\frac59C+32 \)}{\( F=\frac59 C-\frac{160}9 \)}{\( F=32C+100 \)}{}{}}
\End

\Data\ID{1003171601}
\Updated{2020-12-30 11:12:03}
\Level{C}
\SubArea{Linear functions}
\LANGen{
\Question{
Consider the function \( f \) given by \( f(x)=\frac12x+\frac32 \) and consider the line \( p \) that is parallel to the \( x \)-axis and intersects \( y \)-axis at the point \( \left[0;\frac12\right] \). Find the function \( g \) such that the graph of \( g \) is symmetric with the graph of \( f \) about the line \( p \).}
{\( g(x)=-\frac12x-\frac12 \)}{\( g(x)=2x-\frac12 \)}{\( g(x)=-\frac12x-\frac32 \)}{\( g(x)=\frac12x-\frac32 \)}{}{}}
\LANGpl{
\Question{
Rozważ funkcję \( f \) wyrażoną przez \( f(x)=\frac12x+\frac32 \) i rozważ linię \( p \) która jest równoległa do osi współrzędnych  \( x \) i przecina oś współrzędnych \( y \) w punkcie  \( \left[0;\frac12\right] \). Wyznacz funkcję  \( g \) taką, dla której wykres funkcji \( g \) jest symetryczny z wykresem funkcji \( f \) względem linii \( p \).}
{\( g(x)=-\frac12x-\frac12 \)}{\( g(x)=2x-\frac12 \)}{\( g(x)=-\frac12x-\frac32 \)}{\( g(x)=\frac12x-\frac32 \)}{}{}}
\LANGcs{
\Question{
Je dána funkce \( f \) s předpisem \( f(x)=\frac12x+\frac32 \) a přímka \( p \), která je rovnoběžná s osou \( x \) a protíná osu \( y \) v bodě \( \left[0;\frac12\right] \). Vyberte předpis funkce \( g \), jejíž graf bude souměrný s grafem funkce \( f \) podle přímky \( p \).}
{\( g(x)=-\frac12x-\frac12 \)}{\( g(x)=2x-\frac12 \)}{\( g(x)=-\frac12x-\frac32 \)}{\( g(x)=\frac12x-\frac32 \)}{}{}}
\LANGsk{
\Question{
Daná je funkcia \( f \) s predpisom \( f(x)=\frac12x+\frac32 \) a priamka \( p \) rovnobežná s osou \( x \), ktorá pretína os \( y \) v bode \( \left[0;\frac12\right] \). Určte predpis funkcie \( g \), ktorej graf je osovo súmerný  s grafom funkcie \( f \) podľa priamky \( p \).}
{\( g(x)=-\frac12x-\frac12 \)}{\( g(x)=2x-\frac12 \)}{\( g(x)=-\frac12x-\frac32 \)}{\( g(x)=\frac12x-\frac32 \)}{}{}}
\LANGes{
\Question{
Considera la función \( f \) dada por \( f(x)=\frac12x+\frac32 \) y la recta \( p \) que es paralela al eje \( x \) e intersecta al eje \( y \) en el punto \( \left[0;\frac12\right] \). Halla la función \( g \) tal que la gráfica de \( g \) es simétrica a la gráfica de \( f \) respecto a la recta \( p \).}
{\( g(x)=-\frac12x-\frac12 \)}{\( g(x)=2x-\frac12 \)}{\( g(x)=-\frac12x-\frac32 \)}{\( g(x)=\frac12x-\frac32 \)}{}{}}
\End

\Data\ID{1103171201}
\Updated{2020-09-10 01:09:19}
\Level{C}
\SubArea{Linear functions}
\LANGen{
\Question{
Suppose a train comes to stop uniformly. The instantaneous velocity \( v \) of the train in the time \( t \) of slowing down is modelled by \( v=45-25t \). Find the graph of the function \( v \).}
{}{}{}{}{}{}}
\LANGpl{
\Question{
Załóżmy, że pociąg zatrzymuje się równomiernie. Chwilowa prędkość \( v \) pociągu w czasie zwalniania \( t \) przedstawiono za pomocą \( v=45-25t \). Wyznacz wykres funkcji \( v \).}
{}{}{}{}{}{}}
\LANGcs{
\Question{
Okamžitá rychlost vlaku je při rovnoměrném brzdění popsána rovnicí \( v=45-25t \). Z následujících možností vyberte graf odpovídající závislosti okamžité rychlosti vlaku \( v \) na době jeho brzdění \( t \).}
{}{}{}{}{}{}}
\LANGsk{
\Question{
Okamžitú rýchlosť vlaku pri rovnomernom brzdení popisuje rovnica \( v=45-25t \). Z následujúcich možností vyberte graf, ktorý zodpovedá závislosti okamžitej rýchlosti vlaku \( v \) na čase jeho brzdenia \( t \).}
{}{}{}{}{}{}}
\LANGes{
\Question{
Supongamos que un tren para uniformemente. La velocidad instantánea \( v \) del tren en el tiempo \( t \) decreciendo se representa por \( v=45-25t \). Halla la gráfica de la función \( v \).}
{}{}{}{}{}{}}

\IMGanswer{1103171201a.pdf}
\IMGanswer{1103171201b.pdf}
\IMGanswer{1103171201c.pdf}
\IMGanswer{1103171201d.pdf}\End

\Data\ID{1103171501}
\Updated{2020-09-11 10:09:39}
\Level{C}
\SubArea{Linear functions}
\IMG{1103171501.pdf}
\LANGen{
\Question{
Ohm's law states that the current \( I \) through a conductor is directly proportional to the voltage \( U \) between the endpoints of the conductor. This relationship is described by the equation \( I=\frac UR \), where \(  R \) is the resistance of the conductor. Current-voltage characteristics of the conductors \( A \) and \( B \) are in the picture. Which of the conductors has greater resistance?}
{\( A \)}{\( B \)}{Both conductors have the same resistance.}{It is not possible to answer the question based on the graph.}{}{}}
\LANGpl{
\Question{
Prawo Ohm'a głosi, że prąd \( I \) przechodzący przez przewodnik jest wprost proporcjonalny do napięcia \( U \) pomiędzy końcami przewodnika. Ta zależność została zapisana w postaci równania \( I=\frac UR \), gdzie \(  R \) jest oporem przewodnika. Charakterystyka prądowo-napięciowa   \( A \) i \( B \) została przedstawiona na rysunku.  Który z przewodników ma większy opór?}
{\( A \)}{\( B \)}{Obydwa maja taki sam opór.}{Nie można udzielić odpowiedzi na podstawie wykresu.}{}{}}
\LANGcs{
\Question{
Ohmův zákon vyjadřuje vztah přímé úměrnosti mezi proudem \( I \), který prochází vodičem a napětím mezi jeho konci \( U \). Tento vztah je vyjádřen rovnicí \( I=\frac UR \), kde \(  R \) je elektrický odpor vodiče. Na obrázku jsou grafy průběhu proudu v závislosti na napětí ve dvou různých vodičích. Který z vodičů má větší elektrický odpor \( R \)?}
{\( A \)}{\( B \)}{Oba vodiče mají stejný odpor.}{Na základě daného grafu není možné otázku zodpovědět.}{}{}}
\LANGsk{
\Question{
Ohmov zákon vyjadruje vzťah priamej úmernosti medzi prúdom \( I \), ktorý prechádza vodičom a napätím na koncoch vodiča \( U \). Tento vzťah je vyjadrený rovnicou \( I=\frac UR \), kde \(  R \) je elektrický odpor vodiča. Na obrázku sú dva grafy priebehu prúdu v závislosti na napätí v dvoch rôznych vodičoch. Ktorý z vodičov má väčší elektrický odpor \( R \)?}
{\( A \)}{\( B \)}{Obidva vodiče majú rovnaký odpor.}{Na základe daného grafu nie je možné na otázku odpovedať.}{}{}}
\LANGes{
\Question{
La ley de Ohm dice que la corriente \( I \) que pasa por un conductor es directamente proporcional al voltaje \( U \)  entre los puntos finales del conductor. Esta relación se puede describir por la ecuación \( I=\frac UR \), donde \(  R \) es la resistencia del conductor. Las características de la corriente y el voltaje de los conductores \( A \) y \( B \)  están descritas en el dibujo.  ¿Cuál de los conductores tiene mayor resistencia?}
{\( A \)}{\( B \)}{Los conductores tienen la misma resistencia.}{No es posible responder la pregunta con el dibujo.}{}{}}
\End

\Data\ID{1103171504}
\Updated{2021-01-27 05:01:32}
\Level{C}
\SubArea{Linear functions}
\IMG{1103171504.pdf}
\LANGen{
\Question{
The picture shows velocity-time graphs of movements of cars \( A \), \( B \), \( C \) and \( D \).  Which of the cars speeds up with constant acceleration of \( 0.8\,\frac{\mathrm{m}}{\mathrm{s}^2} \)? 
\[ \]
Hint: An acceleration \( a \) is the rate of change of velocity \( \Delta v \) of an object with respect to time \( \Delta t \), i.e. \( a=\frac{\Delta v}{\Delta t} \).}
{\( A \)}{\( B \)}{\( C \)}{\( D \)}{}{}}
\LANGpl{
\Question{
Rysunek przedstawia wykresy prędkości i czasu ruchu samochodów  \( A \), \( B \), \( C \) i \( D \).  Który samochód przyspiesza ze stałym przyspieszeniem \( 0{,}8\,\frac{\mathrm{m}}{\mathrm{s}^2} \)? 
\[ \]
Wskazówka: Przyspieszenie \( a \) jest wskaźnikiem zmiany prędkości \( \Delta v \) przedmiotu w odniesieniu do czasu \( \Delta t \), t.j. \( a=\frac{\Delta v}{\Delta t} \).}
{\( A \)}{\( B \)}{\( C \)}{\( D \)}{}{}}
\LANGcs{
\Question{
Na obrázku je grafická závislost rychlosti na čase pro pohyb aut \( A \), \( B \), \( C \) a \( D \). Které auto se rozjíždí se stálým zrychlením \( 0{,}8\,\frac{\mathrm{m}}{\mathrm{s}^2} \)? 
\[ \]
Nápověda: Zrychlení tělesa \( a \) je definováno jako podíl změny jeho rychlosti \( \Delta v \) a doby \( \Delta t \), ve které tato změna nastala, tj. \( a=\frac{\Delta v}{\Delta t} \).}
{\( A \)}{\( B \)}{\( C \)}{\( D \)}{}{}}
\LANGsk{
\Question{
Na obrázku je grafická závislosť rýchlosti na čase pre pohyb aut \( A \), \( B \), \( C \) a \( D \).  Ktoré auto sa rozbieha so stálym zrýchlením \( 0{,}8\,\frac{\mathrm{m}}{\mathrm{s}^2} \)? 
\[ \]
Nápoveda: Zrýchlenie telesa \( a \) je definované ako podiel zmeny rýchlosti \( \Delta v \) a času \( \Delta t \), t.j. \( a=\frac{\Delta v}{\Delta t} \).}
{\( A \)}{\( B \)}{\( C \)}{\( D \)}{}{}}
\LANGes{
\Question{
Las gráficas representan la relación entre velocidad y tiempo de los coches \( A \), \( B \), \( C \) y \( D \). ¿Cuál de los coches acelera con una acceleración constante de \( 0.8\,\frac{\mathrm{m}}{\mathrm{s}^2} \)? 
Pista: Una acceleración \( a \)  es la proporción de cambio de velocidad \( \Delta v \)  de un objeto respecto al tiempo \( \Delta t \), es decir, \( a=\frac{\Delta v}{\Delta t} \).}
{\( A \)}{\( B \)}{\( C \)}{\( D \)}{}{}}
\End

\Data\ID{2000003109}
\Updated{2022-01-17 06:01:07}
\Level{C}
\SubArea{Linear functions}
\LANGen{
\Question{
In the morning at \(7\,\mathrm{a.m.}\)  we measured \(3^\circ\mathrm{C}\), at \(10\,\mathrm{a.m.}\) we measured \(12^\circ \mathrm{C}\). How many degrees was at \(9\,\mathrm{a.m.}\), if we assume that the temperature rose linearly?}
{\(9^\circ\mathrm{C}\)}{\(10^\circ\mathrm{C}\)}{\(8^\circ\mathrm{C}\)}{\(6^\circ\mathrm{C}\)}{}{}}
\LANGpl{
\Question{
Rano o godzinie 7.00 temperatura wynosiła \(3^\circ\mathrm{C}\), o 10.00 wynosiła \(12^\circ \mathrm{C}\). Ile stopni było o godzinie 9.00 jeśli przyjmiemy, że temperatura rosła liniowo?}
{\(9^\circ\mathrm{C}\)}{\(10^\circ\mathrm{C}\)}{\(8^\circ\mathrm{C}\)}{\(6^\circ\mathrm{C}\)}{}{}}
\LANGcs{
\Question{
V průběhu dopoledne jsme v  \(7\,\mathrm{hodin}\)  naměřili \(3^\circ\mathrm{C}\) a v \(10\,\mathrm{hodin}\) jsme naměřili \(12^\circ \mathrm{C}\). Kolik stupňů bylo v \(9\,\mathrm{hodin}\) za předpokladu, že teplota rostla lineárně?}
{\(9^\circ\mathrm{C}\)}{\(10^\circ\mathrm{C}\)}{\(8^\circ\mathrm{C}\)}{\(6^\circ\mathrm{C}\)}{}{}}
\LANGsk{
\Question{
Ráno o \(7.\,\) hodine sme namerali \(3^\circ\mathrm{C}\), a o \(10.\,\) hodine už \(12^\circ \mathrm{C}\). Koľko stupňov bolo o \(9.\,\) hodine, ak predpokladáme, že teplota rástla lineárne?}
{\(9^\circ\mathrm{C}\)}{\(10^\circ\mathrm{C}\)}{\(8^\circ\mathrm{C}\)}{\(6^\circ\mathrm{C}\)}{}{}}
\LANGes{
\Question{
Durante una mañana hemos medido la temperatura. A las  \(7\,\mathrm{horas}\) hemos medido \(3^\circ\mathrm{C}\) y a las \(10\,\mathrm{horas}\) hemos medido\(12^\circ \mathrm{C}\). ¿Cuántos grados medimos a las \(9\,\mathrm{horas}\), si suponemos que la temperatura creció linealmente?}
{\(9^\circ\mathrm{C}\)}{\(10^\circ\mathrm{C}\)}{\(8^\circ\mathrm{C}\)}{\(6^\circ\mathrm{C}\)}{}{}}
\End

\Data\ID{2000003110}
\Updated{2022-01-17 06:01:54}
\Level{C}
\SubArea{Linear functions}
\LANGen{
\Question{
Consider functions \(f(x)=-2x+3\) and \(g(x)=3x-2\). What is the value of \(f(g(f(-2)))\)?}
{\(-35\)}{\(13\)}{\(-1\)}{\(-8\)}{}{}}
\LANGpl{
\Question{
Dana są funkcje \(f(x)=-2x+3\) i \(g(x)=3x-2\). Oblicz wartość \(f(g(f(-2)))\).}
{\(-35\)}{\(13\)}{\(-1\)}{\(-8\)}{}{}}
\LANGcs{
\Question{
Jsou dány funkce \(f(x)=-2x+3\) a \(g(x)=3x-2\). Jaká je hodnota \(f(g(f(-2)))\)?}
{\(-35\)}{\(13\)}{\(-1\)}{\(-8\)}{}{}}
\LANGsk{
\Question{
Dané sú funkcie  \(f(x)=-2x+3\) a \(g(x)=3x-2\). Vypočítajte hodnotu \(f(g(f(-2)))\).}
{\(-35\)}{\(13\)}{\(-1\)}{\(-8\)}{}{}}
\LANGes{
\Question{
Dadas las funciones \(f(x)=-2x+3\) y \(g(x)=3x-2\). ¿Cuál es el valor \(f(g(f(-2)))\)?}
{\(-35\)}{\(13\)}{\(-1\)}{\(-8\)}{}{}}
\End

\Data\ID{9000005804}
\Updated{2020-09-14 12:09:03}
\Level{C}
\SubArea{Linear functions}
\LANGen{
\Question{
Given the linear function \(f(x) = 5x - 3\),
solve the following equation. 
\[ 
 f(x) = 5a + 2
\]}
{\(x = a + 1\)}{\(x = a\)}{\(x = a + 5\)}{\(x = a - 1\)}{}{}}
\LANGpl{
\Question{
Dana jest funkcja liniowa \(f\colon y = 5x - 3\). 
Rozwiąż podane równanie.
\[ 
 f(x) = 5a + 2
\]}
{\(x = a + 1\)}{\(x = a\)}{\(x = a + 5\)}{\(x = a - 1\)}{}{}}
\LANGcs{
\Question{
Je dána lineární funkce \(f\colon y = 5x - 3\).
Určete, pro které \(x\in \mathbb{R}\) 
platí, že \(f(x) = 5a + 2\).}
{\(x = a + 1\)}{\(x = a\)}{\(x = a + 5\)}{\(x = a - 1\)}{}{}}
\LANGsk{
\Question{
Daná je lineárna funkcia \(f\colon y = 5x - 3\).
Určte, pre ktoré \(x\in \mathbb{R}\) 
platí, že \(f(x) = 5a + 2\).}
{\(x = a + 1\)}{\(x = a\)}{\(x = a + 5\)}{\(x = a - 1\)}{}{}}
\LANGes{
\Question{
Dada la función linear \(f(x) = 5x - 3\),
Resuelve la ecuación. 
\[ 
 f(x) = 5a + 2
\]}
{\(x = a + 1\)}{\(x = a\)}{\(x = a + 5\)}{\(x = a - 1\)}{}{}}
\End

\Data\ID{9000005807}
\Updated{2021-01-27 05:01:49}
\Level{C}
\SubArea{Linear functions}
\LANGen{
\Question{
Consider the function \(f(x) = -x + 4\) 
and a triangle which has one side on the graph of
\(f\) and
two other sides on the axes. Find the area of this triangle.}
{\(8\)}{\(4\)}{\(6\)}{\(10\)}{}{}}
\LANGpl{
\Question{
Rozważ funkcję \(f\colon y = -x + 4\) 
i trójkąt, którego jeden z boków należy do wykresu  tej funkcji
\(f\), a dwa pozostałe znajdują się na osiach układu współrzędnych. Oblicz pole powierzchni tego trójkąta.}
{\(8\)}{\(4\)}{\(6\)}{\(10\)}{}{}}
\LANGcs{
\Question{
Najděte obsah trojúhelníku, jehož jedna strana leží na grafu funkce
\(f\colon y = -x + 4\) a zbylé dvě
strany na osách \(x\) 
a \(y\).}
{\(8\)}{\(4\)}{\(6\)}{\(10\)}{}{}}
\LANGsk{
\Question{
Určte obsah trojuholníka, ktorého jedna strana leží na grafe funkcie \(f\colon y = -x + 4\) 
a zvyšné dve strany ležia na osiach \(x\) a \(y\).}
{\(8\)}{\(4\)}{\(6\)}{\(10\)}{}{}}
\LANGes{
\Question{
Considera la función \(f(x) = -x + 4\) 
y el triángulo formado por la gráfica de la función
\(f\) y los ejes. Halla el área del triángulo.}
{\(8\)}{\(4\)}{\(6\)}{\(10\)}{}{}}
\End

\Data\ID{9000005808}
\Updated{2020-09-18 05:09:59}
\Level{C}
\SubArea{Linear functions}
\LANGen{
\Question{
Consider the linear functions \(f(x) = x\),
\(g(x) = -x\) and
\(h(x)= 3\).
Consider also a triangle which is formed by the graphs of these three functions. Find
the area of this triangle.}
{\(9\)}{\(3\)}{\(5\)}{\(7\)}{}{}}
\LANGpl{
\Question{
Rozważ funkcje liniowe  \(f\colon y = x\),
\(g\colon y = -x\) i
\(h\colon y = 3\).
Wykresy tych trzech funkcji tworzą trójkąt. Oblicz pole powierzchni tego trójkąta.}
{\(9\)}{\(3\)}{\(5\)}{\(7\)}{}{}}
\LANGcs{
\Question{
Jsou dány funkce \(f\colon y = x\),
\(g\colon y = -x\) a
\(h\colon y = 3\).
Najděte obsah trojúhelníku, jehož strany leží na grafech těchto funkcí.}
{\(9\)}{\(3\)}{\(5\)}{\(7\)}{}{}}
\LANGsk{
\Question{
Dané sú funkcie \(f\colon y = x\),
\(g\colon y = -x\) a
\(h\colon y = 3\).
Určte obsah trojuholníka, ktorého strany ležia na grafoch daných funkcií.}
{\(9\)}{\(3\)}{\(5\)}{\(7\)}{}{}}
\LANGes{
\Question{
Considera las funciones lineales \(f(x) = x\),
\(g(x) = -x\) y
\(h(x)= 3\).
Considera también el triángulo formado por las gráficas de estas tres funciones. Halla el área del triángulo.}
{\(9\)}{\(3\)}{\(5\)}{\(7\)}{}{}}
\End

\Data\ID{9000005809}
\Updated{2021-01-27 05:01:26}
\Level{C}
\SubArea{Linear functions}
\LANGen{
\Question{
Consider the functions \(f(x) = x - 1\) 
and \(g(x) y = -x + a\). Find the value
of the real parameter \(a\in \mathbb{R}\) 
which ensure that the functions have a common value at
\(x = 3\), i.e.
\(f(3) = g(3)\).}
{\(a = 5\)}{\(a = -1\)}{\(a = 1\)}{\(a = 2\)}{}{}}
\LANGpl{
\Question{
Rozważ funkcje \(f\colon y = x - 1\) 
i \(g\colon y = -x + a\). Znajdź wartość rzeczywistego parametru 
\(a\in \mathbb{R}\), który gwarantuje, że funkcje mają wspólną wartość przy 
\(x = 3\), tj.
\(f(3) = g(3)\).}
{\(a = 5\)}{\(a = -1\)}{\(a = 1\)}{\(a = 2\)}{}{}}
\LANGcs{
\Question{
Jsou dány funkce \(f\colon y = x - 1\) 
a \(g\colon y = -x + a\).
Určete \(a\in \mathbb{R}\) 
tak, aby obě funkce měly stejnou funkční hodnotu pro
\(x = 3\).}
{\(a = 5\)}{\(a = -1\)}{\(a = 1\)}{\(a = 2\)}{}{}}
\LANGsk{
\Question{
Dané sú funkcie \(f\colon y = x - 1\) 
a \(g\colon y = -x + a\). Určte hodnotu parametra \(a\in \mathbb{R}\),
aby pre \(x = 3\), platilo \(f(3) = g(3)\).}
{\(a = 5\)}{\(a = -1\)}{\(a = 1\)}{\(a = 2\)}{}{}}
\LANGes{
\Question{
Considera las funciones \(f(x) = x - 1\) 
y \(g(x)  = -x + a\). Halla el valor del parámetro real \(a\in \mathbb{R}\) 
para que se cumpla que \(f(3) = g(3)\).}
{\(a = 5\)}{\(a = -1\)}{\(a = 1\)}{\(a = 2\)}{}{}}
\End

\Data\ID{9000005810}
\Updated{2021-01-27 06:01:11}
\Level{C}
\SubArea{Linear functions}
\LANGen{
\Question{
Consider the functions \(f(x) = x + 1\) 
and \(g(x) = ax + 7\). For what values of the real parameter  \(a\in \mathbb{R}\),  
do the functions have the same function value equal to \(3\).}
{\(a = -2\)}{\(a = -7\)}{\(a = -1\)}{\(a = 7\)}{}{}}
\LANGpl{
\Question{
Rozważ funkcje \(f\colon y = x + 1\) 
i \(g\colon y = ax + 7\). Znajdź wartość rzeczywistego parametru
\(a\in \mathbb{R}\), który gwarantuje, że funkcje mają wspólną wartość równą \(3\).}
{\(a = -2\)}{\(a = -7\)}{\(a = -1\)}{\(a = 7\)}{}{}}
\LANGcs{
\Question{
Jsou dány funkce \(f\colon y = x + 1\) 
a \(g\colon y = ax + 7\). Pro
která \(a\in \mathbb{R}\) 
mají obě funkce stejnou funkční hodnotu rovnu
\(3\)?}
{\(a = -2\)}{\(a = -7\)}{\(a = -1\)}{\(a = 7\)}{}{}}
\LANGsk{
\Question{
Dané sú funkcie \(f\colon y = x + 1\) 
a \(g\colon y = ax + 7\). Pre
ktoré hodnoty parametra \(a\in \mathbb{R}\) 
majú obidve funkcie funkčnú hodnotu rovnú
\(3\)?}
{\(a = -2\)}{\(a = -7\)}{\(a = -1\)}{\(a = 7\)}{}{}}
\LANGes{
\Question{
Considera las funciones \(f(x) = x + 1\) 
y \(g(x) = ax + 7\). ¿Para qué valores del parámetro real  \(a\in \mathbb{R}\),  
tienen las funciones el mismo valor y es igual a  \(3\)?}
{\(a = -2\)}{\(a = -7\)}{\(a = -1\)}{\(a = 7\)}{}{}}
\End

\Data\ID{9000007201}
\Updated{2021-01-27 06:01:09}
\Level{C}
\SubArea{Linear functions}
\LANGen{
\Question{
Consider the function 
\[ 
 f(x) = [x + 2]
\]
defined on the domain \(\mathop{\mathrm{Dom}}(f) = (1;2)\).
Find the parameters \(a\) 
and \(b\) 
in the linear function 
\[ 
 g(x) = ax + b
\]
which ensure that the functions \(f\) 
and \(g\) are identical
on the domain of \(f\).
\[ \]
 Hint: The function \(y = [x]\) 
is a floor function: the largest integer less than or equal to
\(x\). For positive
\(x\) it is also called the
integer part of \(x\).}
{\(a = 0\),
\(b = 3\); 
\(\mathop{\mathrm{Dom}}(g) = (1;2)\)}{\(a = 3\),
\(b = 0\);
\(\mathop{\mathrm{Dom}}(g) = (1;2)\)}{\(a = 0\),
\(b = 4\);
\(\mathop{\mathrm{Dom}}(g) = (1;2)\)}{\(a = -3\),
\(b = 0\);
\(\mathop{\mathrm{Dom}}(g) = (1;2)\)}{}{}}
\LANGpl{
\Question{
Rozważ funkcję
\[ 
 f\colon y = [x + 2]
\]
określoną na dziedzinie \(\mathop{\mathrm{Dom}}(f) = (1;2)\).
Znajdź wartości \(a\) 
i \(b\) 
funkcji liniowej
\[ 
 g\colon y = ax + b,
\] które zapewnią, że funkcje \(f\) 
i  \(g\) są tożsamościowe
na dziedzinie funkcji \(f\).
\[ \]
 Wskazówka: Funkcja  \(y = [x]\) 
jest podłogą: największa liczba całkowita jest mniejsza lub równa 
\(x\). Dla dodatniego
\(x\) nazwana jest również  częścią ułamkową liczby całkowitej  \(x\).}
{\(a = 0\),
\(b = 3\);
\(\mathop{\mathrm{Dom}}(g) = (1;2)\)}{\(a = 3\),
\(b = 0\);
\(\mathop{\mathrm{Dom}}(g) = (1;2)\)}{\(a = 0\),
\(b = 4\);
\(\mathop{\mathrm{Dom}}(g) = (1;2)\)}{\(a = -3\),
\(b = 0\);
\(\mathop{\mathrm{Dom}}(g) = (1;2)\)}{}{}}
\LANGcs{
\Question{
Je dána funkce \(f\colon y = [x + 2]\) a
platí \(D(f) = (1;2)\). Co musí
platit pro koeficienty \(a\),
\(b\) a definiční obor
lineární funkce \(g\colon y = ax + b\), aby se
rovnala zadané funkci \(f\)?
\[ \]
 Nápověda: Funkce \(y = [x]\) je celá
část čísla \(x\). Každému
reálnému číslu \(x\) 
přiřadí největší celé číslo, které je menší, nebo rovno
\(x\).}
{\(a = 0\ \wedge \ b = 3 \ ; \ D(g) = (1;2)\)}{\(a = 3\ \wedge \ b = 0\ ; \ D(g) = (1;2)\)}{\(a = 0\ \wedge \ b = 4\ ;\ D(g) = (1;2)\)}{\(a = -3\ \wedge \ b = 0\ ; \ D(g) = (1;2)\)}{}{}}
\LANGsk{
\Question{
Daná je funkcia \(f\colon y = [x + 2]\) a
platí \(D(f) = (1;2)\). Čo musí
platiť pre koeficienty \(a\),
\(b\) a pre definičný obor
lineárnej funkcie \(g\colon y = ax + b\), aby sa
rovnala zadanej funkcii \(f\)? 
\[ \]
Poznámka: Funkcia \(y = [x]\) je celá
časť čísla \(x\). Každému
reálnemu číslu \(x\) 
priradí najväčšie celé číslo, ktoré je menšie alebo rovné
\(x\).}
{\(a = 0\ \wedge \ b = 3\ ; \ D(g) = (1;2)\)}{\(a = 3\ \wedge \ b = 0\ ; \ D(g) = (1;2)\)}{\(a = 0\ \wedge \ b = 4\ ; \ D(g) = (1;2)\)}{\(a = -3\ \wedge \ b = 0\ ; \ D(g) = (1;2)\)}{}{}}
\LANGes{
\Question{
Considera la función
\[ 
 f(x) = [x + 2]
\]
definida en el Dominio \(\mathop{\mathrm{Dom}}(f) = (1;2)\).
Halla los parámetros \(a\) 
y \(b\) 
en la función lineal
\[ 
 g(x) = ax + b
\]
que garantizan que las funciones \(f\) 
y \(g\) son idénticas en el dominio de \(f\).
\[ \]
Pista: La función \(y = [x]\) 
es la función parte entera: cada \(x\) lo relaciona con el mayor número entero igual o menor que
\(x\).}
{\(a = 0\),
\(b = 3\); 
\(\mathop{\mathrm{Dom}}(g) = (1;2)\)}{\(a = 3\),
\(b = 0\);
\(\mathop{\mathrm{Dom}}(g) = (1;2)\)}{\(a = 0\),
\(b = 4\);
\(\mathop{\mathrm{Dom}}(g) = (1;2)\)}{\(a = -3\),
\(b = 0\);
\(\mathop{\mathrm{Dom}}(g) = (1;2)\)}{}{}}
\End

\Data\ID{9000007202}
\Updated{2021-01-27 06:01:25}
\Level{C}
\SubArea{Linear functions}
\LANGen{
\Question{
Consider the function 
\[ 
 f(x) = [x] + 3
\]
on the domain \(\mathop{\mathrm{Dom}}(f) = (1;2)\).
Find the parameters \(a\) 
and \(b\) 
and a domain of the linear function 
\[ 
 g\colon y = ax + b
\]
which ensure that \(f\) 
and \(g\) 
are identical functions. 
\[ \]
Hint: The function \(y = [x]\) 
is a floor function: the largest integer less than or equal to
\(x\). For positive
\(x\) it is also called the
integer part of \(x\).}
{\(a = 0\),
\(b = 4\);
\(\mathop{\mathrm{Dom}}(g) = (1;2)\)}{\(a = 0\),
\(b = 3\);
\(\mathop{\mathrm{Dom}}(g) = (1;2)\)}{\(a = 3\),
\(b = 0\);
\(\mathop{\mathrm{Dom}}(g) = (1;2)\)}{\(a = -3\),
\(b = 0\);
\(\mathop{\mathrm{Dom}}(g) = (1;2)\)}{}{}}
\LANGpl{
\Question{
Rozważmy funkcję
\[ 
 f\colon y = [x] + 3
\]
na dziedzinie \(\mathop{\mathrm{Dom}}(f) = (1;2)\).
Określ wartości parametrów \(a\) 
i \(b\) 
i dziedzinę funkcji liniowej
\[ 
 g\colon y = ax + b,
\] które zapewnią, że funkcje \(f\) 
and \(g\) 
są funkcjami tożsamościowymi. 
\[ \]
Wskazówka: Funkcja \(y = [x]\) 
jest podłogą: największa liczba całkowita jest mniejsza lub równa 
\(x\). Dla dodatniego
\(x\) zwana jest również częścią ułamkową liczby całkowitej \(x\).}
{\(a = 0\),
\(b = 4\);
\(\mathop{\mathrm{Dom}}(g) = (1;2)\)}{\(a = 0\),
\(b = 3\);
\(\mathop{\mathrm{Dom}}(g) = (1;2)\)}{\(a = 3\),
\(b = 0\);
\(\mathop{\mathrm{Dom}}(g) = (1;2)\)}{\(a = -3\),
\(b = 0\);
\(\mathop{\mathrm{Dom}}(g) = (1;2)\)}{}{}}
\LANGcs{
\Question{
Je dána funkce \(f\colon y = [x] + 3\) a
platí \(D(f) = (1;2)\). Co musí
platit pro koeficienty \(a\),
\(b\) a definiční obor
lineární funkce \(g\colon y = ax + b\), aby se
rovnala zadané funkci \(f\)? 
\[ \]
Nápověda: Funkce \(y = [x]\) je celá
část čísla \(x\). Každému
reálnému číslu \(x\) 
přiřadí největší celé číslo, které je menší, nebo rovno
\(x\).}
{\(a = 0\ \wedge \ b = 4\ ; \ D(g) = (1;2)\)}{\(a = 0\ \wedge \ b = 3\ ; \ D(g) = (1;2)\)}{\(a = 3\ \wedge \ b = 0\ ; \ D(g) = (1;2)\)}{\(a = -3\ \wedge \ b = 0\ ; \ D(g) = (1;2)\)}{}{}}
\LANGsk{
\Question{
Daná je funkcia \(f\colon y = [x] + 3\) a
platí \(D(f) = (1;2)\). Čo musí
platiť pre koeficienty \(a\),
\(b\) a pre definičný obor
lineárnej funkcie \(g\colon y = ax + b\), aby sa
rovnala zadanej funkcii \(f\)? 
\[ \]
Poznámka: Funkcia \(y = [x]\) je cel
časť čísla \(x\). Každému
reálnemu číslu \(x\) 
priradí najväčšie celé číslo, ktoré je menšie alebo rovné
\(x\).}
{\(a = 0\ \wedge \ b = 4\ ; \ D(g) = (1;2)\)}{\(a = 0\ \wedge \ b = 3\ ; \ D(g) = (1;2)\)}{\(a = 3\ \wedge \ b = 0\ ; \ D(g) = (1;2)\)}{\(a = -3\ \wedge \ b = 0\ ; \ D(g) = (1;2)\)}{}{}}
\LANGes{
\Question{
Considera la función
\[ 
 f(x) = [x] + 3
\]
definida en el Dominio \(\mathop{\mathrm{Dom}}(f) = (1;2)\).
Halla los parámetros \(a\) 
y \(b\) 
y el Dominio de la función lineal
\[ 
 g\colon y = ax + b
\]
que garantizan que \(f\) 
y \(g\) 
son funciones idénticas
\[ \]
Pista: La función \(y = [x]\)  es la función  parte entera: cada \(x\) lo relaciona con el mayor número entero igual o menor que \(x\).}
{\(a = 0\),
\(b = 4\);
\(\mathop{\mathrm{Dom}}(g) = (1;2)\)}{\(a = 0\),
\(b = 3\);
\(\mathop{\mathrm{Dom}}(g) = (1;2)\)}{\(a = 3\),
\(b = 0\);
\(\mathop{\mathrm{Dom}}(g) = (1;2)\)}{\(a = -3\),
\(b = 0\);
\(\mathop{\mathrm{Dom}}(g) = (1;2)\)}{}{}}
\End

\Data\ID{9000007203}
\Updated{2021-01-27 06:01:40}
\Level{C}
\SubArea{Linear functions}
\LANGen{
\Question{
Consider the function 
\[ 
 f(x) =\mathop{ \mathrm{sgn}}\nolimits (x - 2)
\]
defined on \(\mathop{\mathrm{Dom}}(f) =\mathbb{R} ^{-}\). Find
the parameters \(a\) 
a \(b\) and
domain of the linear function 
\[ 
 g(x) = ax + b
\]
which ensure that \(f\) 
and \(g\) 
are identical functions. 
\[ \]
Hint: The function \(y =\mathop{ \mathrm{sgn}}\nolimits (x)\) 
is the sign function. The values of sign function is
\(1\) for each
positive \(x\),
\(- 1\) for each
negative \(x\) 
and \(0\) if
\(x = 0\).}
{\(a = 0\),
\(b = -1\);
\(\mathop{\mathrm{Dom}}(g) =\mathbb{R} ^{-}\)}{\(a = 0\),
\(b = 1\);
\(\mathop{\mathrm{Dom}}(g) =\mathbb{R} ^{+}\)}{\(a = 1\),
\(b = 0\);
\(\mathop{\mathrm{Dom}}(g) =\mathbb{R} ^{-}\)}{\(a = -1\),
\(b = 0\);
\(\mathop{\mathrm{Dom}}(g) =\mathbb{R} ^{+}\)}{}{}}
\LANGpl{
\Question{
Rozważ funkcję
\[ 
 f\colon y =\mathop{ \mathrm{sgn}}\nolimits (x - 2)
\]
określoną w \(\mathop{\mathrm{Dom}}(f) =\mathbb{R} ^{-}\). Wyznacz
wartości parametrów \(a\) 
i  \(b\) i
dziedzinę funkcji liniowej
\[ 
 g\colon y = ax + b,
\] które zapewnią, że funkcje \(f\) 
i \(g\) 
będą tożsamościowe.
\[ \]
 Wskazówka: Funkcja \(y =\mathop{ \mathrm{sgn}}\nolimits (x)\) 
jest funkcją znaku. Wartością funkcji znaku 
\(1\) dla każdego dodatniego \(x\),
\(- 1\) dla każdego
ujemnego \(x\) 
i \(0\) jeśli
\(x = 0\).}
{\(a = 0\),
\(b = -1\);
\(\mathop{\mathrm{Dom}}(g) =\mathbb{R} ^{-}\)}{\(a = 0\),
\(b = 1\);
\(\mathop{\mathrm{Dom}}(g) =\mathbb{R} ^{+}\)}{\(a = 1\),
\(b = 0\);
\(\mathop{\mathrm{Dom}}(g) =\mathbb{R} ^{-}\)}{\(a = -1\),
\(b = 0\);
\(\mathop{\mathrm{Dom}}(g) =\mathbb{R} ^{+}\)}{}{}}
\LANGcs{
\Question{
Je dána funkce 
\[f\colon y =\mathop{ \mathrm{sgn}}\nolimits (x - 2)
\] a
platí \(D(f) =\mathbb{R} ^{-}\). Co musí
platit pro koeficienty \(a\),
\(b\) a definiční obor
lineární funkce 
\[g\colon y = ax + b,\] 
aby se
rovnala zadané funkci \(f\)? 
\[ \]
Nápověda: Funkce \(y =\mathop{ \mathrm{sgn}}\nolimits (x)\) 
každému kladnému \(x\) 
přiřadí číslo \(1\),
číslu \(0\) přiřadí
\(0\) a zápornému
\(x\) přiřadí
číslo \(- 1\).}
{\(a = 0\ \wedge \ b = -1\ ; \ D(g) =\mathbb{R} ^{-}\)}{\(a = 0\ \wedge \ b = 1\ ; \ D(g) =\mathbb{R} ^{+}\)}{\(a = 1\ \wedge \ b = 0\ ; \ D(g) =\mathbb{R} ^{-}\)}{\(a = -1\ \wedge \ b = 0\ ; \ D(g) =\mathbb{R} ^{+}\)}{}{}}
\LANGsk{
\Question{
Daná je funkcia 
\[f\colon y =\mathop{ \mathrm{sgn}}\nolimits (x - 2)
\]
 a
platí \(D(f) =\mathbb{R} ^{-}\). Čo musí
platiť pre koeficienty \(a\),
\(b\) a definičný obor
lineárnej funkcie
\[g\colon y = ax + b,
\] aby sa
rovnala zadanej funkcii \(f\)? 
\[ \]
Pomôcka: Funkcia \(y =\mathop{ \mathrm{sgn}}\nolimits (x)\) 
každému kladnému \(x\) 
priradí číslo \(1\),
číslu \(0\) priradí
\(0\) a zápornému
\(x\) priradí
číslo \(- 1\).}
{\(a = 0\ \wedge \ b = -1\ ; \ D(g) =\mathbb{R} ^{-}\)}{\(a = 0\ \wedge \ b = 1\ ; \ D(g) =\mathbb{R} ^{+}\)}{\(a = 1\ \wedge \ b = 0\ ; \ D(g) =\mathbb{R} ^{-}\)}{\(a = -1\ \wedge \ b = 0\ ; \ D(g) =\mathbb{R} ^{+}\)}{}{}}
\LANGes{
\Question{
Considera la función
\[ 
 f(x) =\mathop{ \mathrm{sgn}}\nolimits (x - 2)
\]
definida en \(\mathop{\mathrm{Dom}}(f) =\mathbb{R} ^{-}\). Halla los parámetros
 \(a\) 
y \(b\) y
el Dominio de la función lineal
\[ 
 g(x) = ax + b
\]
que garantiza que \(f\) 
y \(g\) 
son funciones idénticas.
\[ \]
Pista: La función \(y =\mathop{ \mathrm{sgn}}\nolimits (x)\) 
es la función signo. Los valores de la función  signo son
\(1\) para cada \(x\) positivo,
\(- 1\) para cada  \(x\) negativo
y \(0\) si
\(x = 0\).}
{\(a = 0\),
\(b = -1\);
\(\mathop{\mathrm{Dom}}(g) =\mathbb{R} ^{-}\)}{\(a = 0\),
\(b = 1\);
\(\mathop{\mathrm{Dom}}(g) =\mathbb{R} ^{+}\)}{\(a = 1\),
\(b = 0\);
\(\mathop{\mathrm{Dom}}(g) =\mathbb{R} ^{-}\)}{\(a = -1\),
\(b = 0\);
\(\mathop{\mathrm{Dom}}(g) =\mathbb{R} ^{+}\)}{}{}}
\End

\Data\ID{9000007207}
\Updated{2021-01-27 06:01:53}
\Level{C}
\SubArea{Linear functions}
\LANGen{
\Question{
Identify a function which has the following three properties: it has at least one minimum or maximum, it is an increasing function and the range of this function is the set of all nonnegative numbers.}
{\(f(x) = 2x - 2\),
\(x\in [ 1;+\infty )\)}{\(f(x) = 2x + 2\),
\(x\in (-1;+\infty )\)}{\(f(x) = -2x + 2\),
\(x\in (-\infty ;1] \)}{\(f(x) = -2x - 2\),
\(x\in \mathbb{R}\)}{}{}}
\LANGpl{
\Question{
Która z poniższych funkcji posiada trzy następujące własności: posiada przynajmniej jedno minimum lub maksimum, jest funkcją rosnącą i zakres tej funkcji jest zbiorem wszystkich liczb nieujemnych.}
{\(f\colon y = 2x - 2\),
\(x\in [ 1;+\infty )\)}{\(f\colon y = 2x + 2\),
\(x\in (-1;+\infty )\)}{\(f\colon y = -2x + 2\),
\(x\in (-\infty ;1] \)}{\(f\colon y = -2x - 2\),
\(x\in \mathbb{R}\)}{}{}}
\LANGcs{
\Question{
Mezi následujícími funkcemi vyberte funkci, která má tyto vlastnosti: má alespoň jeden extrém (minimum nebo maximum), je rostoucí a její obor hodnot jsou nezáporná reálná čísla.}
{\(f\colon y = 2x - 2\), \( x\in \langle 1;+\infty )\)}{\(f\colon y = 2x + 2\), \( x\in(-1;+\infty )\)}{\(f\colon y = -2x + 2\), \( x\in (-\infty ;1\rangle \)}{\(f\colon y = -2x - 2\), \( x\in\mathbb{R}\)}{}{}}
\LANGsk{
\Question{
Vyberte funkciu, ktorá má nasledujúce tri vlastnosti: má aspoň jeden extrém (minimum alebo maximum), je rastúca a obor hodnôt tejto funkcie sú nezáporné reálne čísla.}
{\(f\colon y = 2x - 2\), \( x\in \langle 1;+\infty )\)}{\(f\colon y = 2x + 2\), \( x\in(-1;+\infty )\)}{\(f\colon y = -2x + 2\), \( x\in (-\infty ;1\rangle \)}{\(f\colon y = -2x - 2\), \( x\in\mathbb{R}\)}{}{}}
\LANGes{
\Question{
Identifica cuál de las funciones cumple las siguientes propiedades: tiene por lo mínimo un mínimo o máximo, es una función creciente y el Rango de la función es el conjunto de todos los números no negativos.}
{\(f(x) = 2x - 2\),
\(x\in [ 1;+\infty )\)}{\(f(x) = 2x + 2\),
\(x\in (-1;+\infty )\)}{\(f(x) = -2x + 2\),
\(x\in (-\infty ;1] \)}{\(f(x) = -2x - 2\),
\(x\in \mathbb{R}\)}{}{}}
\End

\Data\ID{9000007208}
\Updated{2021-01-27 06:01:59}
\Level{C}
\SubArea{Linear functions}
\LANGen{
\Question{
Paul's home is \(6\, \mathrm{km}\) from
the school. At the time \(t = 0\) 
Paul starts to walk from his home to the school along a straight street at a constant
velocity \(5\, \mathrm{km}/\mathrm{h}\).
Find the function which describes Paul's remaining distance to the school as a
function of time.}
{\(s = 6 - 5t\)}{\(s = 5t - 6\)}{\(s = 5t\)}{\(s = 5t + 6\)}{}{}}
\LANGpl{
\Question{
Dom Pawła znajduję się w odległości \(6\, \mathrm{km}\) od 
szkoły. W czasie \(t = 0\) 
Paweł zaczyna iść z domu do szkoły prostą ulicą ze stałą prędkością
 \(5\, \mathrm{km}/\mathrm{h}\).
Znajdź funkcję, która określa odległość do szkoły jaką musi jeszcze przebyć Paweł jako funkcję czasu.}
{\(s = 6 - 5t\)}{\(s = 5t - 6\)}{\(s = 5t\)}{\(s = 5t + 6\)}{}{}}
\LANGcs{
\Question{
Tomáš bydlí \(6\, \mathrm{km}\) 
od školy. Vyberte rovnici funkce, která bude vyjadřovat závislost aktuální
Tomášovy vzdálenosti od školy na době jeho chůze, předpokládáme-li, že
Tomáš půjde z domova do školy rovnoměrným přímočarým pohybem
rychlostí \(5\, \mathrm{km}/\mathrm{h}\).}
{\(s = 6 - 5t\)}{\(s = 5t - 6\)}{\(s = 5t\)}{\(s = 5t + 6\)}{}{}}
\LANGsk{
\Question{
Tomáš býva \(6\, \mathrm{km}\) od školy. 
V čase \(t = 0\) Tomáš odchádza z domu do školy konštantnou rýchlosťou \(5\, \mathrm{km}/\mathrm{h}\).
Určte predpis funkcie, ktorá vyjadruje závislosť 
Tomášovej vzdialenosti od školy na čase jeho chôdze.}
{\(s = 6 - 5t\)}{\(s = 5t - 6\)}{\(s = 5t\)}{\(s = 5t + 6\)}{}{}}
\LANGes{
\Question{
Pablo vive a \(6\, \mathrm{km}\) de
su escuela. En el tiempo \(t = 0\) 
Pablo empieza a andar de su casa a la escuela por una calle recta a una velocidad constante \(5\, \mathrm{km}/\mathrm{h}\).
Halla la función que describe la distancia que le queda a Pablo para llegar a la escuela en función del tiempo.}
{\(s = 6 - 5t\)}{\(s = 5t - 6\)}{\(s = 5t\)}{\(s = 5t + 6\)}{}{}}
\End

\Data\ID{9000007209}
\Updated{2021-01-27 06:01:11}
\Level{C}
\SubArea{Linear functions}
\IMG{000072_09-figure0.pdf}
\LANGen{
\Question{
A current-voltage characteristics is shown in the graph. Find the current
\(I\) as a function of
the voltage \(U\).}
{\(I = \frac{2} 
{3}U -\frac{4}{3};U\in [2;\infty)
\)}{\(I = \frac{3} 
{2}U - 2;U\in [2;\infty)
\)}{\(I = \frac{3} 
{2}U + 2;U\in [2;\infty)
\)}{\(I = \frac{2} 
{3}U + 2;U\in [2;\infty)
\)}{}{}}
\LANGpl{
\Question{
Wykres przedstawia charakterystykę napięci prądu.  Określ prąd
\(I\) jako funkcję napięcia
\(U\).}
{\(I = \frac{2} 
{3}U -\frac{4} 
{3};U\in \langle 2,\infty)
\)}{\(I = \frac{3} 
{2}U - 2;U\in \langle 2,\infty)
\)}{\(I = \frac{3} 
{2}U + 2;U\in \langle 2,\infty)
\)}{\(I = \frac{2} 
{3}U + 2;U\in \langle 2,\infty)
\)}{}{}}
\LANGcs{
\Question{
Voltampérová charakteristika elektrolytu má průběh, který
je graficky znázorněn na obrázku. Vyjádřete proud jako
funkci napětí.}
{\(I = \frac{2} 
{3}U -\frac{4} 
{3};U\in \langle 2,\infty)
\)}{\(I = \frac{3} 
{2}U - 2;U\in \langle 2,\infty)
\)}{\(I = \frac{3} 
{2}U + 2;U\in \langle 2,\infty)
\)}{\(I = \frac{2} 
{3}U + 2;U\in \langle 2,\infty)
\)}{}{}}
\LANGsk{
\Question{
Voltampérová charakteristika elektrolytu má priebeh, ktorý
je graficky znázornený na obrázku. Vyjadrite prúd ako
funkciu napätia.}
{\(I = \frac{2} 
{3}U -\frac{4} 
{3};U\in \langle 2,\infty)
\)}{\(I = \frac{3} 
{2}U - 2;U\in \langle 2,\infty)
\)}{\(I = \frac{3} 
{2}U + 2;U\in \langle 2,\infty)
\)}{\(I = \frac{2} 
{3}U + 2;U\in \langle 2,\infty)
\)}{}{}}
\LANGes{
\Question{
En la gráfica se representa la relación entre la corriente y el voltaje.Encuentra la corriente
\(I\) en función del voltaje \(U\).}
{\(I = \frac{2} 
{3}U -\frac{4}{3};U\in [2;\infty)
\)}{\(I = \frac{3} 
{2}U - 2;U\in [2;\infty)
\)}{\(I = \frac{3} 
{2}U + 2;U\in [2;\infty)
\)}{\(I = \frac{2} 
{3}U + 2;U\in [2;\infty)
\)}{}{}}
\End

\Data\ID{9000007210}
\Updated{2021-04-05 06:04:25}
\Level{C}
\SubArea{Linear functions}
\LANGen{
\Question{
Jane has to reach the opposite shore of a lake. She has three possibilities how to
cross. She can use her own boat, start immediately and sail at the velocity
 \(4\, \mathrm{km}/\mathrm{h}\). Another option is to wait for her friend, Peter, with a
 faster boat. Peter's boat is capable to sail at the velocity
 \(10\, \mathrm{km}/\mathrm{h}\).
 However, his boat will be available in
 \(1.5\) 
 hours from now. The last option is to use the regular passenger
 transportation line which will departure in
 \(2.25\) hours
 from now and sails at the speed \(20\, \mathrm{km}/\mathrm{h}\). Find the interval of distances to the opposite shore for which the fastest option is to
use Peter's boat.}
{between \(10\) 
and \(15\) 
kilometers}{up to \(10\) 
kilometers}{between \(15\) 
and \(20\) 
kilometers}{larger that \(20\) 
kilometers}{}{}}
\LANGpl{
\Question{
Janek musi dostać się na przeciwną stronę jeziora. Ma trzy możliwości, żeby tam dotrzeć.  Może użyć własnej łodzi, musi zacząć natychmiast i żeglować z prędkością 
 \(4\, \mathrm{km}/\mathrm{h}\). Może poczekać na przyjaciela Piotra, który ma szybszą łódź. Łódź Piotra może poruszać się z prędkością
 \(10\, \mathrm{km}/\mathrm{h}\).
 Jednakże jego łódź będzie dostępna dopiero za 
 \(1.5\) 
 godziny. Ostatnią możliwością jest użycie zwykłej linii  pasażerskiej, która odpływa za
 \(2.25\) godziny i porusza się z prędkością  \(20\, \mathrm{km}/\mathrm{h}\). Znajdź przedział odległości do przeciwnego brzegu jeziora, dla którego najszybszą opcją jest użycie łodzi Piotra.}
{pomiędzy\(10\) 
i \(15\) 
kilometrów}{do \(10\) 
kilometrów}{pomiędzy \(15\) 
i \(20\) 
kilometrów}{więcej niż \(20\) 
kilometrów}{}{}}
\LANGcs{
\Question{
Petr se potřebuje dostat přes jezero. Zvažuje tři možnosti. Může nasednout
do vlastní loďky a vyplout okamžitě, ale jeho průměrná rychlost bude pouze
\(4\, \mathrm{km}/\mathrm{h}\).
Nebo může požádat kamaráda, aby ho tam zavezl. Kamarád
má rychlejší loď, která může plout průměrnou rychlostí
\(10\, \mathrm{km}/\mathrm{h}\), ale mohli by vyplout
až za \(1{,}5\) hodiny.
Poslední Petrovou možností je využít pravidelnou lodní linku, která vyplouvá za
\(2{,}25\) hodiny. V tomto případě
by cestoval rychlostí \(20\, \mathrm{km}/\mathrm{h}\).
V jaké vzdálenosti musí být přístav na druhém břehu, aby bylo
nejvýhodnější použít kamarádovu loď?}
{mezi \(10\) 
a \(15\) 
kilometry}{do \(10\) 
kilometrů}{mezi \(15\) 
a \(20\) 
kilometry}{větší než \(20\) 
kilometrů}{}{}}
\LANGsk{
\Question{
Jana sa potrebuje dostať do prístavu na druhej strane jazera. Má tri nasledovné možnosti. 
Môže použiť vlastnú loď a vyraziť ihneď priemernou rýchlosťou \(4\, \mathrm{km}/\mathrm{h}\).
Druhá možnosť je počkať na kamaráta Petra, ktorý má rýchlejšiu loď. Petrova loď pláva priemernou rýchlosťou \(10\, \mathrm{km}/\mathrm{h}\), ale môže vyraziť až za \(1{,}5\) hodiny.
Posledná možnosť je využiť pravidelné lodné spojenie, ktoré vyráža za \(2{,}25\) hodiny priemernou rýchlosťou \(20\, \mathrm{km}/\mathrm{h}\).
V akej vzdialenosti musí byť prístav na druhej strane jazera, aby bolo najvýhodnejšie použiť Petrovu loď?}
{medzi \(10\) 
a \(15\) 
kilometrami}{do \(10\) 
kilometrov}{medzi \(15\) 
a \(20\) 
kilometrami}{viac než \(20\) 
kilometrov}{}{}}
\LANGes{
\Question{
Pedro necesita atravesar un lago y tiene tres opciones. Puede ir en su barco e irse inmediatamente
pero su velocidad media será solamente \(4\, \mathrm{km}/\mathrm{h}\).
Otra opción es pedirle a un amigo que le lleve en su barco, que puede ir con la velocidad media  \(10\, \mathrm{km}/\mathrm{h}\) pero tendrían que salir dentro de \(1{,}5\) horas.
La última opción es ir utilizando un barco público que sale dentro de \(2{,}25\) horas y su velocidad media es  \(20\, \mathrm{km}/\mathrm{h}\).
¿Para qué distancia entre las dos orillas sería mejor la opción del barco del amigo de Pedro?}
{entre \(10\) 
y \(15\) 
kilómetros}{menos que \(10\) 
kilómetros}{entre \(15\) 
y \(20\) 
kilómetros}{más que \(20\) 
kilómetros}{}{}}
\End

\Data\ID{9000007809}
\Updated{2021-04-05 07:04:00}
\Level{C}
\SubArea{Linear functions}
\LANGen{
\Question{
The price of a goods in a shop is \(\$15\) 
per item. The Internet price in an e-shop is cheaper by
\(\$2\) per item. The shipping
cost of the e-shop is \(\$125\).
What is the minimal number of items, which makes the total cost for a transaction
smaller in the e-shop?}
{\(63\)}{\(9\)}{\(62\)}{\(125\)}{\(126\)}{}}
\LANGpl{
\Question{
Cena towaru w sklepie wynosi \(\$15\) 
za sztukę. Cena w sklepie internetowym jest niższa o
\(\$2\) za sztukę. Koszt wysyłki ze sklepu internetowego wynosi \(\$125\).
Jaka jest minimalna liczba przedmiotów, która spowoduje, że łączny koszt transakcji będzie mniejszej w sklepie internetowym?}
{\(63\)}{\(9\)}{\(62\)}{\(125\)}{\(126\)}{}}
\LANGcs{
\Question{
Zboží v obchodě stojí \(15\) Kč
za jeden kus. Na internetu se dá stejné zboží pořídit o
\(2\) Kč
za kus levněji, ale je třeba připočítat poštovné a balné, které činí
\(125\) Kč.
Jaký musí být minimální počet kusů v objednávce, aby byl nákup na
internetu výhodnější?}
{\(63\)}{\(9\)}{\(62\)}{\(125\)}{\(126\)}{}}
\LANGsk{
\Question{
Tovar v obchode stojí \(15\) EURO
za jeden kus. Na internete sa dá rovnaký tovar kúpiť o
\(2\) EURO
za kus lacnejšie. Treba však pripočítať poštovné a balné, ktoré je
\(125\) EURO.
Koľko musíme minimálne objednať kusov tovaru, aby bol nákup na internete výhodnejší?}
{\(63\)}{\(9\)}{\(62\)}{\(125\)}{\(126\)}{}}
\LANGes{
\Question{
El precio de todos los artículos de una tienda, si los compras allí presencialmente es de \(\$15\) 
. El precio en la tienda online es
\(\$2\) menos por artículo pero los gastos de envío son \(\$125\).
¿Cuántos artículos necesitaríamos comprar para que el precio total sea menor conprando online que en la tienda física?}
{\(63\)}{\(9\)}{\(62\)}{\(125\)}{\(126\)}{}}
\End

\Data\ID{9000007810}
\Updated{2021-04-05 07:04:43}
\Level{C}
\SubArea{Linear functions}
\LANGen{
\Question{
A fuel tank in a car has the capacity \(40\) 
litres. The current volume of the fuel in the fuel tank is
\(6\) litres. The speed
of fuelling is \(1\) litre
of gasoline each \(3\) 
seconds. Find the function which describes the volume of the gasoline in the fuel tank (in litres)
as a function of time (in seconds).}
{\(V = \frac{1} 
{3}t + 6,\ t\in [ 0;102] \)}{\(V = 3t + 6,\ t\in [ 0;102] \)}{\(V = 3t + 6,\ t\in [ 0;40] \)}{\(V = 3t + 6,\ t\in \mathbb{R}_{0}^{+}\)}{\(V = \frac{1} 
{3}t + 6,\ t\in [ 0;40] \)}{}}
\LANGpl{
\Question{
Zbiornik paliwa w samochodzie ma pojemność \(40\) 
litrów. Obecna objętość paliwa w zbiorniku paliwa wynosi
\(6\) litrów.  Prędkość tankowania wynosi \(1\) litr benzyny co 
 \(3\) sekundy. Wyznacz funkcję, która opisuje objętość benzyny w zbiorniku paliwa jako funkcję  czasu.}
{\(V = \frac{1} 
{3}t + 6,\ t\in [ 0;102] \)}{\(V = 3t + 6,\ t\in [ 0;102] \)}{\(V = 3t + 6,\ t\in [ 0;40] \)}{\(V = 3t + 6,\ t\in \mathbb{R}_{0}^{+}\)}{\(V = \frac{1} 
{3}t + 6,\ t\in [ 0;40] \)}{}}
\LANGcs{
\Question{
V nádrži automobilu o celkové kapacitě
\(40\) litrů zůstalo
pouze \(6\) 
litrů benzínu. Při tankování přitéká
\(1\) litr benzínu
každé \(3\) sekundy.
Určete předpis funkce, která vyjadřuje závislost množství benzínu v nádrži
(\(V \) - v litrech) na čase
(\(t\) - v sekundách).}
{\(V = \frac{1} 
{3}t + 6,\ t\in \langle 0;102\rangle \)}{\(V = 3t + 6,\ t\in \langle 0;102\rangle \)}{\(V = 3t + 6,\ t\in \langle 0;40\rangle \)}{\(V = 3t + 6,\ t\in \mathbb{R}_{0}^{+}\)}{\(V = \frac{1} 
{3}t + 6,\ t\in \langle 0;40\rangle \)}{}}
\LANGsk{
\Question{
Nádrž auta má celkovú kapacitu \(40\) litrov ale zostalo v nej len \(6\) litrov benzínu. 
Počas tankovania priteká \(1\) liter benzínu každé \(3\) sekundy.
Určte predpis funkcie, ktorá vyjadruje závislosť množstva benzínu v nádrži
(\(V \) - v litroch) na čase
(\(t\) - v sekundách).}
{\(V = \frac{1} 
{3}t + 6,\ t\in \langle 0;102\rangle \)}{\(V = 3t + 6,\ t\in \langle 0;102\rangle \)}{\(V = 3t + 6,\ t\in \langle 0;40\rangle \)}{\(V = 3t + 6,\ t\in \mathbb{R}_{0}^{+}\)}{\(V = \frac{1} 
{3}t + 6,\ t\in \langle 0;40\rangle \)}{}}
\LANGes{
\Question{
Un depósito de combustible de un coche tiene una capacidad de \(40\) litros.
El volumen de combustible en el coche en este momento es 
\(6\) litros. La velocidad de consumo es \(1\) litros de gasolina cada \(3\) 
segundos. Halla la función que describe el volumen de gasolina en el depósito (en litros) en función del tiempo (en segundos).}
{\(V = \frac{1} 
{3}t + 6,\ t\in [ 0;102] \)}{\(V = 3t + 6,\ t\in [ 0;102] \)}{\(V = 3t + 6,\ t\in [ 0;40] \)}{\(V = 3t + 6,\ t\in \mathbb{R}_{0}^{+}\)}{\(V = \frac{1} 
{3}t + 6,\ t\in [ 0;40] \)}{}}
\End

\Data\ID{9000009301}
\Updated{2021-04-05 07:04:44}
\Level{C}
\SubArea{Linear functions}
\LANGen{
\Question{
An automatic machine produces \(12\) 
components per minute and stores them in a box with capacity
\(1\: 500\) 
components. The machine starts with an initial amount of
\(240\) 
components in the box. In what time will be the box full?}
{\(1\, \mathrm{h}\) 
\(45\, \mathrm{min}\)}{\(1\, \mathrm{h}\) 
\(55\, \mathrm{min}\)}{\(2\, \mathrm{h}\) 
\(5\, \mathrm{min}\)}{\(2\, \mathrm{h}\) 
\(15\, \mathrm{min}\)}{}{}}
\LANGpl{
\Question{
Automatyczna maszyna produkuje \(12\) 
elementów na minutę i umieszcza je w pudełku o pojemności 
\(1\: 500\) 
elementów. Maszyna zaczyna od 
\(240\) 
elementów w pudełku. W jakim czasie pudełko zapełni się?}
{\(1\, \mathrm{h}\) 
\(45\, \mathrm{min}\)}{\(1\, \mathrm{h}\) 
\(55\, \mathrm{min}\)}{\(2\, \mathrm{h}\) 
\(5\, \mathrm{min}\)}{\(2\, \mathrm{h}\) 
\(15\, \mathrm{min}\)}{}{}}
\LANGcs{
\Question{
Automat vyrobí \(12\) 
součástek za minutu a ukládá je do zásobníku, jehož kapacita je
\(1\: 500\) kusů. Na začátku
směny je v zásobníku \(240\) 
kusů. Za jak dlouho bude zásobník plný?}
{\(1\, \mathrm{h}\) 
\(45\, \mathrm{min}\)}{\(1\, \mathrm{h}\) 
\(55\, \mathrm{min}\)}{\(2\, \mathrm{h}\) 
\(5\, \mathrm{min}\)}{\(2\, \mathrm{h}\) 
\(15\, \mathrm{min}\)}{}{}}
\LANGsk{
\Question{
Automat vyrobí \(12\) 
súčiastok za minútu a ukladá ich do zásobníka, ktorého kapacita je
\(1\: 500\) kusov. Automat začína pracovať s počtom \(240\) 
kusov súčiastok v zásobníku. Ako dlho bude trvať, kým bude zásobník plný?}
{\(1\, \mathrm{h}\) 
\(45\, \mathrm{min}\)}{\(1\, \mathrm{h}\) 
\(55\, \mathrm{min}\)}{\(2\, \mathrm{h}\) 
\(5\, \mathrm{min}\)}{\(2\, \mathrm{h}\) 
\(15\, \mathrm{min}\)}{}{}}
\LANGes{
\Question{
Una máquina automática produce \(12\) 
componentes por minuto y los pone en una caja cuya capacidad es de  
\(1\: 500\) 
componentes. Si la máquina empieza con 
\(240\) 
componentes en la caja. ¿Cuánto tardará en llenar la caja?}
{\(1\, \mathrm{h}\) 
\(45\, \mathrm{min}\)}{\(1\, \mathrm{h}\) 
\(55\, \mathrm{min}\)}{\(2\, \mathrm{h}\) 
\(5\, \mathrm{min}\)}{\(2\, \mathrm{h}\) 
\(15\, \mathrm{min}\)}{}{}}
\End

\Data\ID{9000009302}
\Updated{2021-04-05 07:04:59}
\Level{C}
\SubArea{Linear functions}
\LANGen{
\Question{
An automatic machine produces \(12\) 
components per minute and stores them in a box with capacity
\(1\: 500\) 
components. The machine starts with an initial amount of
\(240\) 
components in the box. In what time will the box contain
\(1\: 020\) 
components?}
{\(1\, \mathrm{h}\) 
\(5\, \mathrm{min}\)}{\(55\, \mathrm{min}\)}{\(1\, \mathrm{h}\)}{\(1\, \mathrm{h}\) 
\(10\, \mathrm{min}\)}{}{}}
\LANGpl{
\Question{
Automatyczna maszyna produkuje \(12\) 
elementów na minutę i umieszcza je w pudełku o pojemności
\(1\: 500\) 
elementów. Maszyna zaczyna z początkową ilością
\(240\) 
elementów w pudełku. W jakim czasie w pudełku znajdzie się
\(1\: 020\) 
elementów?}
{\(1\, \mathrm{h}\) 
\(5\, \mathrm{min}\)}{\(55\, \mathrm{min}\)}{\(1\, \mathrm{h}\)}{\(1\, \mathrm{h}\) 
\(10\, \mathrm{min}\)}{}{}}
\LANGcs{
\Question{
Automat vyrobí \(12\) 
součástek za minutu a ukládá je do zásobníku, jehož kapacita je
\(1\: 500\) 
kusů. Na začátku směny je v zásobníku
\(240\) kusů. Za jak dlouho
bude v zásobníku \(1\: 020\) 
součástek?}
{\(1\, \mathrm{h}\) 
\(5\, \mathrm{min}\)}{\(55\, \mathrm{min}\)}{\(1\, \mathrm{h}\)}{\(1\, \mathrm{h}\) 
\(10\, \mathrm{min}\)}{}{}}
\LANGsk{
\Question{
Automat vyrobí \(12\) súčiastok za minútu a ukladá ich do zásobníka, ktorého kapacita je \(1\: 500\) kusov. 
Na začiatku zmeny je v zásobníku \(240\) kusov. 
Za aký dlhý čas bude v zásobníku \(1\: 020\) súčiastok?}
{\(1\, \mathrm{h}\) 
\(5\, \mathrm{min}\)}{\(55\, \mathrm{min}\)}{\(1\, \mathrm{h}\)}{\(1\, \mathrm{h}\) 
\(10\, \mathrm{min}\)}{}{}}
\LANGes{
\Question{
Una máquina automática produce \(12\) 
componentes por minuto y los pone en una caja con una capacidad de
\(1\: 500\) 
componentes. Si la máquina empieza con
\(240\) 
componentes en la caja. ¿Cuánto tardará hasta que hayan
\(1\: 020\) 
componentes en la caja?}
{\(1\, \mathrm{h}\) 
\(5\, \mathrm{min}\)}{\(55\, \mathrm{min}\)}{\(1\, \mathrm{h}\)}{\(1\, \mathrm{h}\) 
\(10\, \mathrm{min}\)}{}{}}
\End

\Data\ID{9000009303}
\Updated{2021-04-05 07:04:58}
\Level{C}
\SubArea{Linear functions}
\LANGen{
\Question{
A tank contains \(1\: 000\) 
litres of petrol. The petrol escapes at a constant speed
\(20\) litres
per minute. In what time will be the tank empty?}
{\(50\, \mathrm{min}\)}{\(60\, \mathrm{min}\)}{\(40\, \mathrm{min}\)}{\(30\, \mathrm{min}\)}{}{}}
\LANGpl{
\Question{
Zbiornik zawiera \(1\: 000\) 
litrów benzyny. Benzyna wycieka ze stałą prędkością
\(20\) litrów na minutę.
W jakim czasie zbiornik opróżni się?}
{\(50\, \mathrm{min}\)}{\(60\, \mathrm{min}\)}{\(40\, \mathrm{min}\)}{\(30\, \mathrm{min}\)}{}{}}
\LANGcs{
\Question{
V cisterně je \(1\: 000\) 
litrů nafty. Nafta vytéká stálou rychlostí
\(20\) litrů
za minutu. Za jak dlouho bude cisterna prázdná?}
{\(50\, \mathrm{min}\)}{\(60\, \mathrm{min}\)}{\(40\, \mathrm{min}\)}{\(30\, \mathrm{min}\)}{}{}}
\LANGsk{
\Question{
V nádrži je \(1\: 000\)  litrov nafty. Nafta vyteká stálou rýchlosťou \(20\) litrov za minútu. 
Ako dlho bude trvať vyprázdnenie nádrže?}
{\(50\, \mathrm{min}\)}{\(60\, \mathrm{min}\)}{\(40\, \mathrm{min}\)}{\(30\, \mathrm{min}\)}{}{}}
\LANGes{
\Question{
Una cisterna contiene \(1\: 000\) 
litros de petróleo. El petróleo sale a una velocidad constante de 
\(20\) litros por minuto. ¿Cuánto tardará la cisterna en estar vacía?}
{\(50\, \mathrm{min}\)}{\(60\, \mathrm{min}\)}{\(40\, \mathrm{min}\)}{\(30\, \mathrm{min}\)}{}{}}
\End

\Data\ID{9000009304}
\Updated{2021-04-05 07:04:08}
\Level{C}
\SubArea{Linear functions}
\LANGen{
\Question{
A tank contains \(1\: 000\) 
litres of petrol. The petrol escapes at a constant speed
\(20\) litres per minute. In what
time will there be just \(200\) 
litres of the petrol in the tank?}
{\(40\, \mathrm{min}\)}{\(10\, \mathrm{min}\)}{\(20\, \mathrm{min}\)}{\(30\, \mathrm{min}\)}{}{}}
\LANGpl{
\Question{
Zbiornik zawiera \(1\: 000\) 
litrów benzyny. Benzyna wycieka ze stałą prędkością 
\(20\) litrów na minutę. W jakim czasie w zbiorniku pozostanie tylko \(200\) litrów
paliwa?}
{\(40\, \mathrm{min}\)}{\(10\, \mathrm{min}\)}{\(20\, \mathrm{min}\)}{\(30\, \mathrm{min}\)}{}{}}
\LANGcs{
\Question{
V cisterně je \(1\: 000\) 
litrů nafty. Nafta vytéká stálou rychlostí
\(20\) litrů za minutu. Za jak
dlouho bude v cisterně \(200\) 
litrů nafty?}
{\(40\, \mathrm{min}\)}{\(10\, \mathrm{min}\)}{\(20\, \mathrm{min}\)}{\(30\, \mathrm{min}\)}{}{}}
\LANGsk{
\Question{
V nádrži je \(1\: 000\) 
litrov nafty. Nafta vyteká stálou rýchlosťou
\(20\) litrov za minútu. Za aký čas
bude v nádrži \(200\) 
litrov nafty?}
{\(40\, \mathrm{min}\)}{\(10\, \mathrm{min}\)}{\(20\, \mathrm{min}\)}{\(30\, \mathrm{min}\)}{}{}}
\LANGes{
\Question{
Una cisterna contiene \(1\: 000\) 
litros de petróleo. El petróleo sale a velocidad
\(20\) litros por minuto. ¿Cuántos minutos han de pasar para que queden solamente \(200\) 
litros de petróleo en la cisterna?}
{\(40\, \mathrm{min}\)}{\(10\, \mathrm{min}\)}{\(20\, \mathrm{min}\)}{\(30\, \mathrm{min}\)}{}{}}
\End

\Data\ID{9000009305}
\Updated{2021-04-05 07:04:59}
\Level{C}
\SubArea{Linear functions}
\LANGen{
\Question{
Anne decided to make a bicycle trip with his friend which lives
\(10\, \mathrm{km}\) 
from Anne house. Anne went from her house to the house of her friend first.
Then they started to measure the time and went on a constant velocity
\(18\, \mathrm{km}/\mathrm{h}\).
In what time will be the total distance traveled by Anne equal to
\(34\, \mathrm{km}\)?}
{\(1\, \mathrm{h}\) 
\(20\, \mathrm{min}\)}{\(1\, \mathrm{h}\) 
\(58\, \mathrm{min}\)}{\(2\, \mathrm{h}\) 
\(26\, \mathrm{min}\)}{\(2\, \mathrm{h}\) 
\(30\, \mathrm{min}\)}{}{}}
\LANGpl{
\Question{
Ania postanowiła wybrać się na wycieczkę rowerową ze swoją przyjaciółką Zosią, która mieszka  
\(10\, \mathrm{km}\) 
od domu Ani. Najpierw Ania przebyła trasę do domu przyjaciółki. Następnie wspólnie zaczęły mierzyć czas i poruszały się ze stałą prędkością
\(18\, \mathrm{km}/\mathrm{h}\).
W jakim czasie całkowita odległość pokonana przez Anię wyniesie
\(34\, \mathrm{km}\)?}
{\(1\, \mathrm{h}\) 
\(20\, \mathrm{min}\)}{\(1\, \mathrm{h}\) 
\(58\, \mathrm{min}\)}{\(2\, \mathrm{h}\) 
\(26\, \mathrm{min}\)}{\(2\, \mathrm{h}\) 
\(30\, \mathrm{min}\)}{}{}}
\LANGcs{
\Question{
Martina si domluvila cyklistický výlet s kamarádem Pavlem, který bydlí
\(10\, \mathrm{km}\) 
od Martinina domu. Martina nejprve jela z domu k Pavlovi, kde si začali
měřit čas a rychlost. Od Pavlova domu jeli společně konstantní rychlostí 
\(18\, \mathrm{km}/\mathrm{h}\).
Za jak dlouho od odjezdu od Pavlova domu bude mít Martina ujeto \(34\,
\mathrm{km}\)?}
{\(1\, \mathrm{h}\) 
\(20\, \mathrm{min}\)}{\(1\, \mathrm{h}\) 
\(58\, \mathrm{min}\)}{\(2\, \mathrm{h}\) 
\(26\, \mathrm{min}\)}{\(2\, \mathrm{h}\) 
\(30\, \mathrm{min}\)}{}{}}
\LANGsk{
\Question{
Janka si dohodla cyklistický výlet s kamarátom Petrom, ktorý býva \(10\, \mathrm{km}\) od Jankinho domu.
Janka najprv išla na bicykli z domu k Petrovi. Odtiaľ si spoločne začali merať čas, pri konštantnej rýchlosti \(18\, \mathrm{km}/\mathrm{h}\).
Za aký čas od odchodu z Petrovho domu bude mať Janka najazdených \(34\,\mathrm{km}\)?}
{\(1\, \mathrm{h}\) 
\(20\, \mathrm{min}\)}{\(1\, \mathrm{h}\) 
\(58\, \mathrm{min}\)}{\(2\, \mathrm{h}\) 
\(26\, \mathrm{min}\)}{\(2\, \mathrm{h}\) 
\(30\, \mathrm{min}\)}{}{}}
\LANGes{
\Question{
Anne decidió ir de viaje en bici con un amigo que vive a 
\(10\, \mathrm{km}\) 
de la casa de Anna. Anna fue primero a la casa de su amigo. Luego empezaron a medir el tiempo y fueron a una velocidad constante de
\(18\, \mathrm{km}/\mathrm{h}\).
¿Cuánto tiempo tardará Anna en recorrer la distancia total de
\(34\, \mathrm{km}\)?}
{\(1\, \mathrm{h}\) 
\(20\, \mathrm{min}\)}{\(1\, \mathrm{h}\) 
\(58\, \mathrm{min}\)}{\(2\, \mathrm{h}\) 
\(26\, \mathrm{min}\)}{\(2\, \mathrm{h}\) 
\(30\, \mathrm{min}\)}{}{}}
\End

\Data\ID{9000009306}
\Updated{2021-04-05 07:04:56}
\Level{C}
\SubArea{Linear functions}
\LANGen{
\Question{
Anne decided to make a bicycle trip with his friend which lives
\(10\, \mathrm{km}\) 
from Anne house. Anne went from her house to the house of her friend first.
Then they started to measure the time and went on a constant velocity
\(18\, \mathrm{km}/\mathrm{h}\) for
\(2\, \mathrm{h}\) 
\(10\, \mathrm{min}\).
What is the total distance traveled by Anne?}
{\(49\, \mathrm{km}\)}{\(39\, \mathrm{km}\)}{\(35\, \mathrm{km}\)}{\(45\, \mathrm{km}\)}{}{}}
\LANGpl{
\Question{
Anna postanowiła wybrać się na wycieczkę rowerową z przyjaciółką, która mieszka
\(10\, \mathrm{km}\) 
od domu Anny. Anna przebyła najpierw trasę z własnego domu do domu przyjaciółki. 
Następnie wraz z przyjaciółka zaczęły mierzyć czas i poruszały się ze stałą prędkością
\(18\, \mathrm{km}/\mathrm{h}\) przebyli wybraną trasę w ciągu 
\(2\, \mathrm{h}\) 
\(10\, \mathrm{min}\).
Ile wyniosła całkowita odległość pokonana przez Annę.}
{\(49\, \mathrm{km}\)}{\(39\, \mathrm{km}\)}{\(35\, \mathrm{km}\)}{\(45\, \mathrm{km}\)}{}{}}
\LANGcs{
\Question{
Martina si domluvila cyklistický výlet s kamarádem Pavlem, který bydlí
\(10\, \mathrm{km}\) od Martinina domu. Martina nejprve jela z domu k 
Pavlovi, kde si začali měřit čas a rychlost. Od Pavlova domu jeli 
společně \(2\) hodiny a \(10\) minut konstantní rychlostí 
\(18\, \mathrm{km}/\mathrm{h}\). Jakou celkovou vzdálenost Martina ujela?}
{\(49\, \mathrm{km}\)}{\(39\, \mathrm{km}\)}{\(35\, \mathrm{km}\)}{\(45\, \mathrm{km}\)}{}{}}
\LANGsk{
\Question{
Janka si dohodla cyklistický výlet s kamarátom Petrom, ktorý býva \(10\, \mathrm{km}\) od Jankinho domu.
Janka najprv išla na bicykli z domu k Petrovi. Odtiaľ si spoločne začali merať čas, pri konštantnej rýchlosti \(18\, \mathrm{km}/\mathrm{h}\).
Od Petrovho domu bicyklovali spoločne \(2\) hodiny a \(10\) minút. Akú celkovú vzdialenosť Janka najazdila?}
{\(49\, \mathrm{km}\)}{\(39\, \mathrm{km}\)}{\(35\, \mathrm{km}\)}{\(45\, \mathrm{km}\)}{}{}}
\LANGes{
\Question{
Anna decidió ir de viaje en bici con un amigo que vive 
\(10\, \mathrm{km}\) 
de casa de Anna.  Anna fue de su casa a la casa de su amigo primero. 
Luego empezaron a medir el tiempo y fueron a una velocidad constante 
\(18\, \mathrm{km}/\mathrm{h}\) durante
\(2\, \mathrm{h}\) 
\(10\, \mathrm{min}\).
¿Qué distancia ha viajado Anna en total?}
{\(49\, \mathrm{km}\)}{\(39\, \mathrm{km}\)}{\(35\, \mathrm{km}\)}{\(45\, \mathrm{km}\)}{}{}}
\End

\Data\ID{9000009307}
\Updated{2021-04-05 07:04:19}
\Level{C}
\SubArea{Linear functions}
\LANGen{
\Question{
The sound velocity at the temperature
\(0\, ^{\circ } \mathrm{C}\) is
\(331\, \mathrm{m}/\mathrm{s}\). An increase of the
temperature by \(1\, ^{\circ } \mathrm{C}\) increases
the speed of velocity by \(0.6\, \mathrm{m}/\mathrm{s}\).
Estimate the sound speed at the temperature
\(18\, ^{\circ } \mathrm{C}\).}
{\(341.8\, \mathrm{m}/\mathrm{s}\)}{\(341.2\, \mathrm{m}/\mathrm{s}\)}{\(348\, \mathrm{m}/\mathrm{s}\)}{\(349\, \mathrm{m}/\mathrm{s}\)}{}{}}
\LANGpl{
\Question{
Prędkość dźwięku w temperaturze
\(0\, ^{\circ } \mathrm{C}\) wynosi 
\(331\, \mathrm{m}/\mathrm{s}\). Wzrost temperatury o  \(1\, ^{\circ } \mathrm{C}\) zwiększa prędkość o \(0.6\, \mathrm{m}/\mathrm{s}\).
Oblicz prędkość dźwięku w temperaturze
\(18\, ^{\circ } \mathrm{C}\).}
{\(341.8\, \mathrm{m}/\mathrm{s}\)}{\(341.2\, \mathrm{m}/\mathrm{s}\)}{\(348\, \mathrm{m}/\mathrm{s}\)}{\(349\, \mathrm{m}/\mathrm{s}\)}{}{}}
\LANGcs{
\Question{
Rychlost zvuku ve vzduchu je při teplotě
\(0\, ^{\circ } \mathrm{C}\) přibližně
\(331\, \mathrm{m}/\mathrm{s}\). Zvýší-li se
teplota o \(1\, ^{\circ } \mathrm{C}\), zvýší
se rychlost zvuku o \(0{,}6\, \mathrm{m}/\mathrm{s}\).
Jaká je rychlost zvuku ve vzduchu při teplotě
\(18\, ^{\circ } \mathrm{C}\)?}
{\(341{,}8\, \mathrm{m}/\mathrm{s}\)}{\(341{,}2\, \mathrm{m}/\mathrm{s}\)}{\(348\, \mathrm{m}/\mathrm{s}\)}{\(349\, \mathrm{m}/\mathrm{s}\)}{}{}}
\LANGsk{
\Question{
Rýchlosť zvuku vo vzduchu je pri teplote
\(0\, ^{\circ } \mathrm{C}\) približne
\(331\, \mathrm{m}/\mathrm{s}\). Ak sa zvýši
teplota o \(1\, ^{\circ } \mathrm{C}\), zvýši
sa rýchlosť zvuku o \(0{,}6\, \mathrm{m}/\mathrm{s}\).
Aká je rýchlosť zvuku vo vzduchu pri teplote
\(18\, ^{\circ } \mathrm{C}\)?}
{\(341{,}8\, \mathrm{m}/\mathrm{s}\)}{\(341{,}2\, \mathrm{m}/\mathrm{s}\)}{\(348\, \mathrm{m}/\mathrm{s}\)}{\(349\, \mathrm{m}/\mathrm{s}\)}{}{}}
\LANGes{
\Question{
La velocidad del sonido a una temperatura de 
\(0\, ^{\circ } \mathrm{C}\) es de 
\(331\, \mathrm{m}/\mathrm{s}\). El aumento de la temperatura en \(1\, ^{\circ } \mathrm{C}\) aumenta la velocidad de sonido en \(0.6\, \mathrm{m}/\mathrm{s}\).
Estima la velocidad de sonido a una temperatura de
\(18\, ^{\circ } \mathrm{C}\).}
{\(341.8\, \mathrm{m}/\mathrm{s}\)}{\(341.2\, \mathrm{m}/\mathrm{s}\)}{\(348\, \mathrm{m}/\mathrm{s}\)}{\(349\, \mathrm{m}/\mathrm{s}\)}{}{}}
\End

\Data\ID{9000009308}
\Updated{2021-04-05 07:04:51}
\Level{C}
\SubArea{Linear functions}
\LANGen{
\Question{
A car moving at a speed \(90\) km/h
starts to brake. It slows down at a constant acceleration
\(2\, \mathrm{m}/\mathrm{s}^{2}\).
Which time is required to stop the car?}
{\(12.5\, \mathrm{s}\)}{\(45\, \mathrm{s}\)}{\(45\, \mathrm{min}\)}{\(12.5\, \mathrm{min}\)}{}{}}
\LANGpl{
\Question{
Samochód poruszający się ze stałą prędkością \(90\) km/h
zaczyna hamować. Zwalnia ze stałym opóźnieniem \(2\, \mathrm{m}/\mathrm{s}^{2}\).
Ile czasu potrzebuje, żeby się zatrzymać?}
{\(12.5\, \mathrm{s}\)}{\(45\, \mathrm{s}\)}{\(45\, \mathrm{min}\)}{\(12.5\, \mathrm{min}\)}{}{}}
\LANGcs{
\Question{
Automobil se pohybuje stálou rychlostí
\(90\) km/h.
Začne rovnoměrně brzdit se stálým zrychlením
\(2\, \mathrm{m}/\mathrm{s}^{2}\). Za
jak dlouho automobil zastaví?}
{\(12{,}5\, \mathrm{s}\)}{\(45\, \mathrm{s}\)}{\(45\, \mathrm{min}\)}{\(12{,}5\, \mathrm{min}\)}{}{}}
\LANGsk{
\Question{
Auto sa pohybuje stálou rýchlosťou
\(90\) km/h.
Ak začne auto brzdiť  rovnomerným spomalením
\(2\, \mathrm{m}/\mathrm{s}^{2}\), za aký dlhý čas zastaví?}
{\(12{,}5\, \mathrm{s}\)}{\(45\, \mathrm{s}\)}{\(45\, \mathrm{min}\)}{\(12{,}5\, \mathrm{min}\)}{}{}}
\LANGes{
\Question{
Un coche que viaja a  \(90\) km/h
empieza a frenar. Disminuye su velocidad con una acceleracin constante de
\(2\, \mathrm{m}/\mathrm{s}^{2}\).
¿Cuánto tiempo necesita para parar?}
{\(12.5\, \mathrm{s}\)}{\(45\, \mathrm{s}\)}{\(45\, \mathrm{min}\)}{\(12.5\, \mathrm{min}\)}{}{}}
\End

\Data\ID{9000009309}
\Updated{2021-04-05 07:04:19}
\Level{C}
\SubArea{Linear functions}
\LANGen{
\Question{
The speed of a swimmer in a \(50\, \mathrm{m}\)-pool
is \(0.8\, \mathrm{m}/\mathrm{s}\).
How fast will he swim two pool lengths (one pool length is
\(50\) meters), if
he/she needs \(2\, \mathrm{s}\) 
to turn at the end of the pool?}
{\(127\, \mathrm{s}\)}{\(82\, \mathrm{s}\)}{\(84\, \mathrm{s}\)}{\(129\, \mathrm{s}\)}{}{}}
\LANGpl{
\Question{
Prędkość pływaka w  basenie o długości \(50\, \mathrm{m}\) 
wynosi \(0.8\, \mathrm{m}/\mathrm{s}\).
Jak szybko dany pływak przepłynie 2 długości basenu (jedna długość wynosi 
\(50\) metrów) jeśli 
potrzebuje na zawrócenie na końcu basenu \(2\, \mathrm{s}\)?}
{\(127\, \mathrm{s}\)}{\(82\, \mathrm{s}\)}{\(84\, \mathrm{s}\)}{\(129\, \mathrm{s}\)}{}{}}
\LANGcs{
\Question{
Rychlost plavce v bazénu o délce \(50\, \mathrm{m}\) 
je \(0{,}8\, \mathrm{m}/\mathrm{s}\).
Za jak dlouho uplave dva bazény (jeden bazén měří
\(50\) 
metrů), trvá-li mu jedna otočka na jeho konci
\(2\, \mathrm{s}\)?}
{\(127\, \mathrm{s}\)}{\(82\, \mathrm{s}\)}{\(84\, \mathrm{s}\)}{\(129\, \mathrm{s}\)}{}{}}
\LANGsk{
\Question{
Rýchlosť plavca v bazéne dĺžky \(50\, \mathrm{m}\) 
je \(0{,}8\, \mathrm{m}/\mathrm{s}\).
Ako dlho mu bude trvať preplávať dva bazény (jeden bazén meria
\(50\) 
metrov), ak mu trvá jedna otočka na konci bazéna
\(2\, \mathrm{s}\)?}
{\(127\, \mathrm{s}\)}{\(82\, \mathrm{s}\)}{\(84\, \mathrm{s}\)}{\(129\, \mathrm{s}\)}{}{}}
\LANGes{
\Question{
La velocidad de un nadador en una piscina de  \(50\, \mathrm{m}\)
es \(0.8\, \mathrm{m}/\mathrm{s}\).
¿Cuánto tardará en nadar dos largos de piscina (un largo de piscina son \(50\) metros),si necesita \(2\, \mathrm{s}\) 
para dar la vuelta al final de la piscina?}
{\(127\, \mathrm{s}\)}{\(82\, \mathrm{s}\)}{\(84\, \mathrm{s}\)}{\(129\, \mathrm{s}\)}{}{}}
\End

\Data\ID{9000009310}
\Updated{2021-04-05 07:04:09}
\Level{C}
\SubArea{Linear functions}
\IMG{000093_10-figure0.pdf}
\LANGen{
\Question{
The graph describes a distance traveled by a motorbike as a function of time. Find
an analytic expression for this function.}
{\(s = 50 + 5t,\ t\in [ 0;30] \)}{\(s = 5 + 50t,\ t\in [ 0;30] \)}{\(s = 50t,\ t\in [ 0;30] \)}{\(s = 5t - 50,\ t\in [ 0;30] \)}{}{}}
\LANGpl{
\Question{
Wykres przedstawia odległość pokonywaną przez motocykl jako funkcję czasu. Wyznacz wzór tej funkcji.}
{\(s = 50 + 5t,\ t\in [ 0;30] \)}{\(s = 5 + 50t,\ t\in [ 0;30] \)}{\(s = 50t,\ t\in [ 0;30] \)}{\(s = 5t - 50,\ t\in [ 0;30] \)}{}{}}
\LANGcs{
\Question{
Na obrázku je graf závislosti dráhy motocyklu na čase. Který předpis
vyjadřuje tuto závislost?}
{\(s = 50 + 5t,\ t\in \langle 0;30\rangle \)}{\(s = 5 + 50t,\ t\in \langle 0;30\rangle \)}{\(s = 50t,\ t\in \langle 0;30\rangle \)}{\(s = 5t - 50,\ t\in \langle 0;30\rangle \)}{}{}}
\LANGsk{
\Question{
Na obrázku je graf závislosti rýchlosti nákladného vlaku od času. Ktorý predpis vyjadruje túto závislosť?}
{\(s = 50 + 5t,\ t\in [ 0;30] \)}{\(s = 5 + 50t,\ t\in [ 0;30] \)}{\(s = 50t,\ t\in [ 0;30] \)}{\(s = 5t - 50,\ t\in [ 0;30] \)}{}{}}
\LANGes{
\Question{
La gráfica representa la distancia viajada por una moto en función de tiempo. Encuentra la expresión analítica para esta función.}
{\(s = 50 + 5t,\ t\in [ 0;30] \)}{\(s = 5 + 50t,\ t\in [ 0;30] \)}{\(s = 50t,\ t\in [ 0;30] \)}{\(s = 5t - 50,\ t\in [ 0;30] \)}{}{}}
\End

\Data\ID{9000009311}
\Updated{2021-04-05 07:04:31}
\Level{C}
\SubArea{Linear functions}
\IMG{000093_11-figure0.pdf}
\LANGen{
\Question{
The graph shows the speed of a train as a function of time. Find an analytic
expression for this function.}
{\(v = 30 - \frac{3}{4}t,\ t\in [ 0;20] \)}{\(v = 30 + \frac{3}{4}t,\ t\in [ 0;20] \)}{\(v = 15 +\frac{3}{4}t,\ t\in [ 0;20] \)}{\(v = 30 - \frac{4}{3}t,\ t\in [ 0;20] \)}{}{}}
\LANGpl{
\Question{
Wykres przedstawia prędkość pociągu jako funkcję czasu. Znajdź wzór tej funkcji.}
{\(v = 30 - \frac{3}{4}t,\ t\in [ 0;20] \)}{\(v = 30 + \frac{3}{4}t,\ t\in [ 0;20] \)}{\(v = 15 + \frac{3}{4}t,\ t\in [ 0;20] \)}{\(v = 30 - \frac{4}{3}t,\ t\in [ 0;20] \)}{}{}}
\LANGcs{
\Question{
Na obrázku je graf závislosti rychlosti nákladního vlaku na čase. Který
předpis vyjadřuje tuto závislost?}
{\(v = 30 - \frac{3}{4}t,\ t\in \langle 0;20\rangle \)}{\(v = 30 + \frac{3}{4}t,\ t\in \langle 0;20\rangle \)}{\(v = 15 + \frac{3}{4}t,\ t\in \langle 0;20\rangle \)}{\(v = 30 - \frac{4}{3}t,\ t\in \langle 0;20\rangle \)}{}{}}
\LANGsk{
\Question{
Na obrázku je graf závislosti rýchlosti nákladného vlaku od času. Ktorý
predpis vyjadruje túto závislosť?}
{\(v = 30 - \frac{3}{4}t,\ t\in \langle 0;20\rangle \)}{\(v = 30 + \frac{3}{4}t,\ t\in \langle 0;20\rangle \)}{\(v = 15 + \frac{3}{4}t,\ t\in \langle 0;20\rangle \)}{\(v = 30 - \frac{4}{3}t,\ t\in \langle 0;20\rangle \)}{}{}}
\LANGes{
\Question{
La gráfica representa la velocidad de un tren en función de tiempo. Encuentra la expresión analítica para esta función.}
{\(v = 30 - \frac{3}{4}t,\ t\in [ 0;20] \)}{\(v = 30 + \frac{3}{4}t,\ t\in [ 0;20] \)}{\(v = 15 +\frac{3}{4}t,\ t\in [ 0;20] \)}{\(v = 30 - \frac{4}{3}t,\ t\in [ 0;20] \)}{}{}}
\End
